% Spectres lumineux et sources de lumière — Physique 1ère D — Cours signé OnBuch+
% Programme officiel MINESEC. Compiler avec : tectonic N12-spectres-lumineux-sources.tex
\newcommand{\DOCMATIERE}{Physique}
\newcommand{\DOCNIVEAU}{1ère D}
\newcommand{\DOCMODULE}{Module 3 : Optique}
\newcommand{\DOCLECON}{N12}
\newcommand{\DOCTITRE}{Spectres lumineux et sources de lumière}
% OnBuch+ — préambule « cours signés » ENRICHI (compilé avec tectonic / XeLaTeX)
% Système visuel « Pop » d'OnBuch+ : fond crème (jamais blanc), bordures encre
% épaisses, ombres « dures » décalées, titres Archivo Black en capitales.
% Version MATHÉMATIQUES : graphes (pgfplots/tikz), boîtes pédagogiques
% (définition, propriété, méthode, attention, le savais-tu, exercice résolu,
% exercice, TP, bilan), page de garde et signature OnBuch+ ancrée en bas.
%
% Macros attendues dans main.tex AVANT \input{preamble} :
%   \DOCMATIERE  \DOCNIVEAU  \DOCMODULE  \DOCLECON (numéro)  \DOCTITRE
\documentclass[11pt]{article}

\usepackage{fontspec}
\usepackage[french]{babel}
\usepackage[a4paper,margin=2cm,top=3.4cm,bottom=2.8cm,headheight=1.4cm,footskip=1.3cm]{geometry}
\usepackage{amsmath,amssymb,amsfonts}
\usepackage{mathtools} % \xmapsto, \coloneqq... souvent utilisés par les modèles
\usepackage{cancel} % simplifications barrées (bilans de forces)
% \abs{z} (module, valeur absolue) et \norm{u} (norme), utilisés par les modèles
\providecommand{\abs}{}\let\abs\relax\DeclarePairedDelimiter{\abs}{\lvert}{\rvert}
\providecommand{\norm}{}\let\norm\relax\DeclarePairedDelimiter{\norm}{\lVert}{\rVert}
\usepackage[table]{xcolor} % [table] : fournit aussi \rowcolors (alternance de lignes)
\usepackage{pagecolor}
\usepackage{tikz}
\usetikzlibrary{arrows.meta,positioning,calc,shapes.geometric,shapes.misc,shapes.arrows,decorations.pathmorphing,decorations.pathreplacing,decorations.markings,patterns,shadows,mindmap,trees,fit,backgrounds,matrix,intersections,angles,quotes}
\usepackage{pgfplots}
\usepackage[version=4]{mhchem} % équations nucléaires \ce{^{235}_{92}U}
\usepackage[european,siunitx]{circuitikz} % schémas électriques
\pgfplotsset{compat=1.18}
\usepgfplotslibrary{fillbetween,groupplots} % aires entre courbes (calcul intégral)
\usepackage{enumitem}
\usepackage{fancyhdr}
% [most] charge toutes les bibliothèques tcolorbox (listings, minted, raster,
% documentation...) et gonfle inutilement la mémoire TeX sur un document long
% (~300 boîtes) : on ne charge que ce qui est réellement utilisé.
\usepackage[skins,breakable]{tcolorbox}
\usepackage{graphicx}
\usepackage{colortbl}
\usepackage{booktabs}
\usepackage{tabularx}
\usepackage{diagbox} % case d'en-tête diagonale des tableaux à double entrée
% Colonnes extensibles centrée / gauche / droite (C, L, R), souvent utilisées par les modèles
\newcolumntype{C}{>{\centering\arraybackslash}X}
\newcolumntype{L}{>{\raggedright\arraybackslash}X}
\newcolumntype{R}{>{\raggedleft\arraybackslash}X}
\usepackage{array}
\usepackage{multicol}
\usepackage{siunitx}
% parse-numbers=false : accepte aussi \SI{2,5 \times 10^{-3}}{...}
\sisetup{locale=FR,output-decimal-marker={,},group-minimum-digits=4,parse-numbers=false}
\DeclareSIUnit\ohm{\ensuremath{\symup{\Omega}}} % U+2126 absent de la police
\DeclareSIUnit\division{div} % réglages d'oscilloscope (ms/div, V/div)
\DeclareSIUnit\tr{tr} % tours (vitesse de rotation, ex. \SI{1500}{\tr\per\minute})
\DeclareSIUnit\ch{ch} % cheval-vapeur (puissance moteur, unité usuelle hors SI)
\DeclareSIUnit\franccfa{FCFA} % franc CFA (coûts, contexte camerounais)
\DeclareSIUnit\dioptre{\ensuremath{\delta}} % vergence d'une lentille (dioptrie)
\usepackage{qrcode}
% \degree : commande fréquemment utilisée par les modèles pour un angle ou une
% température en dehors de \SI{}{} (non fournie par les paquets chargés).
\providecommand{\degree}{^{\circ}}
% \qed (fin de démonstration), utilisé par les modèles sans amsthm
\providecommand{\qed}{\hfill$\square$}
% \barre : double barre des valeurs interdites dans un tableau de variations (tkz-tab)
\providecommand{\barre}{\ensuremath{\|}}
% environnement proof (amsthm non chargé) utilisé par les modèles
\ifdefined\proof\else\newenvironment{proof}[1][Démonstration]{\par\noindent\textit{#1.}\ }{\qed\par}\fi
\usepackage{lastpage}
\usepackage{hyperref}
\hypersetup{hidelinks}

% ============================================================
% Palette « Pop » OnBuch+ — mêmes hex que la classe P (plus_ui.dart)
% ============================================================
\definecolor{poporange}{HTML}{F59321}
\definecolor{popdark}{HTML}{E07A0C}
\definecolor{popcream}{HTML}{FFF6E8}
\definecolor{popink}{HTML}{1C1714}
\definecolor{poppurple}{HTML}{7A5AE0}
\definecolor{popgreen}{HTML}{2E9E6A}
\definecolor{poppink}{HTML}{E0457B}
\definecolor{popblue}{HTML}{2F7DE1}
\definecolor{popgold}{HTML}{C98A12}
\definecolor{popmuted}{HTML}{8A7F76}
\definecolor{popline}{HTML}{EADFD2}
\definecolor{popgray}{HTML}{8A7F76}
\colorlet{popgrey}{popgray}
% Alias tolérants (le modèle nomme parfois les couleurs différemment)
\colorlet{popwhite}{white}
\colorlet{popblack}{popink}
\colorlet{popred}{poppink}
\colorlet{popyellow}{popgold}
% Teintes claires pour fonds de tableaux / zones
\colorlet{popblueL}{popblue!10!white}
\colorlet{poporangeL}{poporange!14!white}
\colorlet{popgreenL}{popgreen!12!white}
\colorlet{poppurpleL}{poppurple!12!white}
\colorlet{poppinkL}{poppink!12!white}
\colorlet{popgoldL}{popgold!14!white}

\pagecolor{popcream}

% ============================================================
% Typographie
% ============================================================
\defaultfontfeatures{Ligatures=TeX}
\setmainfont{PlusJakartaSans-Medium.ttf}[Path=../../../fonts/,BoldFont=PlusJakartaSans-Bold.ttf]
% Maths en sans-serif (Fira Math) avec chiffres et lettres droites de Plus
% Jakarta, pour que les formules se fondent dans le texte.
\usepackage[math-style=ISO,bold-style=ISO]{unicode-math}
\setmathfont{FiraMath-Regular.otf}[Path=../../../fonts/]
\setmathfont{PlusJakartaSans-Medium.ttf}[Path=../../../fonts/,range={up/num,up/{Latin,latin},\mathup}]
\usepackage{icomma} % virgule décimale sans espace en mode math
% Caractères mathématiques Unicode absents des polices de texte (Plus Jakarta,
% Archivo Black) : rendus par leur équivalent mathématique, en texte comme en
% formule — sinon ils s'affichent en carré vide (ex. « Arithmétique dans ℤ »).
\usepackage{newunicodechar}
\newunicodechar{ℝ}{\ensuremath{\mathbb{R}}}
\newunicodechar{ℤ}{\ensuremath{\mathbb{Z}}}
\newunicodechar{ℂ}{\ensuremath{\mathbb{C}}}
\newunicodechar{ℕ}{\ensuremath{\mathbb{N}}}
\newunicodechar{ℚ}{\ensuremath{\mathbb{Q}}}
\newunicodechar{𝔻}{\ensuremath{\mathbb{D}}}
\newunicodechar{α}{\ensuremath{\alpha}}
\newunicodechar{β}{\ensuremath{\beta}}
\newunicodechar{γ}{\ensuremath{\gamma}}
\newunicodechar{δ}{\ensuremath{\delta}}
\newunicodechar{ε}{\ensuremath{\varepsilon}}
\newunicodechar{θ}{\ensuremath{\theta}}
\newunicodechar{λ}{\ensuremath{\lambda}}
\newunicodechar{μ}{\ensuremath{\mu}}
\newunicodechar{σ}{\ensuremath{\sigma}}
\newunicodechar{τ}{\ensuremath{\tau}}
\newunicodechar{φ}{\ensuremath{\varphi}}
\newunicodechar{ψ}{\ensuremath{\psi}}
\newunicodechar{ω}{\ensuremath{\omega}}
\newunicodechar{Σ}{\ensuremath{\Sigma}}
% majuscules produites par \MakeUppercase dans les titres (α-aminés -> Α)
\newunicodechar{Α}{\ensuremath{\alpha}}
\newunicodechar{Β}{\ensuremath{\beta}}
\newunicodechar{Γ}{\ensuremath{\Gamma}}
\newunicodechar{Δ}{\ensuremath{\Delta}}
\newunicodechar{Ε}{\ensuremath{\varepsilon}}
\newunicodechar{Θ}{\ensuremath{\Theta}}
\newunicodechar{Λ}{\ensuremath{\Lambda}}
\newunicodechar{Μ}{\ensuremath{\mu}}
\newunicodechar{Τ}{\ensuremath{\tau}}
\newunicodechar{Φ}{\ensuremath{\Phi}}
\newunicodechar{Ψ}{\ensuremath{\Psi}}
\newunicodechar{Ω}{\ensuremath{\Omega}}
\newunicodechar{↦}{\ensuremath{\mapsto}}
\newunicodechar{∈}{\ensuremath{\in}}
\newunicodechar{∉}{\ensuremath{\notin}}
\newunicodechar{∧}{\ensuremath{\wedge}}
\newunicodechar{⇔}{\ensuremath{\Leftrightarrow}}
\newunicodechar{⟺}{\ensuremath{\Longleftrightarrow}}
\newunicodechar{⇒}{\ensuremath{\Rightarrow}}
\newunicodechar{⟹}{\ensuremath{\Longrightarrow}}
\newunicodechar{∘}{\ensuremath{\circ}}
\newunicodechar{∪}{\ensuremath{\cup}}
\newunicodechar{∩}{\ensuremath{\cap}}
\newunicodechar{≡}{\ensuremath{\equiv}}
\newunicodechar{⊂}{\ensuremath{\subset}}
\newunicodechar{⊥}{\ensuremath{\perp}}
\newunicodechar{∃}{\ensuremath{\exists}}
\newunicodechar{∀}{\ensuremath{\forall}}
\newunicodechar{∛}{\ensuremath{\sqrt[3]{\,}}}
\newunicodechar{∜}{\ensuremath{\sqrt[4]{\,}}}
\newunicodechar{⃗}{\ensuremath{{}^{\to}}}
\newunicodechar{ⁿ}{\ifmmode^{n}\else\textsuperscript{\ensuremath{n}}\fi}
\newunicodechar{ˣ}{\ifmmode^{x}\else\textsuperscript{\ensuremath{x}}\fi}
\newunicodechar{ᵇ}{\ifmmode^{b}\else\textsuperscript{\ensuremath{b}}\fi}
\newunicodechar{ʳ}{\ifmmode^{r}\else\textsuperscript{\ensuremath{r}}\fi}
\newunicodechar{ᵘ}{\ifmmode^{u}\else\textsuperscript{\ensuremath{u}}\fi}
\newunicodechar{ʸ}{\ifmmode^{y}\else\textsuperscript{\ensuremath{y}}\fi}
\newunicodechar{ᵃ}{\ifmmode^{a}\else\textsuperscript{\ensuremath{a}}\fi}
\newunicodechar{ᵉ}{\ifmmode^{e}\else\textsuperscript{\ensuremath{e}}\fi}
\newunicodechar{ᵏ}{\ifmmode^{k}\else\textsuperscript{\ensuremath{k}}\fi}
\newunicodechar{ᵖ}{\ifmmode^{p}\else\textsuperscript{\ensuremath{p}}\fi}
\newunicodechar{ᴺ}{\ifmmode^{N}\else\textsuperscript{\ensuremath{N}}\fi}
\newunicodechar{ᶜ}{\ifmmode^{c}\else\textsuperscript{\ensuremath{c}}\fi}
\newunicodechar{ʰ}{\ifmmode^{h}\else\textsuperscript{\ensuremath{h}}\fi}
\newunicodechar{ᵠ}{\ifmmode^{q}\else\textsuperscript{\ensuremath{q}}\fi}
\newunicodechar{ₙ}{\ifmmode_{n}\else\textsubscript{\ensuremath{n}}\fi}
\newunicodechar{ₐ}{\ifmmode_{a}\else\textsubscript{\ensuremath{a}}\fi}
\newunicodechar{ᵢ}{\ifmmode_{i}\else\textsubscript{\ensuremath{i}}\fi}
\newunicodechar{ᵣ}{\ifmmode_{r}\else\textsubscript{\ensuremath{r}}\fi}
\newunicodechar{ₖ}{\ifmmode_{k}\else\textsubscript{\ensuremath{k}}\fi}
\newunicodechar{ᵦ}{\ifmmode_{\beta}\else\textsubscript{\ensuremath{\beta}}\fi}
\newunicodechar{ₓ}{\ifmmode_{x}\else\textsubscript{\ensuremath{x}}\fi}
\newfontfamily\popdisplay{ArchivoBlack-Regular.ttf}[Path=../../../fonts/]
\newfontfamily\poplogo{SpaceGrotesk-Bold.ttf}[Path=../../../fonts/]
\newfontfamily\popsemi{PlusJakartaSans-SemiBold.ttf}[Path=../../../fonts/]
\newfontfamily\popxbold{PlusJakartaSans-ExtraBold.ttf}[Path=../../../fonts/]
\newfontfamily\popxboldls{PlusJakartaSans-ExtraBold.ttf}[Path=../../../fonts/,LetterSpace=3.0]

\setlist[enumerate]{leftmargin=1.7em,itemsep=5pt,topsep=4pt}
\setlist[itemize]{leftmargin=1.7em,itemsep=4pt,topsep=4pt,label={\color{poporange}\textbullet}}
\setlength{\parindent}{0pt}
\setlength{\parskip}{5pt}
\renewcommand{\arraystretch}{1.25}

% pgfplots : style Pop par défaut
\pgfplotsset{
  every axis/.append style={axis line style={popink,line width=1pt},tick style={popink},
    grid style={popline,line width=0.5pt},label style={font=\small},tick label style={font=\footnotesize}},
  popcurve/.style={poporange,line width=1.8pt,no markers,smooth},
}
% Même style disponible dans un \draw TikZ simple (hors pgfplots)
\tikzset{popcurve/.style={poporange,line width=1.8pt}}
\tikzset{popgrid/.style={popline,very thin}}
\colorlet{popcurve}{poporange} % le nom du style est parfois utilisé comme couleur

% ============================================================
% Pill / squiggle
% ============================================================
\newcommand{\poppill}[3]{%
  \tikz[baseline=-0.55ex]\node[draw=popink,line width=1.2pt,rounded corners=2.6pt,%
    fill=#2,inner xsep=6.5pt,inner ysep=3pt]{%
    \popxboldls\fontsize{7}{7}\selectfont\color{#3}\MakeUppercase{#1}};%
}
\newcommand{\popsquiggle}[1]{%
  \tikz[baseline=-0.4ex]{
    \draw[#1,line width=2.6pt,line cap=round]
      plot [smooth] coordinates {(0,0.09) (0.34,0.01) (0.68,0.21) (1.02,0.01) (1.36,0.10)};
  }%
}
% Mot-clé surligné (marqueur orange sous le texte)
\newcommand{\cle}[1]{{\popxbold\color{popdark}#1}}

% ============================================================
% En-tête / pied
% ============================================================
\pagestyle{fancy}
\fancyhf{}
\renewcommand{\headrulewidth}{0pt}
\renewcommand{\footrulewidth}{0pt}
\lhead{\raisebox{-0.15ex}{\poplogo\fontsize{13}{13}\selectfont\color{popink}OnBuch\color{poporange}+}}
\rhead{\poppill{\DOCMATIERE\ \textbullet\ \DOCNIVEAU\ \textbullet\ Leçon \DOCLECON}{white}{popink}}
\lfoot{\popxbold\fontsize{7.6}{7.6}\selectfont\color{popmuted}\MakeUppercase{Cours signé OnBuch+}\ \textbullet\ {\popsemi\DOCTITRE}}
\rfoot{\tikz[baseline=-0.6ex]\node[rounded corners=8.5pt,fill=popink,minimum height=17pt,minimum width=17pt,inner xsep=5pt,inner sep=0]{\popxbold\fontsize{8}{8}\selectfont\color{popcream}\thepage\,/\,\pageref*{LastPage}};}

% ============================================================
% Titres
% ============================================================
\newcommand{\chaptitre}[2]{%
  \begin{tcolorbox}[enhanced,colback=poporange,colframe=popink,boxrule=2.6pt,arc=11pt,
      drop shadow=popink,left=15pt,right=15pt,top=12pt,bottom=11pt,before skip=3pt,after skip=18pt]
    \poppill{#2}{popink}{popcream}\par\vspace{9pt}
    {\raggedright\hyphenpenalty=10000\exhyphenpenalty=10000\popdisplay\fontsize{22}{24}\selectfont\color{popcream}\MakeUppercase{#1}\par}
  \end{tcolorbox}
}
% Section numérotée : pastille orange avec le numéro + titre Archivo
\newcounter{popsec}
\newcommand{\coursec}[1]{%
  \stepcounter{popsec}\setcounter{popsub}{0}%
  \par\vspace{16pt}\noindent
  \begin{minipage}{\linewidth}
  \tikz[baseline=-0.7ex]\node[draw=popink,line width=1.6pt,fill=poporange,rounded corners=4pt,minimum size=22pt,inner sep=0]{\popdisplay\fontsize{12}{12}\selectfont\color{popcream}\thepopsec};%
  \hspace{8pt}{\popdisplay\fontsize{15}{17}\selectfont\color{popink}\MakeUppercase{#1}}\par
  \vspace{4pt}\noindent{\color{poporange}\rule{36pt}{3.6pt}}
  \end{minipage}\par\nopagebreak\vspace{8pt}
}
\newcounter{popsub}
\newcommand{\courssub}[1]{%
  \stepcounter{popsub}%
  \par\vspace{9pt}\noindent{\popxbold\fontsize{11.5}{13}\selectfont\color{popink}{\color{poporange}\thepopsec.\thepopsub}\hspace{6pt}#1}\par\nopagebreak\vspace{4pt}
}

% ============================================================
% Boîtes pédagogiques — même recette : fond blanc, bordure épaisse colorée,
% ombre dure encre, badge plein. Titre optionnel [#1].
% ============================================================
\tcbset{popbase/.style={breakable,enhanced,colback=white,boxrule=2.3pt,arc=9pt,
  shadow={3pt}{-3pt}{0pt}{popink},left=14pt,right=14pt,top=11pt,bottom=11pt,before skip=11pt,after skip=15pt,
  fontupper=\popsemi\fontsize{10.6}{14.5}\selectfont\color{popink},
  fontlower=\popsemi\fontsize{10.6}{14.5}\selectfont\color{popink}}}
% Coche dessinée (le glyphe ✓ n'existe pas dans les polices du document)
\newcommand{\popcheck}[1]{\tikz[baseline=-0.5ex]\draw[#1,line width=1.6pt,line cap=round,line join=round](-0.13,0.0)--(-0.04,-0.09)--(0.13,0.1);}
\providecommand{\coche}{\popcheck{popgreen}}
\providecommand{\resultat}[1]{\par\smallskip\noindent\fcolorbox{popgreen}{popgreenL}{\parbox{\dimexpr\linewidth-2\fboxsep-2\fboxrule\relax}{\textbf{Résultat.} #1}}\par\smallskip}
\newcommand{\popboxhead}[3]{\poppill{#1}{#2}{white}\ifx\relax#3\relax\else\hspace{6pt}{\popxbold\fontsize{10.6}{12}\selectfont\color{popink}#3}\fi\par\vspace{7pt}}

\newtcolorbox{aretenir}[1][]{popbase,colframe=popgreen,before upper={\popboxhead{À retenir}{popgreen}{#1}}}
\newtcolorbox{exemplebox}[1][]{popbase,colframe=poporange,before upper={\popboxhead{Exemple}{poporange}{#1}}}
\newtcolorbox{definition}[1][]{popbase,colframe=popblue,before upper={\popboxhead{Définition}{popblue}{#1}}}
\newtcolorbox{propriete}[1][]{popbase,colframe=poppurple,before upper={\popboxhead{Propriété}{poppurple}{#1}}}
\newtcolorbox{methode}[1][]{popbase,colframe=popgold,before upper={\popboxhead{Méthode}{popgold}{#1}}}
\newtcolorbox{attention}[1][]{popbase,colframe=poppink,before upper={\popboxhead{Attention}{poppink}{#1}}}
\newtcolorbox{experience}[1][]{popbase,colframe=popblue,colback=popblueL,before upper={\popboxhead{Expérience}{popblue}{#1}}}
% « Le savais-tu ? » : carte violette pleine (contexte, culture, Cameroun)
\newtcolorbox{savaistu}[1][]{popbase,colback=poppurple,colframe=popink,
  fontupper=\popsemi\fontsize{10.4}{14.2}\selectfont\color{white},
  before upper={\poppill{Le savais-tu\,?}{white}{poppurple}\ifx\relax#1\relax\else\hspace{6pt}{\popxbold\color{white}#1}\fi\par\vspace{7pt}}}
% Exercice résolu : énoncé en haut, \tcblower puis la solution sur fond crème
\newcounter{popexo}\newcounter{popexr}
\newtcolorbox{exoresolu}[1][]{popbase,colframe=poporange,colbacklower=popcream,
  segmentation style={popink,line width=1.2pt,dashed},
  before upper={\stepcounter{popexr}\popboxhead{Exercice résolu \thepopexr}{poporange}{#1}},
  before lower={\poppill{Solution}{popink}{popcream}\par\vspace{6pt}}}
% Exercice (non résolu en place) — la correction est dans la partie Corrigés
\newtcolorbox{exercice}[1][]{popbase,colframe=popink,
  before upper={\stepcounter{popexo}\popboxhead{Exercice \thepopexo}{popink}{#1}}}
% Corrigé
\newtcolorbox{corrige}[1][]{popbase,colframe=popgreen,colback=popgreenL,
  before upper={\popboxhead{Corrigé}{popgreen}{#1}}}
% Objectifs / prérequis / bilan
\newtcolorbox{objectifs}{popbase,colframe=popink,colback=poporangeL,
  before upper={\popboxhead{Objectifs de la leçon}{popink}{}}}
\newtcolorbox{prerequis}{popbase,colframe=popink,colback=popblueL,
  before upper={\popboxhead{Prérequis}{popblue}{}}}
\newtcolorbox{bilan}{popbase,colframe=popink,colback=white,boxrule=2.8pt,
  before upper={\popboxhead{Fiche bilan}{poporange}{L'essentiel à connaître pour l'examen}}}
% Figure encadrée avec légende
\newtcolorbox{popfigure}[1][]{enhanced,breakable=false,colback=white,colframe=popline,boxrule=1.4pt,arc=8pt,
  left=10pt,right=10pt,top=10pt,bottom=8pt,before skip=10pt,after skip=12pt,halign=center,
  lower separated=false,halign lower=center,
  fontlower=\popsemi\fontsize{9}{11.5}\selectfont\color{popmuted},
  #1}
\newcommand{\legende}[1]{\par\vspace{6pt}{\popsemi\fontsize{9}{11.5}\selectfont\color{popmuted}#1}}

% ============================================================
% Signature OnBuch+
% ============================================================
\newcommand{\poplogoinline}[1]{{\poplogo\fontsize{#1}{#1}\selectfont\color{popink}OnBuch\color{poporange}+}}
\newcommand{\signatureonbuch}{%
  \begin{tcolorbox}[enhanced,colback=popink,colframe=popink,boxrule=2.2pt,arc=10pt,
      drop shadow=poporange,left=14pt,right=14pt,top=10pt,bottom=10pt,before skip=0pt,after skip=0pt]
    \tikz[baseline=-0.7ex]\node[circle,fill=popgreen,draw=popcream,line width=1.3pt,minimum size=22pt,inner sep=0]{\popcheck{white}};%
    \hspace{9pt}\begin{minipage}[c]{0.62\linewidth}
      {\popxbold\fontsize{12}{13}\selectfont\color{popcream}Signé\ \textbullet\ {\poplogo OnBuch\color{poporange}+}}\par\vspace{2pt}
      {\raggedright\popsemi\fontsize{8}{10.5}\selectfont\color{popline}\MakeUppercase{Équipe pédagogique OnBuch+}\\\MakeUppercase{Conforme au programme officiel MINESEC}\par}
    \end{minipage}\hfill
    {\poplogo\fontsize{22}{22}\selectfont\color{popcream}OnBuch\color{poporange}+}
  \end{tcolorbox}%
}

% ============================================================
% Page de garde
% #1 titre · #2 matière · #3 niveau · #4 module · #5 n° leçon · #6 durée ·
% #7 description / objectifs (texte ou itemize)
% ============================================================
\newcommand{\pagegarde}[7]{%
  \thispagestyle{empty}
  \newgeometry{a4paper,margin=1.7cm,top=1.6cm,bottom=1.4cm}%
  \begin{tikzpicture}[remember picture,overlay]
    \fill[poporange!18] ([xshift=-3.2cm,yshift=-3.6cm]current page.north east) circle (4.6cm);
    \fill[poppurple!14] ([xshift=2.4cm,yshift=7.6cm]current page.south west) circle (3.8cm);
  \end{tikzpicture}%
  {\poplogo\fontsize{30}{30}\selectfont\color{popink}OnBuch\color{poporange}+}\hfill
  \raisebox{0.6ex}{\poppill{Cours signé \textbullet\ Premium}{popink}{popcream}}\par
  \vspace{1pt}\popsquiggle{poppurple}\par
  \vspace{8pt}
  \begin{tcolorbox}[enhanced,colback=poporange,colframe=popink,boxrule=2.8pt,arc=13pt,
      drop shadow=popink,left=18pt,right=18pt,top=12pt,bottom=13pt,after skip=10pt]
    \poppill{#2 \textbullet\ #3}{popink}{popcream}\hspace{5pt}\poppill{#4}{white}{popink}\par\vspace{8pt}
    {\popdisplay\fontsize{14}{14}\selectfont\color{popink}LEÇON #5}\par\vspace{4pt}
    {\raggedright\hyphenpenalty=10000\exhyphenpenalty=10000\popdisplay\fontsize{25}{27}\selectfont\color{popcream}\MakeUppercase{#1}\par}
    \vspace{8pt}
    {\popxbold\fontsize{9.5}{9.5}\selectfont\color{popink}\MakeUppercase{Durée conseillée : #6}}
  \end{tcolorbox}
  \begin{tcolorbox}[enhanced,colback=white,colframe=popink,boxrule=2.2pt,arc=10pt,
      drop shadow=popink,left=16pt,right=16pt,top=10pt,bottom=10pt,after skip=8pt]
    \poppill{Dans cette leçon}{poppurple}{white}\par\vspace{5pt}
    \popsemi\fontsize{9.3}{12.6}\selectfont\color{popink}#7
  \end{tcolorbox}
  \noindent\begin{minipage}[t]{0.487\linewidth}
    \begin{tcolorbox}[enhanced,colback=popgreen,colframe=popink,boxrule=2.2pt,arc=10pt,
        drop shadow=popink,left=11pt,right=11pt,top=10pt,bottom=10pt,height=3.1cm,valign=center]
      \tcbox[colback=white,colframe=popink,boxrule=1.3pt,arc=5pt,boxsep=0pt,nobeforeafter,
        left=3pt,right=3pt,top=3pt,bottom=3pt,valign=center]{\qrcode[height=1.9cm]{https://play.google.com/store/apps/details?id=com.luvvix.onbuch&hl=fr}}%
      \hspace{8pt}\begin{minipage}[c]{0.5\linewidth}
        {\popdisplay\fontsize{11}{12.5}\selectfont\color{white}\MakeUppercase{Télécharge OnBuch}}\par\vspace{4pt}
        {\raggedright\popsemi\fontsize{8.4}{11}\selectfont\color{white}Gratuit \textbullet\ Google Play\\Tuteur IA, annales, concours, résultats\par}
      \end{minipage}
    \end{tcolorbox}
  \end{minipage}\hfill
  \begin{minipage}[t]{0.487\linewidth}
    \begin{tcolorbox}[enhanced,colback=poppurple,colframe=popink,boxrule=2.2pt,arc=10pt,
        drop shadow=popink,left=11pt,right=11pt,top=10pt,bottom=10pt,height=3.1cm,valign=center]
      \tikz[baseline=-0.7ex]\node[circle,fill=white,draw=popink,line width=1.3pt,minimum size=1.6cm,inner sep=0]{\popdisplay\fontsize{24}{24}\selectfont\color{poppurple}+};%
      \hspace{8pt}\begin{minipage}[c]{0.6\linewidth}
        {\popdisplay\fontsize{11}{12.5}\selectfont\color{white}\MakeUppercase{Passe à OnBuch+}}\par\vspace{4pt}
        {\raggedright\popsemi\fontsize{8.4}{11}\selectfont\color{white}Tous les cours signés, Tuteur IA illimité, fiches PDF — onglet OnBuch+ de l'app\par}
      \end{minipage}
    \end{tcolorbox}
  \end{minipage}
  \vspace{8pt}
  \popcontenu
  \vfill
  \signatureonbuch
  \restoregeometry
  \newpage
}

% Rangée « Ce cours contient » (page de garde)
\newcommand{\poptile}[3]{% #1 couleur #2 gros texte #3 légende
  \begin{tcolorbox}[enhanced,colback=#1,colframe=popink,boxrule=2pt,arc=9pt,shadow={2.5pt}{-2.5pt}{0pt}{popink},
      left=6pt,right=6pt,top=9pt,bottom=9pt,halign=center,valign=center,height=1.75cm,width=\linewidth,nobeforeafter]
    {\popdisplay\fontsize{12.5}{13}\selectfont\color{white}\MakeUppercase{#2}}\par\vspace{3pt}
    {\centering\popsemi\fontsize{7.8}{9.5}\selectfont\color{white}#3\par}
  \end{tcolorbox}}
\newcommand{\popcontenu}{%
  \par\noindent{\popxboldls\fontsize{7.5}{7.5}\selectfont\color{popmuted}CE COURS SIGNÉ CONTIENT}\par\vspace{7pt}
  \noindent\begin{minipage}[t]{0.235\linewidth}\poptile{popblue}{Cours}{complet, illustré, pas à pas}\end{minipage}\hfill
  \begin{minipage}[t]{0.235\linewidth}\poptile{poporange}{Méthodes}{et exercices résolus}\end{minipage}\hfill
  \begin{minipage}[t]{0.235\linewidth}\poptile{poppink}{Exercices}{type examen + corrigés}\end{minipage}\hfill
  \begin{minipage}[t]{0.235\linewidth}\poptile{popgreen}{Fiche bilan}{l'essentiel à retenir}\end{minipage}\par}

% Sommaire
\newtcolorbox{sommaire}{enhanced,colback=white,colframe=popink,boxrule=2.2pt,arc=10pt,
  drop shadow=popink,left=18pt,right=18pt,top=13pt,bottom=13pt,before skip=6pt,after skip=20pt,
  before upper={\poppill{Sommaire}{poppurple}{white}\par\vspace{9pt}\popsemi\fontsize{10.6}{14.5}\selectfont\color{popink}}}

% Signature de fin de cours — poussée tout en bas de la dernière page
\newcommand{\findecours}{%
  \par\vfill\nopagebreak
  \begin{center}\popsquiggle{poporange}\end{center}\vspace{4pt}
  \signatureonbuch
}
\AtBeginDocument{\def\llbracket{\mathopen{[\mkern-3.5mu[}}\def\rrbracket{\mathclose{]\mkern-3.5mu]}}} % intervalles d'entiers (glyphe ⟦ absent de la police)
% Glyphes absents de Fira Math : redéfinis à partir de glyphes disponibles
\AtBeginDocument{%
  \def\setminus{\mathbin{\backslash}}\let\smallsetminus\setminus
  \def\vdots{\mathord{\vbox{\baselineskip3.2pt\lineskiplimit0pt\kern4pt\hbox{.}\hbox{.}\hbox{.}}}}%
  \def\ddots{\mathinner{\mkern1mu\raise6.5pt\hbox{.}\mkern2mu\raise3.6pt\hbox{.}\mkern2mu\raise0.7pt\hbox{.}\mkern1mu}}%
  \def\triangle{\mathord{\scalebox{0.9}{$\symup{\Delta}$}}}%
  \def\nabla{\mathord{\raisebox{\depth}{\scalebox{1}[-1]{$\symup{\Delta}$}}}}%
  \def\checkmark{\popcheck{popdark}}%
  \def\star{\mathbin{\ast}}\def\bigstar{\mathord{\ast}}%
  \def\diamond{\mathbin{\scalebox{0.6}{$\Diamond$}}}%
  \def\bigcup{\mathop{\mathchoice{\raisebox{-0.3em}{\scalebox{1.7}{$\cup$}}}{\scalebox{1.25}{$\cup$}}{\cup}{\cup}}\displaylimits}%
  \def\bigcap{\mathop{\mathchoice{\raisebox{-0.3em}{\scalebox{1.7}{$\cap$}}}{\scalebox{1.25}{$\cap$}}{\cap}{\cap}}\displaylimits}%
  \def\lesseqqgtr{\mathrel{\lessgtr}}\def\bigoplus{\mathop{\oplus}\displaylimits}%
}
\providecommand{\ssi}{si et seulement si} % abréviation fréquente
\AtBeginDocument{\providecommand{\wideparen}{\overparen}} % arc de cercle

%% FIGFIX-BEGIN v5 — correctifs globaux des figures (docs/onbuch-generation/tools/figfix)
\usepackage{adjustbox}\usepackage{etoolbox}\tcbuselibrary{raster}
% toute figure TikZ/pgfplots plus large que la ligne est réduite à la largeur disponible
\BeforeBeginEnvironment{tikzpicture}{\begin{adjustbox}{max width=\linewidth}}
\AfterEndEnvironment{tikzpicture}{\end{adjustbox}}
% légendes pgfplots lisibles par-dessus les courbes
\pgfplotsset{every axis/.append style={legend style={fill=white,fill opacity=0.92,text opacity=1,draw=black!15,inner sep=3pt}}}
% étiquettes d'arêtes (syntaxe edge["..."]) avec fond blanc
\tikzset{every edge quotes/.append style={fill=white,inner sep=1.2pt}}
%% FIGFIX-END
\begin{document}

\pagegarde{Spectres lumineux et sources de lumière}{Physique}{1ère D}{Module 3 : Optique}{N12}{7 h}{Cette leçon étudie les spectres lumineux émis par les différentes sources de lumière (étoiles, lampes, laser, DEL) et situe la lumière visible dans le domaine des ondes électromagnétiques. Elle introduit la loi de Wien reliant la couleur d'un corps chauffé à sa température, outil fondamental de l'astrophysique.
\vspace{4pt}
\begin{multicols}{2}\small
\begin{itemize}
\item À la fin de cette leçon, je sais décrire le protocole permettant d'obtenir le spectre d'une lumière à l'aide d'un prisme
\item À la fin de cette leçon, je sais distinguer une source polychromatique d'une source monochromatique caractérisée par sa longueur d'onde dans le vide
\item À la fin de cette leçon, je sais analyser un profil spectral (spectre continu, spectre de raies d'émission)
\item À la fin de cette leçon, je sais citer les domaines des ondes électromagnétiques et situer le visible
\item À la fin de cette leçon, je sais énoncer et exploiter la loi de Wien pour déterminer une température de surface ou une longueur d'onde d'émission maximale
\item À la fin de cette leçon, je sais expliquer l'évolution de la couleur d'un corps chauffé en fonction de sa température
\end{itemize}\end{multicols}}

\begin{sommaire}
\begin{multicols}{2}
\begin{enumerate}
\item Sources de lumière et obtention d'un spectre
\item Sources monochromatiques et polychromatiques
\item Analyse des profils spectraux
\item Les ondes électromagnétiques
\item Couleur des corps chauffés et loi de Wien
\item Démarche expérimentale : lumières colorées
\item Activité d'intégration
\item Méthodes et exercices résolus
\item Exercices
\item Corrigés des exercices
\item Fiche bilan
\end{enumerate}
\end{multicols}
\end{sommaire}

\begin{objectifs}
\begin{itemize}
\item À la fin de cette leçon, je sais décrire le protocole permettant d'obtenir le spectre d'une lumière à l'aide d'un prisme
\item À la fin de cette leçon, je sais distinguer une source polychromatique d'une source monochromatique caractérisée par sa longueur d'onde dans le vide
\item À la fin de cette leçon, je sais analyser un profil spectral (spectre continu, spectre de raies d'émission)
\item À la fin de cette leçon, je sais citer les domaines des ondes électromagnétiques et situer le visible
\item À la fin de cette leçon, je sais énoncer et exploiter la loi de Wien pour déterminer une température de surface ou une longueur d'onde d'émission maximale
\item À la fin de cette leçon, je sais expliquer l'évolution de la couleur d'un corps chauffé en fonction de sa température
\item À la fin de cette leçon, je sais mener une démarche expérimentale simple pour illustrer la notion de lumière colorée
\end{itemize}
\end{objectifs}

\begin{prerequis}
\begin{itemize}
\item Propagation rectiligne de la lumière et réflexion/réfraction (lois de Snell-Descartes, classe de Seconde)
\item Notion d'onde : période, fréquence, célérité, relation λ = cT (début d'année de Première)
\item Dispersion de la lumière blanche par un prisme (notion vue en Seconde)
\item Notion de température et conversion degrés Celsius / kelvins
\item Lecture et exploitation de graphiques (profils, courbes)
\end{itemize}
\end{prerequis}

\newpage

\coursec{Sources de lumière et obtention d'un spectre}

La nuit, à Yaoundé comme à Maroua, les rues s'éclairent de lampadaires orangés, de lampes à incandescence jaunâtres, de DEL blanches et bleutées, tandis que le ciel nous envoie la lumière des étoiles. Pourquoi ces lumières n'ont-elles pas toutes la même couleur ? Pourquoi le fer chauffé dans la forge du quartier devient-il rouge, puis orange, puis blanc ? Et comment les astronomes connaissent-ils la température des étoiles sans jamais les toucher ? Cette leçon répond à ces questions en analysant la lumière émise par les différentes sources.

\courssub{Source de lumière et objet diffusant}

Avant d'étudier la lumière elle-même, distinguons rigoureusement ce qui \emph{produit} la lumière de ce qui se contente de la \emph{renvoyer}.

\begin{definition}[Source primaire et objet diffusant]
Une \textbf{source primaire} (ou source de lumière) émet sa propre lumière par un processus de conversion d'énergie (thermique, électrique, chimique, nucléaire). Un \textbf{objet diffusant} (ou objet éclairé) ne produit pas de lumière ; il intercepte une partie de la lumière incidente et la réémet dans toutes les directions par diffusion.
\end{definition}

\begin{center}
\begin{tabularx}{\linewidth}{|l|X|X|}\hline
\rowcolor{popblueL}
\textbf{Critère} & \textbf{Source primaire} & \textbf{Objet diffusant} \\ \hline
Origine de la lumière & Émission propre & Diffusion (réflexion/diffusion) d'une lumière reçue \\
Exemples naturels & Soleil, étoiles, lucioles, lave volcanique & Lune, planètes, nuages, parois de gorge \\
Exemples artificiels & Lampe à incandescence, tube fluorescent, DEL, laser, flamme & Mur peint, livre, écran de cinéma (réflexion), écran de téléphone (émission mais source secondaire) \\ \hline
\end{tabularx}
\end{center}

\begin{attention}[Piège classique : la Lune n'est pas une source]
Au Probatoire, ne confondez jamais \emph{source de lumière} et \emph{objet lumineux}. La Lune brille dans le ciel nocturne, mais elle ne fait que diffuser la lumière solaire. C'est un objet diffusant, pas une source primaire.
\end{attention}

\courssub{Classification des sources : thermiques et luminescentes}

Les sources primaires se rangent en deux grandes familles selon le mécanisme physique d'émission.

\begin{propriete}[Sources thermiques (incandescence)]
Une \textbf{source thermique} émet de la lumière parce qu'elle est \emph{chaude}. Les atomes/ions vibrent violemment ; les charges accélérées rayonnent un spectre \emph{continu} (toutes les longueurs d'onde présentes) dont la répartition en énergie dépend \emph{uniquement de la température} (loi de Planck, loi de Wien). Plus la température est élevée, plus le pic d'émission se décale vers le bleu.
\begin{itemize}
    \item \textbf{Incandescence solide :} filament de lampe (T \approx \SIrange{2500}{3000}{\kelvin}), fer en fusion, Soleil (surface \approx \SI{5800}{\kelvin}).
    \item \textbf{Incandescence gazeuse / plasma :} flamme de bougie (particules de suie incandescentes), arc électrique de soudage (plasma dense \approx \SI{6000}{\kelvin}).
    \item \textbf{Astrophysique :} étoiles (corps noirs approchés).
\end{itemize}
\end{propriete}

\begin{propriete}[Sources luminescentes (luminescence)]
Une \textbf{source luminescente} émet de la lumière \emph{sans être portée à haute température} (lumière froide). L'émission résulte de transitions quantiques d'électrons entre niveaux d'énergie discrets (atomes, molécules, semi-conducteurs). Le spectre est \emph{discontinu} (lignes ou bandes d'émission).
\begin{itemize}
    \item \textbf{Electroluminescence :} DEL (diode électroluminescente), tube fluorescent, lampe néon, panneau d'affichage.
    \item \textbf{Photoluminescence :} tube fluorescent (UV $\to$ visible par fluorescence), objets « phosphorescents ».
    \item \textbf{Chimiluminescence :} lucioles, bâtons lumineux (chimiluminescence).
    \item \textbf{Emission stimulée :} LASER (Light Amplification by Stimulated Emission of Radiation) : spectre extrêmement étroit, cohérent.
\end{itemize}
\end{propriete}

\begin{savaistu}[Le lampadaire de votre quartier]
À Douala ou à Bertoua, observez les lampadaires :
\begin{itemize}
    \item \textbf{Orange/jaune intense, spectre à lignes larges :} lampe \textbf{sodium haute pression} (très efficace, utilisée par ENEO pour l'éclairage public).
    \item \textbf{Blanc bleuté, spectre continu + lignes :} lampe \textbf{iodures métalliques} ou \textbf{vapeur de mercure}.
    \item \textbf{Blanc pur, spectre large (bande) :} projecteur \textbf{DEL} (remplace progressivement les anciennes technologies).
    \item \textbf{Jaune pâle, spectre continu :} vieille lampe à \textbf{incandescence} (quasi disparue, rendement \approx \SI{5}{\%} seulement).
\end{itemize}
\end{savaistu}

\begin{center}
\begin{tabularx}{\linewidth}{|l|X|X|X|}\hline
\rowcolor{popblueL}
\textbf{Type} & \textbf{Mécanisme} & \textbf{Spectre} & \textbf{Exemples} \\ \hline
Thermique & Agitation thermique (corps noir) & Continu & Soleil, étoile, lampe à filament, flamme (particules), fer rouge \\ \hline
Luminescente & Transitions quantiques discrètes & Discontinu (lignes/bande) & DEL, tube fluo, néon, laser, lucioles, lampe sodium \\ \hline
\end{tabularx}
\end{center}

\courssub{Dispositif expérimental : obtenir un spectre au prisme}

Pour « voir » la composition d'une lumière, il faut la \emph{disperser} : séparer ses différentes radiations selon leur longueur d'onde.

\begin{methode}[Protocole type : obtention du spectre d'une source au prisme]
\textbf{Matériel :} source à étudier (lampe, Soleil, flamme...), diaphragme à fente fine (ou carte noire avec fente), prisme en verre (ou réseau de diffraction), écran blanc (papier, mur).
\begin{enumerate}
    \item Placer la source devant la fente fine. La fente sélectionne un faisceau fin et bien défini (source secondaire ponctuelle).
    \item Orienter le faisceau émergent vers une face du prisme (pas sur l'arête !).
    \item Observer sur l'écran placé derrière le prisme la tache lumineuse déviée et étalée : c'est le \textbf{spectre}.
    \item Noter l'ordre des couleurs, leur étendue, la présence de zones sombres (raies d'absorption) ou de lignes brillantes (raies d'émission).
\end{enumerate}
\textbf{Schéma légendé obligatoire} (voir figure ci-dessous).
\end{methode}

\begin{popfigure}
\begin{tikzpicture}[scale=1.1, every node/.style={font=\small}]
    % Source
    \node[draw, fill=poporange!30, rounded corners, minimum width=1.5cm, minimum height=1cm, align=center] (source) at (0,0){Source\\$S$};
    \draw[poporange, thick, ->] (source.east) -- ++(0.5,0) node[midway, above] {Lumière};
    
    % Fente
    \draw[line width=3pt, popdark] (1.5,-0.15) -- (1.5,0.15);
    \draw[line width=3pt, popdark] (1.55,-0.15) -- (1.55,0.15);
    \node[above] at (1.525,0.25) {Fente fine};
    
    % Prisme (triangle équilatéral)
    \coordinate (A) at (3,0.8);
    \coordinate (B) at (3,-0.8);
    \coordinate (C) at (4.5,0);
    \draw[fill=popblue!10, draw=popblue, thick] (A) -- (B) -- (C) -- cycle;
    \node[below] at (3.75,-1.1) {Prisme en verre ($n_{violet}>n_{rouge}$)};
    
    % Faisceau incident
    \draw[poporange, thick, ->] (1.525,0) -- (3,0) node[midway, above] {$i$};
    % Normale entrée
    \draw[dashed, thin, popmuted] (3,0.8) -- (3,-0.8);
    \node[left] at (3,0.9) {$N_1$};
    
    % Faisceaux émergents dispersés (couleurs RVB simplifié)
    \draw[red, thick, ->] (3.5,0.3) -- (5.5,0.6) node[right] {Rouge ($\lambda \approx \SI{650}{\nano\meter}$)};
    \draw[orange, thick, ->] (3.5,0.1) -- (5.5,0.25) node[right] {Orange};
    \draw[yellow!80!black, thick, ->] (3.5,-0.05) -- (5.5,-0.05) node[right] {Jaune};
    \draw[green, thick, ->] (3.5,-0.2) -- (5.5,-0.3) node[right] {Vert};
    \draw[blue, thick, ->] (3.5,-0.35) -- (5.5,-0.5) node[right] {Bleu};
    \draw[purple, thick, ->] (3.5,-0.5) -- (5.5,-0.7) node[right] {Violet ($\lambda \approx \SI{400}{\nano\meter}$)};
    
    % Normale sortie (approximée au milieu de la face AC)
    \coordinate (M) at ($(A)!0.5!(C)$);
    \draw[dashed, thin, popmuted] ($(M)+(-0.3,0.5)$) -- ($(M)+(0.3,-0.5)$);
    \node[above right] at ($(M)+(0.3,0.5)$) {$N_2$};
    
    % Angle de déviation pour le rouge (exemple)
    \draw[poppurple, thick, <->] (4.5,0) arc (0:15:1.5) node[midway, right=2pt] {$\delta_R$};
    \draw[poppurple, thick, <->] (4.5,0) arc (0:25:1.2) node[midway, right=2pt] {$\delta_V$};
    
    % Écran
    \draw[line width=4pt, popdark] (6.5,-1) -- (6.5,1);
    \node[left] at (6.5,1.2) {Écran blanc ($\perp$ faisceau moyen)};
    
    % Spectre sur l'écran (rectangles colorés)
    \foreach \y/\col in {0.8/red, 0.5/orange, 0.2/yellow!80!black, -0.1/green, -0.4/blue, -0.7/purple} {
        \draw[fill=\col, draw=none] (6.5,\y-0.1) rectangle (6.7,\y+0.1);
    }
    \node[right] at (6.8,0) {Spectre observé};
    
    % Flèche direction non déviée
    \draw[poppurple, dashed, ->] (3,0) -- (4.5,0) node[midway, above] {Direction non déviée};
\end{tikzpicture}
\legende{Dispositif classique d'obtention d'un spectre par dispersion au prisme. La fente assure la résolution spectrale ; le prisme dévie chaque longueur d'onde différemment ($n_{violet} > n_{rouge}$). $\delta_V > \delta_R$.}
\end{popfigure}

\begin{attention}[Conditions de réussite : fente et incidence]
\begin{itemize}
    \item \textbf{Fente fine :} Si la fente est large, les images des différentes couleurs se chevauchent $\to$ spectre flou, couleurs mélangées (perte de résolution spectrale).
    \item \textbf{Incidence sur une face :} Le faisceau doit frapper une \emph{face plane} du prisme. S'il arrive sur l'arête (sommet), pas de dispersion nette, pertes par réflexions parasites.
    \item \textbf{Écran perpendiculaire :} Pour éviter la distorsion (astigmatisme), l'écran doit être perpendiculaire à la direction moyenne du faisceau émergent.
\end{itemize}
\end{attention}

\courssub{Le prisme disperse : pourquoi le violet est-il plus dévié que le rouge ?}

Le prisme en verre est un milieu dispersif : son indice de réfraction $n$ dépend de la longueur d'onde $\lambda$ (dans le vide). Pour le verre couronnement typique, la formule de Cauchy donne une approximation :
\[
n(\lambda) \approx A + \frac{B}{\lambda^2} \quad (A,B>0,\ \lambda \text{ en } \si{\micro\meter})
\]
Donc $n$ \emph{décroît} quand $\lambda$ \emph{croît} : $n_{violet} > n_{bleu} > n_{vert} > n_{jaune} > n_{orange} > n_{rouge}$.
\emph{Analyse dimensionnelle :} $B$ a dimension de $\lambda^2$ (ex. $\si{\micro\meter}^2$) pour que $B/\lambda^2$ soit sans dimension.

La déviation $\delta$ subie par un rayon traversant un prisme d'angle $A$ (au minimum de déviation) vaut :
\[
n = \frac{\sin\frac{A+\delta_m}{2}}{\sin\frac{A}{2}} \quad \Rightarrow \quad \delta_m = 2\arcsin\left(n\sin\frac{A}{2}\right) - A
\]
Puisque $n$ est plus grand pour le violet, $\delta_m$ est plus grand pour le violet. \textbf{Le violet est le plus dévié, le rouge le moins dévié.}

\begin{definition}[Spectre lumineux d'une source]
Le \textbf{spectre lumineux} d'une source est la figure obtenue sur l'écran après dispersion de la lumière émise par cette source. Il représente la répartition de l'énergie lumineuse (ou de l'intensité spectrale) en fonction de la longueur d'onde $\lambda$ (ou de la fréquence $\nu$, ou du nombre d'onde $\tilde{\nu}=1/\lambda$).
\end{definition}

\courssub{Alternative : le réseau de diffraction}

Au Probatoire, on peut citer le \textbf{réseau de diffraction} comme alternative au prisme pour obtenir un spectre.

\begin{propriete}[Réseau de diffraction (rappel)]
Un réseau est une lame de verre (ou métal) portant un grand nombre de fentes parallèles, équidistantes, de pas $a$ (typiquement \SIrange{100}{1000}{\per\milli\meter}). La condition des maxima d'interférence (ordres $k \in \mathbb{Z}$) pour une incidence $i$ et une diffraction $r_k$ (angles algébriques, même côté de la normale) est :
\[
a(\sin i + \sin r_k) = k\lambda
\]
\begin{itemize}
    \item L'ordre $0$ ($k=0$) donne une tache blanche non dispersée (toutes $\lambda$ superposées).
    \item Les ordres $k \neq 0$ donnent des spectres : $\lambda$ plus grande $\to$ $|r_k|$ plus grand. \textbf{Au réseau, le rouge est plus dévié que le violet (sens inverse du prisme).}
    \item Résolution spectrale bien supérieure au prisme (raies très fines).
\end{itemize}
\end{propriete}

\begin{attention}[Prisme vs Réseau : sens de la dispersion]
Ne confondez pas l'ordre des couleurs ! C'est une question classique de QCM ou de rédaction au Probatoire.
\begin{itemize}
    \item \textbf{Prisme (réfraction) :} $n \downarrow$ quand $\lambda \uparrow$ $\Rightarrow$ \textbf{Violet le plus dévié, Rouge le moins dévié}.
    \item \textbf{Réseau (diffraction) :} $\sin r_k \propto \lambda$ (pour $k>0$) $\Rightarrow$ \textbf{Rouge le plus dévié, Violet le moins dévié}.
\end{itemize}
\end{attention}

\courssub{Exemple résolu : analyse d'un dispositif}

\begin{exoresolu}[Spectre d'une lampe sodium]
Énoncé : On souhaite obtenir le spectre d'une lampe à vapeur de sodium (jaune orangé) à l'aide d'un prisme en verre d'angle $A=\SI{60}{\degree}$ et d'indice $n_{jaune}=\num{1,52}$ (pour la raie D du sodium $\lambda=\SI{589,3}{\nano\meter}$). La fente est éclairée par la lampe. On observe l'écran placé perpendiculairement au faisceau émergent moyen.
\begin{enumerate}
    \item Dessiner le schéma annoté du dispositif.
    \item Calculer l'angle de déviation minimum $\delta_m$ pour la raie jaune.
    \item Préciser l'aspect du spectre observé sur l'écran.
\end{enumerate}

\tcblower
\textbf{Solution détaillée}
\begin{enumerate}
    \item \textbf{Schéma :} Source $\to$ Fente fine $\to$ Prisme (face d'entrée) $\to$ Faisceau dévié $\to$ Écran. (Voir figure générale ci-dessus, avec une seule couleur jaune).
    \item \textbf{Calcul de $\delta_m$ :}
    Relation du minimum de déviation :
    \[
    n = \frac{\sin\frac{A+\delta_m}{2}}{\sin\frac{A}{2}}
    \]
    Avec $A=\SI{60}{\degree}$, $\sin\frac{A}{2}=\sin\SI{30}{\degree}=\num{0,5}$.
    \[
    \num{1,52} = \frac{\sin\frac{\SI{60}{\degree}+\delta_m}{2}}{\num{0,5}} \quad \Rightarrow \quad \sin\frac{\SI{60}{\degree}+\delta_m}{2} = \num{0,76}
    \]
    \[
    \frac{\SI{60}{\degree}+\delta_m}{2} = \arcsin(\num{0,76}) \approx \SI{49,46}{\degree}
    \]
    \[
    \SI{60}{\degree}+\delta_m \approx \SI{98,92}{\degree} \quad \Rightarrow \quad \delta_m \approx \SI{38,9}{\degree}
    \]
    \textbf{Résultat :} $\delta_m \approx \SI{38,9}{\degree}$ (ou \SI{39}{\degree} avec 2 chiffres significatifs).
    \item \textbf{Aspect du spectre :} La lampe sodium émet essentiellement deux raies très proches (doublet D : $\lambda_1=\SI{589,0}{\nano\meter}$, $\lambda_2=\SI{589,6}{\nano\meter}$). Avec un prisme de résolution modeste, on observe \textbf{une seule raie jaune fine} (ou un doublet à peine résolu si le prisme est de qualité). Le fond est noir. C'est un \textbf{spectre de raies d'émission} caractéristique d'une source luminescente (décharge électrique dans une vapeur de sodium).
\end{enumerate}
\end{exoresolu}

\courssub{Exercices d'application}

\begin{exercice}[Classification de sources]
Parmi les sources suivantes, identifier les sources thermiques et les sources luminescentes. Justifier en une phrase chacune.
\begin{enumerate}
    \item Le Soleil (surface $\approx \SI{5800}{\kelvin}$).
    \item Une lampe torche à DEL blanche.
    \item Une flamme de chalumeau oxyacétylénique ($T \approx \SI{3200}{\kelvin}$).
    \item Un tube néon rouge (enseigne publicitaire).
    \item Un projecteur de cinéma à lampe xénon (arc électrique).
    \item Un ver luisant (lucioles).
\end{enumerate}
\end{exercice}

\begin{exercice}[Dispositif au prisme : erreurs à corriger]
Un élève veut obtenir le spectre du Soleil. Il perce un trou de \SI{5}{\milli\meter} de diamètre dans un carton (au lieu d'une fente), place le prisme de manière que le soleil frappe l'arête du prisme, et met l'écran parallèlement à la face de sortie du prisme.
\begin{enumerate}
    \item Citer trois erreurs dans ce protocole.
    \item Expliquer la conséquence de chaque erreur sur le spectre observé.
\end{enumerate}
\end{exercice}

\begin{corrige}[Exercice : Classification de sources]
\begin{enumerate}
    \item \textbf{Thermique.} Le Soleil rayonne comme un corps noir à $\approx \SI{5800}{\kelvin}$ (spectre continu + raies d'absorption de l'atmosphère solaire).
    \item \textbf{Luminescente.} La DEL émet par recombinaison électron-trou dans une jonction semi-conductrice (spectre large mais discontinu, pas de corps noir).
    \item \textbf{Thermique (majoritaire).} La flamme est un gaz chaud contenant des particules de suie incandescentes (spectre continu) + raies d'émission moléculaires (CH, C$_2$) ; l'émission dominante en intensité est thermique (incandescence des particules de carbone).
    \item \textbf{Luminescente.} Le néon émet par décharge dans un gaz rare (transitions atomiques du néon) : spectre de raies rouges/oranges.
    \item \textbf{Thermique (principalement).} L'arc xénon produit un plasma très chaud ($\approx \SI{6000}{\kelvin}$) : spectre quasi continu (corps noir) + raies d'émission xénon. Source thermique intense.
    \item \textbf{Luminescente.} Chimiluminescence : réaction chimique (luciférine + luciférase + O$_2$) excite une molécule qui émet un photon visible (vert/jaune). Lumière froide.
\end{enumerate}
\end{corrige}

\begin{corrige}[Exercice : Dispositif au prisme]
\begin{enumerate}
    \item \textbf{Erreur 1 :} Trou circulaire de \SI{5}{\milli\meter} au lieu d'une fente fine.
    \item \textbf{Erreur 2 :} Faisceau incident sur l'arête (sommet) du prisme au lieu d'une face.
    \item \textbf{Erreur 3 :} Écran parallèle à la face de sortie au lieu d'être perpendiculaire au faisceau émergent moyen.
\end{enumerate}
\begin{enumerate}
    \item \textbf{Conséquence 1 :} La source secondaire est étendue $\to$ chaque longueur d'onde donne une tache large $\to$ les couleurs se chevauchent $\to$ \textbf{spectre flou, bandes larges, mauvaise résolution spectrale}.
    \item \textbf{Conséquence 2 :} Pas de traversée du prisme selon la géométrie prévue (deux réfractions) $\to$ réflexions parasites, pas de dispersion nette, lumière diffusée $\to$ \textbf{pas de spectre exploitable}.
    \item \textbf{Conséquence 3 :} Astigmatisme / distorsion : les rayons extrêmes (violet/rouge) arrivent obliquement sur l'écran $\to$ \textbf{taches elliptiques, étalement inégal, mesure d'angles faussée}.
\end{enumerate}
\end{corrige}

\coursec{Sources monochromatiques et polychromatiques}

Dans la section précédente, nous avons appris à obtenir le spectre d'une source lumineuse à l'aide d'un prisme : la lumière blanche du Soleil ou d'une lampe à incandescence se décompose en une bande colorée continue, du violet au rouge, tandis que le faisceau d'un laser donne une tache d'une seule couleur. Cette observation expérimentale pose une question fondamentale : toutes les lumières se valent-elles ? La réponse est non, et c'est précisément l'objet de cette section : certaines sources n'émettent qu'\emph{une seule} radiation, d'autres en émettent une multitude. Nous allons apprendre à les distinguer, à caractériser chaque radiation par une grandeur précise — sa longueur d'onde dans le vide — et à exploiter la relation qui relie longueur d'onde et fréquence.

\courssub{La lumière monochromatique : une seule radiation}

\textbf{Observation.} Dirigeons le faisceau d'un laser rouge (comme celui d'un pointeur utilisé en classe) vers un prisme. Contrairement à la lumière blanche, le faisceau n'est \emph{pas} décomposé en bande colorée : il est simplement dévié, et on observe sur l'écran une unique tache rouge. Le prisme, pourtant capable de disperser la lumière blanche, ne trouve ici « rien à séparer ». La lumière du laser est donc constituée d'une seule « couleur pure », c'est-à-dire d'une seule radiation.

\begin{definition}[Source monochromatique]
Une source de lumière est dite \cle{monochromatique} (du grec \emph{monos}, seul, et \emph{khrôma}, couleur) si elle émet une \textbf{seule radiation}. Cette radiation est caractérisée par sa \cle{longueur d'onde dans le vide}, notée $\lambda_0$, exprimée en mètres (m) dans le système international, et très souvent en nanomètres (nm) en optique :
\[ 1\ \text{nm} = 10^{-9}\ \text{m}. \]
À chaque valeur de $\lambda_0$ du domaine visible correspond une couleur bien déterminée : on parle de radiation de couleur pure.
\end{definition}

\begin{exemplebox}[Quelques sources monochromatiques ou quasi monochromatiques]
\begin{itemize}
\item Le \textbf{laser hélium-néon} (He-Ne), utilisé dans les pointeurs et certains instruments de laboratoire, émet une radiation rouge de longueur d'onde $\lambda_0 \approx 633\ \text{nm}$.
\item La \textbf{lampe à vapeur de sodium}, qui éclaire en orange certaines rues de Yaoundé et de Douala, émet essentiellement une radiation jaune-orangée de $\lambda_0 = 589\ \text{nm}$ (en réalité un doublet très resserré, assimilé à une seule radiation à notre niveau).
\item Les \textbf{DEL colorées} (rouge, verte, bleue) des panneaux publicitaires et des feux de signalisation sont \emph{quasi} monochromatiques : leur lumière est concentrée dans un intervalle très étroit de longueurs d'onde autour d'une valeur centrale.
\end{itemize}
\end{exemplebox}

\begin{attention}[Le cas piégeux de la DEL blanche]
Une DEL blanche, malgré son aspect uniforme et « pur », n'est \textbf{pas} monochromatique ! Elle est fabriquée à partir d'une DEL bleue recouverte d'un luminophore qui convertit une partie de la lumière bleue en lumière jaune : sa lumière est donc un mélange de plusieurs radiations. Son spectre, loin d'être une raie unique, présente un pic dans le bleu et une large bande dans le jaune-vert-rouge. Retiens bien : \emph{l'aspect visuel uniforme d'une lumière ne prouve rien ; seule l'analyse spectrale (au prisme ou au réseau) permet de conclure}.
\end{attention}

\courssub{La lumière polychromatique : une superposition de radiations}

\textbf{Observation.} Reprenons l'expérience du prisme avec la lumière du Soleil ou celle d'une lampe à incandescence : nous obtenons un spectre \emph{continu}, bande colorée où toutes les couleurs de l'arc-en-ciel se succèdent sans interruption. Cela signifie que cette lumière contient une infinité de radiations de longueurs d'onde différentes, superposées.

\begin{definition}[Source polychromatique]
Une source de lumière est dite \cle{polychromatique} si elle émet une \textbf{superposition de plusieurs radiations} de longueurs d'onde différentes. Son spectre peut être :
\begin{itemize}
\item \textbf{continu} : toutes les radiations d'un intervalle sont présentes (Soleil, lampe à incandescence, flamme) ;
\item \textbf{discontinu (ou de raies)} : seules certaines radiations bien précises sont présentes (lampes à vapeur de mercure, tubes fluorescents, enseignes au néon).
\end{itemize}
La \textbf{lumière blanche} est le cas polychromatique le plus célèbre : elle contient toutes les radiations visibles, du violet au rouge.
\end{definition}

\begin{popfigure}
\begin{center}
\begin{tikzpicture}[x=0.016cm,y=1cm]
% Bande spectrale du visible de 400 à 800 nm
\shade[left color=violet!80!black, right color=red] (400,0) rectangle (800,1.2);
\draw[thick] (400,0) rectangle (800,1.2);
\foreach \x in {400,450,500,550,600,650,700,750,800}{
  \draw[thick] (\x,0) -- (\x,-0.15);
  \node[below] at (\x,-0.2) {\small \pgfmathprintnumber{\x}};
}
\node[below] at (600,-0.85) {\small $\lambda_0$ (nm)};
\node[above] at (475,1.25) {\small violet};
\node[above] at (500,1.25) {\small bleu};
\node[above] at (545,1.25) {\small vert};
\node[above] at (585,1.25) {\small jaune};
\node[above] at (630,1.25) {\small orange};
\node[above] at (730,1.25) {\small rouge};
\end{tikzpicture}
\end{center}
\legende{Le spectre continu de la lumière blanche, gradué en nanomètres : chaque longueur d'onde correspond à une couleur. Le spectre visible s'étend environ de 400 nm (violet) à 800 nm (rouge).}
\end{popfigure}

\begin{popfigure}
\begin{center}
\begin{tikzpicture}[x=0.016cm,y=0.9cm]
% Spectre de la lumière blanche (haut)
\shade[left color=violet!80!black, right color=red] (400,2.6) rectangle (800,3.6);
\draw[thick] (400,2.6) rectangle (800,3.6);
\node[left] at (395,3.1) {\small Lumière blanche};
% Spectre du laser (bas) : une seule raie
\fill[black] (400,0.6) rectangle (800,1.6);
\fill[red] (631,0.6) rectangle (635,1.6);
\draw[thick] (400,0.6) rectangle (800,1.6);
\node[left] at (395,1.1) {\small Laser He-Ne};
\node[above, red] at (633,1.7) {\small $\lambda_0 = 633$ nm};
\foreach \x in {400,500,600,700,800}{
  \draw[thick] (\x,0.6) -- (\x,0.45);
  \node[below] at (\x,0.4) {\small \pgfmathprintnumber{\x}};
}
\node[below] at (600,-0.15) {\small $\lambda_0$ (nm)};
\end{tikzpicture}
\end{center}
\legende{Comparaison des spectres : la lumière blanche donne un spectre continu (toutes les radiations), le laser He-Ne donne une seule raie rouge à 633 nm (source monochromatique).}
\end{popfigure}

\courssub{Longueur d'onde et fréquence : la relation fondamentale}

\textbf{Intuition.} Une radiation lumineuse est une onde électromagnétique : elle se propage dans le vide à la célérité $c$, en vibrant à une fréquence $\nu$ (lettre grecque « nu »). Comme pour une vague qui avance sur l'eau, la distance parcourue pendant une période de vibration est la longueur d'onde. Or en une période $T = 1/\nu$, l'onde parcourt la distance $c \times T$. D'où la relation fondamentale :

\begin{propriete}[Relation entre longueur d'onde et fréquence]
Toute radiation monochromatique est caractérisée par sa \cle{fréquence} $\nu$ (en hertz, Hz) et sa \cle{longueur d'onde dans le vide} $\lambda_0$ (en m), liées par :
\[ \boxed{\lambda_0 = \dfrac{c}{\nu}} \qquad \text{ou de façon équivalente} \qquad \nu = \dfrac{c}{\lambda_0}, \]
où $c = 3{,}00 \times 10^{8}\ \text{m}\cdot\text{s}^{-1}$ est la célérité de la lumière dans le vide.
\end{propriete}

\textbf{Vérification par analyse dimensionnelle.} $[\lambda_0] = \dfrac{[c]}{[\nu]} = \dfrac{\text{m}\cdot\text{s}^{-1}}{\text{s}^{-1}} = \text{m}$. La relation est homogène : une longueur d'onde s'exprime bien en mètres.

\begin{propriete}[Invariance de la fréquence — admise]
La \textbf{fréquence} $\nu$ d'une radiation est une caractéristique \emph{intrinsèque} de cette radiation : elle \textbf{ne dépend pas du milieu} de propagation. C'est elle qui définit la couleur de la radiation. En revanche, la célérité de la lumière et la longueur d'onde changent lorsque la radiation passe du vide (ou de l'air) dans un milieu transparent d'indice $n$ : la longueur d'onde dans le milieu devient
\[ \lambda = \dfrac{\lambda_0}{n} \qquad (n \geq 1). \]
Ainsi, une radiation rouge ($\lambda_0 = 633$ nm) reste rouge dans l'eau ou le verre, mais sa longueur d'onde y est plus petite.
\end{propriete}

\begin{exemplebox}[Fréquence de la radiation jaune du sodium]
La lampe à vapeur de sodium émet une radiation de longueur d'onde dans le vide $\lambda_0 = 589\ \text{nm}$. Calculons sa fréquence :
\begin{align*}
\nu &= \dfrac{c}{\lambda_0} = \dfrac{3{,}00 \times 10^{8}\ \text{m}\cdot\text{s}^{-1}}{589 \times 10^{-9}\ \text{m}} \\
&= \dfrac{3{,}00 \times 10^{8}}{5{,}89 \times 10^{-7}}\ \text{Hz} \\
&\approx 5{,}09 \times 10^{14}\ \text{Hz}.
\end{align*}
Cette radiation vibre donc environ 509 000 milliards de fois par seconde ! Les fréquences lumineuses sont considérables, ce qui explique qu'on préfère souvent raisonner en longueurs d'onde.
\end{exemplebox}

\begin{attention}[Le piège classique : nm contre m]
Dans la relation $\nu = c/\lambda_0$, la longueur d'onde doit être exprimée en \textbf{mètres}, puisque $c$ est en m/s. Oublier la conversion $1\ \text{nm} = 10^{-9}\ \text{m}$ est l'erreur la plus fréquente au Probatoire : écrire $\nu = \dfrac{3 \times 10^{8}}{589}$ donne un résultat faux de neuf ordres de grandeur ! Réflexe à acquérir : \emph{dès que tu lis des nanomètres, convertis immédiatement en mètres avant tout calcul}.
\end{attention}

\courssub{Le spectre visible et la correspondance couleur--longueur d'onde}

L'œil humain n'est sensible qu'à un étroit intervalle de longueurs d'onde : c'est le domaine du \cle{visible}.

\begin{aretenir}[Limites du spectre visible]
Le spectre visible s'étend approximativement de $\lambda_0 \approx 400\ \text{nm}$ (violet) à $\lambda_0 \approx 800\ \text{nm}$ (rouge). En deçà de 400 nm se trouve l'\textbf{ultraviolet} (UV), au-delà de 800 nm l'\textbf{infrarouge} (IR) : ces radiations sont invisibles pour l'œil, mais bien réelles (elles seront étudiées dans la section sur les ondes électromagnétiques).
\end{aretenir}

\begin{center}
\begin{tabularx}{\linewidth}{|l|X|X|}\hline
\rowcolor{popblueL} \textbf{Couleur} & \textbf{Domaine de $\lambda_0$ (nm)} & \textbf{Exemple de source} \\ \hline
Violet & 400 -- 430 & Lampes UV proches (détection de faux billets) \\ \hline
Bleu & 430 -- 490 & DEL bleue, $\lambda_0 \approx 470$ nm \\ \hline
Vert & 490 -- 560 & Laser vert, $\lambda_0 \approx 532$ nm \\ \hline
Jaune & 560 -- 590 & Lampe à vapeur de sodium, 589 nm \\ \hline
Orange & 590 -- 620 & Enseignes lumineuses \\ \hline
Rouge & 620 -- 800 & Laser He-Ne, 633 nm ; DEL rouge, 650 nm \\ \hline
\end{tabularx}
\end{center}

Remarque l'ordre : le violet correspond aux \emph{plus petites} longueurs d'onde (donc aux plus \emph{grandes} fréquences, puisque $\nu = c/\lambda_0$), le rouge aux plus grandes longueurs d'onde (plus petites fréquences). C'est d'ailleurs pourquoi, dans le prisme, le violet est plus dévié que le rouge : l'indice du verre dépend de la longueur d'onde.

\courssub{Exercices résolus}

\begin{exoresolu}[Fréquence d'un laser He-Ne]
\textbf{Énoncé.} Un laser hélium-néon émet une radiation rouge de longueur d'onde dans le vide $\lambda_0 = 633\ \text{nm}$.
\begin{enumerate}
\item Calculer la fréquence $\nu$ de cette radiation.
\item Ce faisceau traverse un bloc de verre d'indice $n = 1{,}50$. La fréquence change-t-elle ? Calculer la longueur d'onde $\lambda$ de la radiation dans le verre.
\end{enumerate}
\tcblower
\textbf{Solution détaillée.}
\begin{enumerate}
\item Conversion préalable : $\lambda_0 = 633\ \text{nm} = 633 \times 10^{-9}\ \text{m} = 6{,}33 \times 10^{-7}\ \text{m}$.
\[ \nu = \dfrac{c}{\lambda_0} = \dfrac{3{,}00 \times 10^{8}}{6{,}33 \times 10^{-7}} \approx 4{,}74 \times 10^{14}\ \text{Hz}. \]
\item La fréquence est une caractéristique invariante de la radiation : elle \textbf{ne change pas} dans le verre. C'est la longueur d'onde qui est modifiée :
\[ \lambda = \dfrac{\lambda_0}{n} = \dfrac{633\ \text{nm}}{1{,}50} = 422\ \text{nm}. \]
Dans le verre, la radiation a une longueur d'onde de 422 nm, mais elle garde sa fréquence de $4{,}74 \times 10^{14}$ Hz et donc sa couleur rouge.
\end{enumerate}
\end{exoresolu}

\begin{exoresolu}[Longueur d'onde à partir de la fréquence]
\textbf{Énoncé.} Une DEL bleue émet une radiation quasi monochromatique de fréquence $\nu = 6{,}38 \times 10^{14}\ \text{Hz}$.
\begin{enumerate}
\item Calculer sa longueur d'onde dans le vide $\lambda_0$, en m puis en nm.
\item Cette radiation appartient-elle au domaine visible ? Si oui, à quelle couleur correspond-elle ?
\end{enumerate}
\tcblower
\textbf{Solution détaillée.}
\begin{enumerate}
\item On inverse la relation fondamentale :
\[ \lambda_0 = \dfrac{c}{\nu} = \dfrac{3{,}00 \times 10^{8}}{6{,}38 \times 10^{14}} \approx 4{,}70 \times 10^{-7}\ \text{m} = 470\ \text{nm}. \]
\item Comme $400\ \text{nm} < 470\ \text{nm} < 800\ \text{nm}$, la radiation appartient bien au visible. D'après le tableau de correspondance, 470 nm se situe dans le domaine du \textbf{bleu} (430 -- 490 nm), ce qui est cohérent avec l'aspect de la DEL.
\end{enumerate}
\end{exoresolu}

\begin{methode}[Exploiter la relation $\lambda_0 = c/\nu$ sans se tromper]
\begin{enumerate}
\item \textbf{Identifier} la grandeur connue ($\lambda_0$ ou $\nu$) et la grandeur cherchée.
\item \textbf{Convertir} systématiquement les longueurs d'onde en mètres : $1\ \text{nm} = 10^{-9}\ \text{m}$.
\item \textbf{Appliquer} $\nu = c/\lambda_0$ ou $\lambda_0 = c/\nu$ avec $c = 3{,}00 \times 10^{8}\ \text{m}\cdot\text{s}^{-1}$.
\item \textbf{Vérifier} la cohérence du résultat : les fréquences lumineuses visibles sont de l'ordre de $10^{14}$ à $10^{15}$ Hz ; les longueurs d'onde visibles, de 400 à 800 nm.
\item \textbf{Conclure} avec l'unité SI et un nombre raisonnable de chiffres significatifs (trois en général).
\end{enumerate}
\end{methode}

\begin{exercice}[Pour t'entraîner]
\begin{enumerate}
\item Un laser vert de pointeur émet une radiation de longueur d'onde $\lambda_0 = 532\ \text{nm}$. Calculer sa fréquence.
\item Une station de radio communautaire de Bafoussam émet à la fréquence $\nu = 98{,}5\ \text{MHz}$. Calculer la longueur d'onde dans le vide de cette onde radio. Appartient-elle au domaine visible ?
\item Le spectre d'une lampe à vapeur de mercure présente des raies isolées à 405 nm (violet), 546 nm (vert) et 578 nm (jaune). Cette lampe est-elle monochromatique ou polychromatique ? Justifier.
\end{enumerate}
\end{exercice}

\begin{corrige}[Exercice]
\begin{enumerate}
\item $\nu = \dfrac{c}{\lambda_0} = \dfrac{3{,}00 \times 10^{8}}{532 \times 10^{-9}} \approx 5{,}64 \times 10^{14}\ \text{Hz}$.
\item $\lambda_0 = \dfrac{c}{\nu} = \dfrac{3{,}00 \times 10^{8}}{98{,}5 \times 10^{6}} \approx 3{,}05\ \text{m}$. Cette longueur d'onde (environ 3 mètres !) est immensément plus grande que 800 nm : il ne s'agit pas du tout de lumière visible, mais d'une onde radio, invisible pour l'œil.
\item La lampe émet \emph{plusieurs} radiations distinctes (trois raies) : elle est donc \textbf{polychromatique}, avec un spectre discontinu (spectre de raies). Une source monochromatique n'émettrait qu'une seule raie.
\end{enumerate}
\end{corrige}

\begin{savaistu}[Le sodium, des rues de Douala aux étoiles]
La radiation jaune à 589 nm des lampes à vapeur de sodium est aussi celle que les astronomes observent dans la lumière des étoiles contenant du sodium : l'analyse spectrale permet d'identifier la composition chimique des astres sans jamais y aller ! Mieux : des lasers accordés sur 589 nm sont tirés vers la haute atmosphère depuis les grands observatoires pour exciter les atomes de sodium présents à 90 km d'altitude et créer une « étoile artificielle » servant à corriger la turbulence de l'air. La même physique que celle des lampadaires, au service des télescopes géants.
\end{savaistu}

\begin{aretenir}[L'essentiel de la section]
\begin{itemize}
\item Une source \textbf{monochromatique} émet une seule radiation, caractérisée par sa longueur d'onde dans le vide $\lambda_0$ (ex. : laser He-Ne, 633 nm).
\item Une source \textbf{polychromatique} émet plusieurs radiations : spectre continu (lumière blanche) ou de raies (vapeur de mercure).
\item Relation fondamentale : $\lambda_0 = \dfrac{c}{\nu}$ avec $c = 3{,}00 \times 10^{8}\ \text{m}\cdot\text{s}^{-1}$.
\item La fréquence $\nu$ caractérise la radiation et est \textbf{invariante} d'un milieu à l'autre ; dans un milieu d'indice $n$, $\lambda = \lambda_0/n$.
\item Le visible s'étend de 400 nm (violet) à 800 nm (rouge) ; toujours convertir les nm en m avant tout calcul.
\end{itemize}
\end{aretenir}

\coursec{Analyse des profils spectraux}

Dans la section précédente, nous avons appris à distinguer une source monochromatique (le laser rouge, qui n'émet qu'une seule longueur d'onde $\lambda_0 \approx \SI{633}{nm}$) d'une source polychromatique (le Soleil, la lampe à incandescence, qui émettent un mélange de radiations). Mais dire qu'une lumière est « polychromatique » reste vague : quelles radiations contient-elle exactement, et en quelles proportions ? Pour répondre à cette question, les physiciens ont inventé un outil remarquable : le \cle{profil spectral}. C'est en quelque sorte la « radiographie » d'une lumière. Cette section t'apprend à le lire, à l'interpréter, et à t'en servir pour identifier une source lumineuse ou même une espèce chimique.

\courssub{Le profil spectral : définition}

\begin{experience}[Une observation pour démarrer]
On décompose la lumière d'une lampe à incandescence avec un prisme : on obtient un arc-en-ciel continu. Mais toutes les couleurs ne sont pas également brillantes : le rouge-orangé y est éclatant, le bleu-violet beaucoup plus faible. Avec une lampe à vapeur de sodium (les lampadaires orangés qu'on rencontre encore sur certains axes routiers camerounais), on n'observe presque rien... sauf une raie jaune très intense. Il ne suffit donc pas de savoir \emph{quelles} couleurs sont présentes : il faut aussi mesurer \emph{combien} de chaque couleur la source émet. Un appareil appelé \cle{spectrophotomètre} (ou spectromètre) fait exactement cela : il mesure, pour chaque longueur d'onde $\lambda$, l'intensité lumineuse reçue.
\end{experience}

\begin{definition}[Profil spectral]
Le \cle{profil spectral} d'une source de lumière est la courbe donnant l'\textbf{intensité lumineuse relative} $I$ émise par la source en fonction de la \textbf{longueur d'onde} $\lambda$ (généralement exprimée en nanomètres, nm).
\begin{itemize}
\item L'intensité est dite « relative » car on la rapporte à l'intensité maximale : $I$ varie alors entre $0$ et $1$ (ou entre $0$ et $100\,\%$).
\item Le profil spectral est la \textbf{carte d'identité lumineuse} de la source : deux sources de profils différents sont de natures différentes.
\end{itemize}
\end{definition}

\begin{attention}[Intensité relative, pas absolue]
Sur un profil spectral, l'axe vertical ne donne pas une puissance en watts mais une grandeur sans unité, normalisée. Une lampe de poche et un projecteur de stade peuvent avoir le \emph{même} profil spectral (même forme de courbe) tout en ayant des puissances très différentes : le profil renseigne sur la \textbf{composition} de la lumière, pas sur sa \textbf{puissance totale}.
\end{attention}

\courssub{Le spectre continu : le cas des sources thermiques}

Chauffe un clou dans la flamme d'un bec Bunsen : il devient rouge, puis orange, puis jaune-orangé. Le filament de tungstène d'une lampe à incandescence, porté à environ \SI{2700}{K}, et la surface du Soleil, à environ \SI{5800}{K}, se comportent de la même façon : ce sont des \cle{sources thermiques}. Leur émission est due à l'agitation thermique des charges (électrons, ions) dans la matière.

\begin{propriete}[Profil spectral d'une source thermique]
Une source thermique (corps porté à haute température : filament, surface d'étoile, métal incandescent) émet un \textbf{spectre continu} : son profil spectral est une courbe \textbf{en cloche}, sans discontinuité, qui s'étend sur tout le visible (et bien au-delà, dans l'infrarouge et l'ultraviolet). Elle présente :
\begin{itemize}
\item un \textbf{maximum} d'émission pour une longueur d'onde $\lambda_{\max}$ qui dépend de la température du corps (nous y reviendrons avec la loi de Wien, section 5) ;
\item des intensités qui décroissent progressivement de part et d'autre du maximum.
\end{itemize}
Plus le corps est chaud, plus le maximum se déplace vers les courtes longueurs d'onde (vers le bleu).
\end{propriete}

\begin{popfigure}
\begin{center}
\begin{tikzpicture}
\begin{axis}[
width=0.9\linewidth, height=7cm,
xlabel={$\lambda$ (nm)}, ylabel={$I$ (intensité relative)},
xmin=300, xmax=1100, ymin=0, ymax=1.1,
xtick={400,500,600,700,800,900,1000},
ytick={0,0.5,1},
grid=both, grid style={popline!40},
legend style={font=\small, at={(0.98,0.98)}, anchor=north east},
clip=false
]
% Domaine visible ombragé
\addplot[popblueL!30, domain=380:780] {1.1} \closedcycle;
\node[font=\small, popblue, rotate=90] at (axis cs:380,0.5) {Visible};

% Courbe Soleil T=5800K -> lambda_max ~ 500 nm (Wien)
\addplot[popcurve, domain=300:1100, samples=300, poporange, very thick]
{exp(-((x-500)/130)^2)};
\addlegendentry{Soleil ($T \approx \SI{5800}{K}$)}

% Courbe Incandescence T=2700K -> lambda_max ~ 1070 nm (IR)
% Dans le visible (380-780), c'est le flanc montant de la cloche.
\addplot[popcurve, domain=300:1100, samples=300, popblue, very thick]
{exp(-((x-1070)/350)^2)};
\addlegendentry{Incandescence ($T \approx \SI{2700}{K}$)}

% Marqueurs lambda_max
\draw[dashed, poporange] (axis cs:500,0) -- (axis cs:500,1.02);
\node[font=\small, poporange] at (axis cs:500,1.08) {$\lambda_{\max,\odot}$};
\draw[dashed, popblue] (axis cs:1070,0) -- (axis cs:1070,1.02);
\node[font=\small, popblue] at (axis cs:1070,1.08) {$\lambda_{\max,\text{inc}}$ (IR)};

\end{axis}
\end{tikzpicture}
\end{center}
\legende{Profils spectraux continus de deux sources thermiques. Le Soleil ($5800\,\text{K}$) a son maximum dans le visible (vert). L'incandescence ($2700\,\text{K}$) a son maximum dans l'infrarouge ($\approx 1070\,\text{nm}$) : dans le visible, on ne voit que le flanc montant de la cloche (prédominance rouge-orangé). La zone bleue clair délimite le domaine visible.}
\end{popfigure}

\courssub{Le spectre de raies d'émission : la signature des atomes}

Changeons radicalement de source. Une \cle{lampe spectrale} contient un gaz (vapeur de sodium, de mercure, d'hydrogène...) à basse pression, traversé par une décharge électrique. Les atomes du gaz, excités par les collisions avec les électrons, se désexcitent en émettant de la lumière. Mais — fait capital — ils n'émettent pas n'importe quelles longueurs d'onde.

\begin{propriete}[Spectre de raies d'émission]
Le profil spectral d'une lampe à vapeur atomique est constitué de \textbf{raies fines et intenses sur fond noir} : l'intensité est nulle partout, sauf pour quelques longueurs d'onde précises appelées \cle{raies d'émission}.\\
\textbf{Chaque élément chimique possède un spectre de raies qui lui est propre} : c'est sa « carte d'identité » lumineuse. Deux éléments différents n'ont jamais le même ensemble de raies.
\end{propriete}

L'explication profonde (que tu découvriras en Terminale avec le modèle de Bohr) est que l'énergie d'un atome est \emph{quantifiée} : l'atome ne peut perdre son énergie que par « sauts » bien définis entre niveaux d'énergie discrets, et chaque saut correspond à une radiation de longueur d'onde précise, donnée par $\lambda = \dfrac{hc}{\Delta E}$. Comme les niveaux d'énergie dépendent du noyau et du nuage électronique, chaque élément a ses propres sauts, donc ses propres raies.

\begin{exemplebox}[Le doublet jaune du sodium]
L'atome de sodium émet deux raies jaunes très proches, de longueurs d'onde $\lambda_1 = \SI{589,0}{nm}$ et $\lambda_2 = \SI{589,6}{nm}$ : c'est le fameux \textbf{doublet du sodium}. À l'œil nu, ces deux raies se confondent en une seule raie jaune intense — d'où la couleur orangée caractéristique des lampadaires au sodium. Un bon spectroscope (ou un réseau de diffraction) les sépare nettement.
\end{exemplebox}

\begin{popfigure}
\begin{center}
\begin{tikzpicture}
\begin{axis}[
width=0.9\linewidth, height=6cm,
xlabel={$\lambda$ (nm)}, ylabel={$I$ (intensité relative)},
xmin=580, xmax=600, ymin=0, ymax=1.15,
xtick={580,585,589,590,595,600}, ytick={0,0.5,1},
grid=both, grid style={popline!40},
tick label style={font=\small}
]
% Doublet sodium : deux gaussiennes étroites centrées sur 589.0 et 589.6
\addplot[popcurve, domain=580:600, samples=400, popgold, very thick]
{exp(-((x-589.0)/0.3)^2) + 0.95*exp(-((x-589.6)/0.3)^2)};
\draw[dashed, popdark] (axis cs:589.0,0) -- (axis cs:589.0,1.02);
\draw[dashed, popdark] (axis cs:589.6,0) -- (axis cs:589.6,1.02);
\node[font=\small, popdark, align=center] at (axis cs:595,0.8) {Doublet Na\\ $\lambda_1=\SI{589,0}{nm}$\\ $\lambda_2=\SI{589,6}{nm}$};
\end{axis}
\end{tikzpicture}
\end{center}
\legende{Profil spectral d'une lampe à vapeur de sodium (zoom sur le jaune) : deux raies fines et intenses (le doublet) sur un fond pratiquement nul.}
\end{popfigure}

\begin{popfigure}
\begin{center}
\begin{tikzpicture}[x=0.016cm]
% Fond noir
\fill[black] (380,0) rectangle (750,3.2);
% Echelle
\foreach \x in {400,450,500,550,600,650,700}{
  \draw[white] (\x,0) -- (\x,-0.15);
  \node[white, font=\scriptsize, below] at (\x,-0.15) {\x};
}
\node[white, font=\scriptsize, below] at (565,-0.75) {$\lambda$ (nm)};
% Raies du sodium (doublet 589)
\fill[yellow] (588.6,0.4) rectangle (589.4,1.3);
\fill[yellow] (589.2,0.4) rectangle (590.0,1.3);
\node[white, font=\small] at (380,0.85) [anchor=west] {\textbf{Na}};
% Raies du mercure : 405 violet, 436 bleu, 546 vert, 578 jaune
\fill[poppurple] (404.5,1.9) rectangle (405.5,2.8);
\fill[popblue] (435.5,1.9) rectangle (436.5,2.8);
\fill[popgreen] (545.5,1.9) rectangle (546.5,2.8);
\fill[yellow] (577.5,1.9) rectangle (579.5,2.8);
\node[white, font=\small] at (380,2.35) [anchor=west] {\textbf{Hg}};
\end{tikzpicture}
\end{center}
\legende{Spectres de raies d'émission du sodium (Na) et du mercure (Hg) : raies colorées sur fond noir. Les positions des raies sont différentes : chaque élément a sa signature.}
\end{popfigure}

\courssub{Le spectre de bandes : DEL et tubes fluorescents}

Entre le continu « parfait » de l'incandescence et les raies ultra-fines des lampes spectrales, il existe un cas intermédiaire : le \cle{spectre de bandes}. Une \textbf{bande} est un domaine de longueurs d'onde assez large (quelques dizaines de nanomètres) sur lequel l'intensité est notable, formant une « bosse » dans le profil.

\begin{itemize}
\item Une \textbf{DEL} (diode électroluminescente) colorée émet une bande assez étroite centrée sur une longueur d'onde donnée : elle est « presque monochromatique », mais pas tout à fait — contrairement au laser dont la raie est infiniment plus fine.
\item Une \textbf{DEL blanche} est en réalité une DEL bleue (InGaN) recouverte d'un luminophore (phosphore YAG:Ce) qui convertit une partie de la lumière bleue en lumière jaune (large bande). Son profil présente donc \textbf{deux bosses} : un pic étroit dans le bleu (vers $\SI{450}{nm}$) et une large bosse dans le jaune-vert (vers $\SI{560}{nm}$). Le mélange bleu + jaune est perçu comme blanc par notre œil (synthèse additive).
\item Les \textbf{tubes fluorescents} (les « néons » des salles de classe et bureaux) combinent des raies fines (vapeur de mercure à $\SI{254}{nm}$ UV et raies visibles $\SI{436}{nm}$, $\SI{546}{nm}$...) et des bandes larges (poudres fluorescentes — tri-phosphores — déposées sur la paroi qui convertissent l'UV en rouge, vert, bleu).
\end{itemize}

\begin{popfigure}
\begin{center}
\begin{tikzpicture}
\begin{axis}[
width=0.9\linewidth, height=6cm,
xlabel={$\lambda$ (nm)}, ylabel={$I$ (intensité relative)},
xmin=380, xmax=750, ymin=0, ymax=1.15,
xtick={400,450,500,550,600,700}, ytick={0,0.5,1},
grid=both, grid style={popline!40}
]
% DEL blanche : pic bleu étroit (FWHM ~20nm) + bosse jaune large (FWHM ~100nm)
\addplot[popcurve, domain=380:750, samples=300, popblue, very thick]
{exp(-((x-450)/10)^2) + 0.65*exp(-((x-560)/50)^2)};
\draw[dashed, popdark] (axis cs:450,0) -- (axis cs:450,1.02);
\draw[dashed, popdark] (axis cs:560,0) -- (axis cs:560,0.68);
\node[font=\small, popdark] at (axis cs:450,1.1) {Pic bleu ($\sim\SI{450}{nm}$)};
\node[font=\small, popdark] at (axis cs:560,0.75) {Bosse jaune ($\sim\SI{560}{nm}$)};
\end{axis}
\end{tikzpicture}
\end{center}
\legende{Profil spectral typique d'une DEL blanche : un pic étroit dans le bleu (émetteur InGaN) et une large bosse dans le jaune (luminophore YAG:Ce).}
\end{popfigure}

\courssub{Lecture et exploitation d'un profil spectral}

\begin{methode}[Analyser le profil spectral d'une source]
Face à un profil spectral, procède toujours dans le même ordre :
\begin{enumerate}
\item \textbf{Observer la forme globale} : courbe en cloche continue (source thermique), raies fines isolées sur fond nul (vapeur atomique), ou bandes/bosses (DEL, tube fluorescent) ?
\item \textbf{Repérer le maximum} : noter la longueur d'onde $\lambda_{\max}$ pour laquelle l'intensité est la plus grande. Pour une source thermique, elle renseigne sur la température (loi de Wien, section suivante).
\item \textbf{Estimer la largeur} : une raie très fine (largeur $\ll \SI{1}{nm}$) indique une émission quasi monochromatique (laser, lampe spectrale) ; une bosse large (quelques nm à dizaines de nm) indique une bande (DEL, fluorescence).
\item \textbf{Relever les raies caractéristiques} s'il y en a : comparer leurs longueurs d'onde à des tables de raies connues pour identifier l'élément chimique émetteur.
\item \textbf{Conclure} sur la nature de la source (thermique, spectrale, DEL, laser...).
\end{enumerate}
\end{methode}

\begin{propriete}[Principe de l'analyse spectrale]
Puisque chaque élément chimique possède un spectre de raies qui lui est propre, on peut \textbf{identifier une espèce chimique} présente dans une source lumineuse en comparant les raies observées aux raies de référence des éléments connus. C'est le principe de l'\cle{analyse spectrale}, utilisée en chimie (dosage), en astrophysique (composition des étoiles, décalage vers le rouge) et dans le contrôle industriel (qualité de l'air, soudure).
\end{propriete}

\begin{exoresolu}[Identifier un gaz inconnu]
\textbf{Énoncé.} On réalise le spectre de la lumière émise par une lampe contenant un gaz inconnu. On observe trois raies intenses aux longueurs d'onde $\lambda_1 = \SI{405}{nm}$, $\lambda_2 = \SI{436}{nm}$ et $\lambda_3 = \SI{546}{nm}$. Les tables donnent : sodium : $\SI{589}{nm}$ (doublet) ; mercure : $\SI{405}{nm}$, $\SI{436}{nm}$, $\SI{546}{nm}$, $\SI{578}{nm}$ ; hydrogène : $\SI{410}{nm}$, $\SI{434}{nm}$, $\SI{486}{nm}$, $\SI{656}{nm}$. Identifier le gaz.
\tcblower
\textbf{Solution.} Appliquons la méthode :
\begin{enumerate}
\item Forme : raies fines sur fond noir $\Rightarrow$ il s'agit d'une vapeur atomique excitée (lampe spectrale).
\item Comparaison des raies observées aux tables de référence :
\begin{itemize}
\item Sodium : sa raie principale (doublet à $\SI{589}{nm}$) est absente du spectre observé. Ce n'est pas le sodium.
\item Hydrogène : ses raies caractéristiques (série de Balmer : $\SI{486}{nm}$ bleu-vert, $\SI{656}{nm}$ rouge) sont absentes. Ce n'est pas l'hydrogène.
\item Mercure : les trois raies observées ($\SI{405}{nm}$ violet, $\SI{436}{nm}$ bleu, $\SI{546}{nm}$ vert) figurent toutes dans son spectre de référence.
\end{itemize}
\item Conclusion : le gaz inconnu est la \textbf{vapeur de mercure}. (La raie à $\SI{578}{nm}$ jaune est peut-être simplement trop faible pour avoir été détectée avec le matériel utilisé : l'absence d'une raie secondaire ne contredit pas l'identification, mais la présence des raies principales la confirme solidement.)
\end{enumerate}
\end{exoresolu}

\begin{attention}[Raies d'émission et raies d'absorption : ne pas confondre]
Un spectre de raies d'émission n'est \textbf{pas} un spectre continu « troué » : c'est un fond \emph{noir} (obscurité) sur lequel se détachent quelques raies brillantes. Il ne faut pas le confondre avec le \textbf{spectre d'absorption} : lorsqu'une lumière blanche (spectre continu, source thermique) traverse un gaz \emph{froid}, celui-ci absorbe certaines radiations correspondant à ses transitions possibles. Le spectre obtenu est alors un \emph{fond continu coloré} traversé de \emph{raies noires} (lignes d'absorption). Fait remarquable (loi de Kirchhoff) : un élément absorbe exactement les longueurs d'onde qu'il est capable d'émettre. \textbf{Raies brillantes sur fond noir = émission} ; \textbf{raies noires sur fond continu = absorption}.
\end{attention}

\courssub{Application : comparer trois sources courantes}

Mettons en pratique tout ce qui précède en comparant les profils spectraux de trois sources que tu croises chaque jour.

\begin{center}
\begin{tabularx}{\linewidth}{|l|X|X|X|}
\hline
\rowcolor{popblueL} \textbf{Critère} & \textbf{Incandescence / Halogène} & \textbf{DEL blanche} & \textbf{Laser (ex. He-Ne)} \\
\hline
Forme du profil & Cloche continue (corps noir) & Pic bleu étroit + bosse jaune large & Raie unique ultra-fine \\
\hline
$\lambda_{\max}$ (max global) & $\approx \SI{1000}{nm}$ (Infrarouge !) & $\approx \SI{450}{nm}$ (Pic bleu) & $\lambda_0$ unique (ex. $\SI{633}{nm}$) \\
\hline
Nature du spectre & Polychromatique (continu) & Polychromatique (bandes) & Quasi monochromatique \\
\hline
Rendement lumineux & Faible ($\sim \SI{15}{lm/W}$) & Élevé ($\sim \SI{100}{lm/W}$) & N/A (spécifique) \\
\hline
\end{tabularx}
\end{center}

Un point frappe immédiatement : pour la lampe à incandescence, le maximum d'émission se situe dans l'\textbf{infrarouge} ($\lambda_{\max} \approx \SI{1000}{nm}$ pour $T=\SI{2700}{K}$), c'est-à-dire \emph{en dehors} du domaine visible ! C'est pourquoi ces lampes chauffent énormément et ont un si mauvais rendement lumineux — raison pour laquelle elles sont progressivement remplacées par les DEL, y compris dans les villes et villages du Cameroun où l'économie d'énergie est précieuse (programme d'électrification rurale, éclairage public SONATREL/ENEO). Le laser, lui, concentre toute son énergie sur une seule longueur d'onde : son profil est une raie si fine (largeur $\Delta\lambda \ll \SI{1}{nm}$) qu'on le considère comme parfaitement monochromatique.

\begin{savaistu}[L'hélium : l'élément découvert dans le Soleil]
En 1868, lors d'une éclipse totale de Soleil observée en Inde, l'astronome français Jules Janssen analyse le spectre de la couronne solaire (prominences) et y remarque une raie jaune brillante à $\SI{587,5}{nm}$... qui ne correspond à \emph{aucun} élément connu sur Terre ! L'astronome anglais Norman Lockyer, étudiant ce spectre, en conclut qu'il s'agit d'un élément nouveau, qu'il baptise \textbf{hélium} (du grec \emph{hélios}, Soleil). Il faudra attendre 1895 pour que le chimiste William Ramsay isole cet élément sur Terre, dans un minéral d'uranium (cléveïte). L'analyse spectrale avait donc permis de découvrir un élément chimique à 150 millions de kilomètres, avant même qu'on le tienne en main : une des plus belles victoires de la physique et de la chimie !
\end{savaistu}

\begin{exercice}[Lecture de profil spectral]
Le profil spectral d'une ampoule présente une courbe en cloche continue dont le maximum se situe à $\lambda_{\max} = \SI{580}{nm}$, sans aucune raie fine.
\begin{enumerate}
\item Quelle est la nature de cette source ?
\item Cette ampoule est-elle une lampe à incandescence classique ou une DEL blanche ? Justifier.
\end{enumerate}
\end{exercice}

\begin{corrige}[Exercice : lecture de profil spectral]
\begin{enumerate}
\item Le profil est une cloche \textbf{continue} : c'est la signature d'une \textbf{source thermique} (corps chauffé à haute température).
\item Une DEL blanche présenterait un pic étroit dans le bleu ($\sim \SI{450}{nm}$) et une bosse large dans le jaune ($\sim \SI{560}{nm}$), donc un profil \emph{non continu} avec deux maxima locaux distincts. Ici, la cloche est unique, continue, avec $\lambda_{\max} = \SI{580}{nm}$ dans le visible (jaune-orangé) : il s'agit d'une \textbf{lampe à incandescence} (ou halogène) relativement chaude ($T \approx \SI{5000}{K}$ d'après Wien, ou filament surchauffé). La DEL blanche est exclue par la forme du profil.
\end{enumerate}
\end{corrige}

\begin{aretenir}[L'essentiel de la section]
\begin{itemize}
\item Le \textbf{profil spectral} donne l'intensité lumineuse relative en fonction de la longueur d'onde $\lambda$ (nm).
\item \textbf{Spectre continu en cloche} (avec un maximum en $\lambda_{\max}$) : sources thermiques (Soleil, incandescence, halogène). $\lambda_{\max}$ dépend de la température (loi de Wien).
\item \textbf{Raies fines sur fond noir} : lampes à vapeur atomique (Na, Hg, H, He...) ; chaque élément a son spectre propre $\rightarrow$ \textbf{analyse spectrale} (identification chimique).
\item \textbf{Bandes} (bosses larges, quelques dizaines de nm) : DEL colorées, DEL blanche (pic bleu + bosse jaune), tubes fluorescents (raies Hg + bandes phosphores).
\item \textbf{Laser} : raie unique ultra-fine (quasi monochromatique, cohérent).
\item Méthode d'analyse : Forme globale $\rightarrow$ $\lambda_{\max}$ $\rightarrow$ Largeur/Raies $\rightarrow$ Identification de la source.
\end{itemize}
\end{aretenir}

\coursec{Les ondes électromagnétiques}

Dans la section précédente, nous avons analysé les \cle{profils spectraux} : nous avons appris à décomposer la lumière, à repérer des raies d'émission ou d'absorption et à identifier des éléments chimiques. Mais une question fondamentale reste en suspens : qu'est-ce que la lumière, au juste ? Et surtout, le spectre que nous observons entre le rouge et le violet est-il tout ce qui existe ? La réponse est non : la lumière visible n'est qu'une infime partie d'une immense famille d'ondes, les \cle{ondes électromagnétiques}, qui vont des rayons gamma aux ondes radio. C'est cette famille que nous allons explorer.

\courssub{La lumière : une onde électromagnétique}

\textbf{Observation.} Le Soleil est situé à environ 150 millions de kilomètres de la Terre. Entre lui et nous, il n'y a \emph{rien} : le vide interplanétaire. Pourtant, sa lumière nous parvient, nous éclaire et nous chauffe. Contrairement au son, qui a besoin d'un milieu matériel (air, eau, solide) pour se propager, la lumière se propage donc \textbf{dans le vide}. C'est une onde d'une nature particulière : elle est constituée d'un champ électrique et d'un champ magnétique qui oscillent et se propagent ensemble, sans support matériel. D'où son nom : \cle{onde électromagnétique}.

\begin{definition}[Onde électromagnétique]
Une \cle{onde électromagnétique} est une onde constituée de la propagation couplée d'un champ électrique $\vec{E}$ et d'un champ magnétique $\vec{B}$ oscillants, perpendiculaires entre eux et perpendiculaires à la direction de propagation. Elle se propage \textbf{dans le vide} (contrairement aux ondes mécaniques comme le son) à la célérité $c$. Elle est caractérisée par :
\begin{itemize}
\item sa \cle{longueur d'onde dans le vide} $\lambda_0$, exprimée en mètres (m) ;
\item sa \cle{fréquence} $\nu$ (lettre grecque « nu »), exprimée en hertz (Hz).
\end{itemize}
Ces deux grandeurs sont liées par la relation fondamentale :
\[ \lambda_0 = \frac{c}{\nu} \qquad \text{soit} \qquad \nu = \frac{c}{\lambda_0} \]
\end{definition}

\begin{propriete}[Célérité de la lumière dans le vide]
Dans le vide, \textbf{toutes} les ondes électromagnétiques — de la lumière visible aux ondes radio en passant par les rayons X — se propagent à la même célérité :
\[ c = \num{3,00e8}~\text{m/s} \quad (\text{soit } \num{300 000}~\text{km/s}) \]
C'est une constante universelle, l'une des plus importantes de toute la physique. Aucun objet matériel, aucun signal ne peut la dépasser.
\end{propriete}

\textbf{Vérification par analyse dimensionnelle.} La relation $\lambda_0 = c/\nu$ est-elle homogène ? $c$ s'exprime en m/s et $\nu$ en Hz, c'est-à-dire en s$^{-1}$. Donc :
\[ \frac{[c]}{[\nu]} = \frac{\text{m}\cdot\text{s}^{-1}}{\text{s}^{-1}} = \text{m} \]
On obtient bien des mètres, unité d'une longueur d'onde. La relation est homogène.

\begin{attention}[La fréquence est une caractéristique de la source]
La fréquence $\nu$ d'une radiation est fixée par la source qui l'émet et \textbf{ne change jamais}, quel que soit le milieu traversé. En revanche, la longueur d'onde dépend du milieu : dans un milieu transparent d'indice $n$, la célérité devient $v = c/n$ et la longueur d'onde devient $\lambda = \lambda_0/n$. C'est pourquoi on précise toujours « longueur d'onde \emph{dans le vide} », notée $\lambda_0$.
\end{attention}

\begin{exemplebox}[La lumière du Soleil met du temps à nous parvenir]
La distance Terre--Soleil vaut environ $d = \num{1,5e11}$ m. La lumière met :
\[ \Delta t = \frac{d}{c} = \frac{\num{1,5e11}}{\num{3,00e8}} = \num{500}~\text{s} \approx 8~\text{min}~20~\text{s} \]
pour nous parvenir. Autrement dit, tu vois le Soleil \emph{tel qu'il était il y a plus de 8 minutes} ! Si le Soleil s'éteignait brutalement, nous continuerions à le voir briller pendant encore 8 minutes.
\end{exemplebox}

\courssub{Les domaines des ondes électromagnétiques}

\textbf{Intuition.} Une même corde de guitare peut produire des notes graves ou aiguës selon sa tension : ce sont toujours des ondes sonores, mais de fréquences différentes. De même, les ondes électromagnétiques forment une \emph{grande famille unique} : toutes ont la même nature (champs $\vec{E}$ et $\vec{B}$ oscillants), toutes se propagent dans le vide à la célérité $c$. Ce qui les distingue, c'est uniquement leur \cle{longueur d'onde} $\lambda_0$ (ou, de manière équivalente, leur fréquence $\nu$, puisque $\nu = c/\lambda_0$).

On classe les ondes électromagnétiques par \textbf{longueur d'onde croissante} (c'est-à-dire par fréquence \emph{décroissante}) :

\begin{aretenir}[Les sept domaines, par longueur d'onde croissante]
\[ \text{rayons } \gamma \;<\; \text{rayons X} \;<\; \text{UV} \;<\; \text{visible} \;<\; \text{IR} \;<\; \text{micro-ondes} \;<\; \text{ondes radio} \]
Plus la longueur d'onde est \textbf{petite}, plus la fréquence est \textbf{grande} et plus l'onde est \textbf{énergétique} (et potentiellement dangereuse).
\end{aretenir}

\begin{center}
\begin{tabularx}{\linewidth}{|l|l|X|}
\hline
\rowcolor{popblueL} \textbf{Domaine} & \textbf{Ordre de grandeur de $\lambda_0$} & \textbf{Applications} \\
\hline
Rayons $\gamma$ & $\lambda_0 < 10^{-12}$ m & Radiothérapie (traitement des cancers, hôpitaux de Yaoundé et Douala), stérilisation médicale \\
\hline
Rayons X & $10^{-12}$ m à $10^{-9}$ m & Radiographie médicale, scanners, contrôle des bagages dans les aéroports \\
\hline
Ultraviolet (UV) & $10^{-9}$ m à $400$ nm & Stérilisation de l'eau et du matériel, fluorescence (détection de faux billets), bronzage \\
\hline
Visible & $400$ nm à $800$ nm & Vision, éclairage, photographie, fibres optiques \\
\hline
Infrarouge (IR) & $800$ nm à $1$ mm & Télécommandes, caméras thermiques, chauffage, vision nocturne \\
\hline
Micro-ondes & $1$ mm à $30$ cm & Four à micro-ondes, téléphonie mobile (MTN, Orange), satellites, radar \\
\hline
Ondes radio & $30$ cm à plusieurs km & Radio (CRTV, FM), télévision, communications, radioastronomie \\
\hline
\end{tabularx}
\end{center}

\begin{popfigure}
\begin{tikzpicture}[x=1cm,y=1cm]
% Axe logarithmique de 10^-12 a 10^3 (15 decades, 0.9 cm par decade)
\draw[->, thick, popdark] (0,0) -- (14.2,0);
% Graduations
\foreach \i/\p in {0/$10^{-12}$, 1.8/$10^{-10}$, 3.6/$10^{-8}$, 5.4/$10^{-6}$, 7.2/$10^{-4}$, 9/$10^{-2}$, 10.8/$10^{0}$, 12.6/$10^{2}$} {
  \draw[thick, popdark] (\i,0) -- (\i,0.15);
  \node[below, font=\scriptsize, popdark] at (\i,-0.08) {\p};
}
\node[below, font=\small, popdark] at (13.6,-0.35) {$\lambda_0$ (m)};
% Bandes de domaines (echelle log : position = (log10(lambda)+12)*0.9)
% gamma : < 10^-11 -> 0 a 0.9 ; X : 10^-12 a 10^-9 -> 0.9 a 2.7 ; UV : 10^-9 a 4.10^-7 -> 2.7 a 5.04
% visible : 4.10^-7 a 8.10^-7 -> 5.04 a 5.31 ; IR : 8.10^-7 a 10^-3 -> 5.31 a 8.1
% micro-ondes : 10^-3 a 0.3 -> 8.1 a 10.33 ; radio : > 0.3 -> 10.33 a 13.5
\fill[poppurple!60] (0,0.25) rectangle (0.9,0.75);
\fill[popblue!50] (0.9,0.25) rectangle (2.7,0.75);
\fill[poppink!50] (2.7,0.25) rectangle (5.04,0.75);
\fill[popgreen!80] (5.04,0.25) rectangle (5.31,0.75);
\fill[popgreen!40] (5.31,0.25) rectangle (8.1,0.75);
\fill[popgold!60] (8.1,0.25) rectangle (10.33,0.75);
\fill[poporange!55] (10.33,0.25) rectangle (13.5,0.75);
\node[font=\scriptsize\bfseries, popdark] at (0.45,0.5) {$\gamma$};
\node[font=\scriptsize\bfseries, popdark] at (1.8,0.5) {X};
\node[font=\scriptsize\bfseries, popdark] at (3.85,0.5) {UV};
\node[font=\scriptsize\bfseries, popdark] at (6.7,0.5) {IR};
\node[font=\scriptsize\bfseries, popdark] at (9.2,0.5) {Micro-ondes};
\node[font=\scriptsize\bfseries, popdark] at (11.9,0.5) {Ondes radio};
% Zoom visible : bande arc-en-ciel
\shade[left color=violet, right color=red] (4.1,-1.6) rectangle (7.1,-1.1);
\draw[thick, popdark] (4.1,-1.6) rectangle (7.1,-1.1);
\draw[->, thick, popmuted] (5.04,0.25) -- (4.3,-1.05);
\draw[->, thick, popmuted] (5.31,0.25) -- (6.9,-1.05);
\node[font=\scriptsize, popdark] at (5.6,-1.85) {Visible : de $400$ nm (violet) à $800$ nm (rouge)};
% Fleche frequence
\draw[->, thick, popink] (13.5,1.15) -- (0.5,1.15);
\node[font=\scriptsize, popink, above] at (7,1.15) {Fréquence $\nu$ croissante (énergie croissante)};
\end{tikzpicture}
\legende{Échelle logarithmique des domaines des ondes électromagnétiques. Le domaine visible (fine bande verte), zoomé en bande colorée, n'occupe qu'une fraction infime du spectre complet.}
\end{popfigure}

\begin{attention}[L'échelle est logarithmique !]
Sur l'échelle ci-dessus, chaque graduation correspond à une multiplication par 100 de la longueur d'onde. Le domaine visible ($400$ à $800$ nm) est si étroit qu'il serait \emph{invisible} sur une échelle linéaire : c'est pourquoi on le « zoome » sous forme de bande colorée. Retiens bien : le visible n'est qu'une toute petite fenêtre du spectre électromagnétique.
\end{attention}

\courssub{Zoom sur les domaines voisins du visible : UV et IR}

\textbf{Les ultraviolets (UV).} Leur nom signifie « au-delà du violet » : ce sont les ondes de longueur d'onde \emph{inférieure} à celle du violet ($\lambda_0 < 400$ nm). Le Soleil en émet abondamment. Sous le soleil tropical du Cameroun, les UV sont intenses : ce sont eux qui provoquent les coups de soleil et qui, à forte dose, endommagent la peau et les yeux. Mais ils ont aussi des applications utiles : les UV \textbf{stérilisent} l'eau et le matériel médical (ils détruisent bactéries et virus), et ils provoquent la \textbf{fluorescence} de certaines substances — c'est le principe des lampes utilisées pour vérifier l'authenticité des billets de banque.

\textbf{Les infrarouges (IR).} Leur nom signifie « au-dessous du rouge » (en fréquence) : ce sont les ondes de longueur d'onde \emph{supérieure} à celle du rouge ($\lambda_0 > 800$ nm). Tout corps chaud émet des infrarouges : c'est pourquoi on les associe à la \textbf{chaleur}. Tu sens les IR sur ta peau quand tu t'approches d'un feu de bois ou d'une braise de maquis. Applications quotidiennes : les \textbf{télécommandes} de téléviseur, les \textbf{caméras thermiques} (qui « voient » la chaleur des corps dans le noir), les lampes chauffantes.

\begin{attention}[Les UV ne sont pas violets, les IR ne sont pas rouges !]
Erreur classique au Probatoire : croire que les ultraviolets sont « de la lumière violette » et les infrarouges « de la lumière rouge ». C'est \textbf{faux} : les UV et les IR sont \emph{invisibles} pour l'œil humain. Leurs noms indiquent seulement leur \emph{position} dans le spectre, juste à côté du violet ou du rouge. Si tu « vois » une lumière violette émise par une lampe UV (boîte de nuit, détecteur de billets), c'est la petite part de lumière \emph{visible} que la lampe émet aussi, ou la fluorescence qu'elle provoque — jamais les UV eux-mêmes.
\end{attention}

\courssub{Une même nature, des longueurs d'onde différentes}

\begin{propriete}[Unité de la famille électromagnétique]
Toutes les ondes électromagnétiques — rayons $\gamma$, rayons X, UV, lumière visible, IR, micro-ondes, ondes radio — ont la \textbf{même nature} (champs électrique et magnétique oscillants) et la \textbf{même célérité dans le vide} $c = \num{3,00e8}$ m/s. Seules changent leur longueur d'onde $\lambda_0$ et leur fréquence $\nu$, liées par $\lambda_0 = c/\nu$.
\end{propriete}

\begin{methode}[Classer deux radiations]
Pour comparer deux radiations électromagnétiques :
\begin{enumerate}
\item Identifier leurs longueurs d'onde $\lambda_0$ (convertir dans la même unité, par exemple le mètre ou le nanomètre, avec $1$ nm $= 10^{-9}$ m).
\item La radiation de \textbf{plus petite} longueur d'onde est la plus \textbf{énergétique} et se situe le plus près du domaine des rayons $\gamma$.
\item Si l'on compare les fréquences, l'ordre est \textbf{inversé} (puisque $\nu = c/\lambda_0$) : la plus grande fréquence correspond à la plus petite longueur d'onde.
\end{enumerate}
\end{methode}

\begin{exoresolu}[Fréquences d'une onde radio et d'une lumière rouge]
\textbf{Énoncé.} (1) La radio CRTV émet sur la fréquence $\nu_1 = 92{,}0$ MHz. Calculer la longueur d'onde dans le vide de cette onde radio. (2) Un laser rouge a pour longueur d'onde $\lambda_2 = 633$ nm. Calculer sa fréquence. (3) Comparer les deux radiations et conclure.
\tcblower
\textbf{Solution.} (1) Pour l'onde radio : $\nu_1 = 92{,}0$ MHz $= \num{9,20e7}$ Hz.
\[ \lambda_1 = \frac{c}{\nu_1} = \frac{\num{3,00e8}}{\num{9,20e7}} \approx \num{3,26}~\text{m} \]
C'est bien une longueur d'onde de l'ordre du mètre, caractéristique des ondes radio FM.

(2) Pour le laser rouge : $\lambda_2 = 633$ nm $= \num{6,33e-7}$ m.
\[ \nu_2 = \frac{c}{\lambda_2} = \frac{\num{3,00e8}}{\num{6,33e-7}} \approx \num{4,74e14}~\text{Hz} \]

(3) Comparaison : $\lambda_1 \approx 3{,}26$ m $\gg \lambda_2 = 633$ nm. Le rapport vaut :
\[ \frac{\lambda_1}{\lambda_2} = \frac{3{,}26}{\num{6,33e-7}} \approx \num{5,2e6} \]
L'onde radio a une longueur d'onde environ \textbf{5 millions de fois plus grande} que la lumière rouge, et donc une fréquence 5 millions de fois plus petite. Pourtant, ces deux ondes ont exactement la \emph{même nature} et voyagent dans le vide à la \emph{même célérité} $c$.
\end{exoresolu}

\begin{savaistu}[De la CRTV au Soleil : une seule et même famille]
Les ondes radio émises par les émetteurs de la CRTV, les signaux de ton téléphone mobile MTN ou Orange, la chaleur d'un feu de bois, la lumière du Soleil qui éclaire les champs de cacao, les rayons X de l'hôpital central de Yaoundé et les rayons gamma utilisés en radiothérapie : tout cela, c'est \emph{la même chose} — des ondes électromagnétiques — qui ne diffèrent que par leur longueur d'onde. Quand tu zappes de la radio à la lumière du jour, tu explores le même spectre, à 15 ordres de grandeur d'intervalle ! C'est l'une des plus belles unifications de la physique, réalisée par James Clerk Maxwell au XIX\textsuperscript{e} siècle, qui a montré que la lumière est une onde électromagnétique.
\end{savaistu}

\begin{exercice}[À toi de jouer]
Un four à micro-ondes utilise des ondes de fréquence $\nu = \num{2,45}$ GHz. (1) Calculer la longueur d'onde correspondante dans le vide. (2) Situer cette radiation sur l'échelle des domaines. (3) Un rayonnement UV a une longueur d'onde de $300$ nm : lequel des deux rayonnements est le plus énergétique ? Justifier.
\end{exercice}

\begin{corrige}[Exercice]
(1) $\nu = \num{2,45}$ GHz $= \num{2,45e9}$ Hz, donc :
\[ \lambda_0 = \frac{c}{\nu} = \frac{\num{3,00e8}}{\num{2,45e9}} \approx \num{0,122}~\text{m} \approx 12~\text{cm} \]
(2) $\lambda_0 \approx 12$ cm se situe bien dans le domaine des micro-ondes (1 mm à 30 cm). (3) Comparons les longueurs d'onde :
\[ \frac{\lambda_{\text{micro-onde}}}{\lambda_{\text{UV}}} = \frac{0{,}122}{\num{3,00e-7}} \approx \num{4,1e5} \]
L'UV ($300$ nm) a une longueur d'onde environ $400\,000$ fois plus petite que la micro-onde : sa fréquence est donc $400\,000$ fois plus grande. C'est le rayonnement \textbf{UV} qui est le plus énergétique — c'est d'ailleurs pour cela qu'il est dangereux pour la peau, alors que les micro-ondes du four ne font qu'agiter les molécules d'eau des aliments.
\end{corrige}

\textbf{En résumé.} La lumière visible n'est qu'une fenêtre étroite ($400$ à $800$ nm) dans l'immense spectre des ondes électromagnétiques. Toutes ces ondes partagent la même nature et la même célérité $c$ dans le vide ; seules $\lambda_0$ et $\nu$, liées par $\lambda_0 = c/\nu$, les distinguent. Mais pourquoi les corps chauffés — un filament de lampe, une braise, le Soleil — émettent-ils de la lumière, et pourquoi leur couleur change-t-elle avec la température ? C'est l'objet de la section suivante, où nous découvrirons la \cle{loi de Wien}.

\coursec{Couleur des corps chauffés et loi de Wien}

Dans la section précédente, nous avons classé les ondes électromagnétiques selon leur longueur d'onde, de l'infrarouge au visible puis à l'ultraviolet. Une question naturelle se pose alors : \emph{qu'est-ce qui détermine le domaine dans lequel une source émet le plus de lumière ?} Pour les lampes spectrales, la réponse tient à la nature du gaz. Mais il existe une famille immense de sources — le Soleil, le filament d'une lampe à incandescence, la braise d'un feu de bois, le fer chauffé par le forgeron — dont le spectre continu ne dépend que d'\textbf{un seul paramètre : la température}. C'est ce que nous allons découvrir maintenant, et cela débouchera sur une loi simple et puissante : la \cle{loi de Wien}.

\courssub{Observation : la couleur d'un corps chauffé change avec la température}

\begin{experience}[Le fer du forgeron et le filament de la lampe]
\textbf{Situation 1 — Le forgeron.} Dans les quartiers de Yaoundé ou de Douala, le forgeron chauffe une barre de fer dans sa forge. Observons l'évolution de sa couleur :
\begin{itemize}
\item à la température ambiante, le fer est gris et n'émet aucune lumière visible ;
\item vers \SI{1000}{K} (environ \SI{730}{\celsius}), il devient \textbf{rouge sombre} ;
\item en chauffant davantage, il passe au \textbf{rouge vif}, puis à l'\textbf{orange} (vers \SI{1500}{K}) ;
\item très chaud, il devient \textbf{jaune} puis \textbf{blanc} (« fer porté au blanc », vers \SI{2000}{K} et plus).
\end{itemize}
\textbf{Situation 2 — La lampe à incandescence.} Le filament de tungstène d'une vieille ampoule, porté à environ \SI{2700}{K}, brille d'une lumière jaunâtre. Le Soleil, dont la surface est à environ \SI{5800}{K}, nous apparaît blanc-jaune.

\textbf{Interprétation.} La couleur de la lumière émise par un corps chauffé ne dépend pas de sa nature chimique (fer, tungstène, gaz solaire...) mais uniquement de sa \cle{température}. Plus le corps est chaud, plus la lumière émise est riche en courtes longueurs d'onde (bleu) ; moins il est chaud, plus elle est dominée par les grandes longueurs d'onde (rouge), voire par l'infrarouge invisible.
\end{experience}

\begin{popfigure}
\begin{center}
\begin{tikzpicture}
% Trois états du fer chauffé
\foreach \x/\col/\temp in {0/popmuted/{300 K (gris, aucune émission visible)}, 4.6/poporange!80!black/{1000 K (rouge sombre)}, 9.2/popgold/{1800 K (orange vif)}}{
\fill[\col] (\x,0) rectangle (\x+2.6,0.55);
\draw[popdark, thick] (\x,0) rectangle (\x+2.6,0.55);
\node[align=center, font=\small] at (\x+1.3,-0.55) {\temp};
}
% Flèche de chauffage
\draw[-{Stealth[length=3mm]}, very thick, popink] (-0.3,-1.4) -- (12.3,-1.4) node[midway, below, font=\small\itshape] {la température augmente};
% Flammes stylisées
\foreach \x in {0.6, 5.2, 9.8}{
\draw[poporange, thick, decorate, decoration={coil, aspect=0.5, segment length=4pt, amplitude=2pt}] (\x+0.7,-0.05) -- (\x+0.7,-1.0);
}
\node[font=\small\itshape, popmuted] at (6,-2.1) {rouge $\longrightarrow$ orange $\longrightarrow$ jaune $\longrightarrow$ blanc quand $T$ croît};
\end{tikzpicture}
\end{center}
\legende{Évolution de la couleur d'un morceau de fer lorsque sa température augmente. À température ambiante, le fer émet dans l'infrarouge lointain, invisible pour nos yeux.}
\end{popfigure}

\begin{aretenir}[Le point essentiel]
Tout corps porté à une température $T$ suffisamment élevée émet un rayonnement électromagnétique dont le spectre est \textbf{continu}. La répartition de la lumière entre les différentes longueurs d'onde — le \cle{profil spectral} — ne dépend pratiquement que de la température du corps, pas de sa composition. C'est ce qu'on appelle le \textbf{rayonnement thermique} (ou rayonnement du \emph{corps noir} dans le cas idéal).
\end{aretenir}

\courssub{Évolution du profil spectral avec la température}

Analysons le spectre continu d'un corps chauffé à l'aide d'un prisme et d'un capteur sensible à toutes les longueurs d'onde (y compris infrarouge et ultraviolet). On mesure, pour chaque longueur d'onde $\lambda$, l'intensité lumineuse émise, et on trace le profil spectral : intensité en fonction de $\lambda$.

\begin{popfigure}
\begin{center}
\begin{tikzpicture}
\begin{axis}[
width=0.95\linewidth, height=7.2cm,
xlabel={$\lambda$ (nm)}, ylabel={intensité émise (unité arbitraire)},
xmin=0, xmax=1600, ymin=0, ymax=1.15,
xtick={380, 580, 780, 1000, 1400},
legend style={font=\small, at={(0.97,0.97)}, anchor=north east},
axis line style={popdark},
every axis plot/.append style={very thick},
]
% Courbe 3000 K : max à 967 nm
\addplot[popink, domain=100:1600, samples=150] {1.9 * (x/967)^(-4) * exp(4*(1 - 967/x)) * (x>200)};
\addlegendentry{$T = \SI{3000}{K}$ \; ($\lambda_{max} \approx \SI{967}{nm}$)}
% Courbe 5000 K : max à 580 nm
\addplot[poporange, domain=100:1600, samples=150] {0.55 * (x/580)^(-4) * exp(4*(1 - 580/x))};
\addlegendentry{$T = \SI{5000}{K}$ \; ($\lambda_{max} \approx \SI{580}{nm}$)}
% Courbe 6500 K : max à 446 nm
\addplot[popblue, domain=100:1600, samples=150] {0.30 * (x/446)^(-4) * exp(4*(1 - 446/x))};
\addlegendentry{$T = \SI{6500}{K}$ \; ($\lambda_{max} \approx \SI{446}{nm}$)}
% Bande visible
\draw[popmuted, dashed] (axis cs:380,0) -- (axis cs:380,1.1);
\draw[popmuted, dashed] (axis cs:780,0) -- (axis cs:780,1.1);
\node[font=\footnotesize, popmuted, rotate=90, anchor=south] at (axis cs:400,0.02) {visible};
\node[font=\footnotesize, popmuted] at (axis cs:1200,0.06) {infrarouge};
\node[font=\footnotesize, popmuted] at (axis cs:200,0.06) {UV};
\end{axis}
\end{tikzpicture}
\end{center}
\legende{Profils spectraux (allure simplifiée) de corps chauffés à trois températures. Quand $T$ augmente, deux phénomènes se produisent : le maximum d'émission se déplace vers les \textbf{courtes} longueurs d'onde, et l'intensité totale émise augmente fortement.}
\end{popfigure}

Deux constatations s'imposent à la lecture de ces courbes :

\begin{enumerate}
\item Le profil spectral présente un \textbf{maximum d'émission} pour une longueur d'onde particulière, notée $\lambda_{max}$.
\item Quand la température $T$ augmente, ce maximum \textbf{se déplace vers les courtes longueurs d'onde} (vers le bleu, puis l'ultraviolet), tandis que l'énergie totale rayonnée augmente considérablement.
\end{enumerate}

C'est exactement ce qu'observait le forgeron : à \SI{1000}{K}, le maximum est dans l'infrarouge, mais la « queue » du spectre atteint le rouge — le fer paraît rouge. À \SI{5800}{K} comme le Soleil, le maximum tombe en plein dans le visible et toutes les couleurs sont présentes en quantités comparables — la lumière paraît blanche.

\courssub{La loi de Wien}

Le physicien allemand Wilhelm Wien a établi en 1893 une relation remarquablement simple entre la température d'un corps chauffé et la longueur d'onde de son maximum d'émission.

\begin{propriete}[Loi de Wien (admise)]
Pour un corps chauffé à la température absolue $T$ et émettant un spectre continu, la longueur d'onde $\lambda_{max}$ correspondant au maximum d'émission vérifie :
\[ \lambda_{max} \times T = \num{2,90e-3} \ \si{m.K} \]
avec :
\begin{itemize}
\item $\lambda_{max}$ en \textbf{mètres} (m) ;
\item $T$ en \textbf{kelvins} (K), température absolue : $T(\text{K}) = \theta(^\circ\text{C}) + 273$.
\end{itemize}
\emph{Au Probatoire, l'expression de la loi de Wien est donnée dans l'énoncé ; il faut savoir l'exploiter. Aucune démonstration n'est exigible.}
\end{propriete}

\textbf{Signification physique.} Le produit $\lambda_{max} \times T$ est constant : $\lambda_{max}$ et $T$ sont donc \textbf{inversement proportionnels}. Plus un corps est \cle{chaud}, plus sa longueur d'onde d'émission maximale est \cle{petite} (le maximum glisse vers le bleu et l'UV) ; plus il est froid, plus $\lambda_{max}$ est grande (le maximum glisse vers le rouge et l'infrarouge).

\textbf{Vérification par analyse dimensionnelle.} Le membre de gauche s'exprime en $\si{m} \times \si{K} = \si{m.K}$, ce qui correspond bien à l'unité de la constante de Wien $\num{2,90e-3}\ \si{m.K}$. La relation est homogène.

\begin{attention}[La température doit être en kelvins !]
L'erreur la plus fréquente au Probatoire est d'utiliser la température en degrés Celsius dans la loi de Wien. La loi fait intervenir la \textbf{température absolue}. Effectue toujours la conversion \emph{avant} tout calcul :
\[ T(\text{K}) = \theta(^\circ\text{C}) + 273 \]
Exemple : un filament à $\theta = \SI{2527}{\celsius}$ correspond à $T = 2527 + 273 = \SI{2800}{K}$, et non à \SI{2527}{K}. De même, convertis $\lambda_{max}$ en \textbf{mètres} : $\SI{500}{nm} = \num{5,00e-7}\ \si{m}$.
\end{attention}

\begin{attention}[Domaine de validité de la loi de Wien]
La loi de Wien ne s'applique \textbf{que} aux corps émettant un \textbf{spectre continu} d'origine thermique : Soleil et étoiles, filaments de lampes à incandescence, braises, lave, corps humain... Elle \textbf{ne s'applique pas} :
\begin{itemize}
\item aux \textbf{lampes spectrales} (vapeur de mercure, de sodium), dont le spectre est formé de \emph{raies} caractéristiques du gaz et non d'un profil continu ;
\item aux \textbf{lasers} et aux \textbf{DEL}, dont la lumière quasi monochromatique provient de transitions électroniques et non de l'agitation thermique.
\end{itemize}
Une DEL rouge n'est pas un corps « froid » et une DEL bleue n'est pas un corps « chaud » : leur couleur n'a rien à voir avec leur température !
\end{attention}

\courssub{Exploitation de la loi de Wien}

\begin{methode}[Exploiter la loi de Wien en trois étapes]
\begin{enumerate}
\item \textbf{Convertir} toutes les données dans les bonnes unités : la température en kelvins ($T = \theta + 273$) et la longueur d'onde en mètres ($1\ \text{nm} = 10^{-9}\ \text{m}$).
\item \textbf{Isoler} la grandeur cherchée dans la relation $\lambda_{max} \times T = \num{2,90e-3}\ \si{m.K}$ : soit $T = \dfrac{\num{2,90e-3}}{\lambda_{max}}$, soit $\lambda_{max} = \dfrac{\num{2,90e-3}}{T}$.
\item \textbf{Calculer} puis \textbf{vérifier l'ordre de grandeur} : une température de surface stellaire se compte en milliers de kelvins ; un corps à quelques centaines de kelvins émet son maximum dans l'infrarouge ($\lambda_{max} \approx \SI{10}{\micro m}$) ; seul un corps à plusieurs milliers de kelvins a son maximum dans le visible.
\end{enumerate}
\end{methode}

\begin{exemplebox}[Température de surface du Soleil]
Le spectre continu de la lumière solaire présente un maximum d'émission pour $\lambda_{max} \approx \SI{500}{nm}$ (lumière verte-bleue, en plein dans le visible). Déterminons la température de surface du Soleil.

\textbf{Étape 1 — Conversion :} $\lambda_{max} = \SI{500}{nm} = \num{5,00e-7}\ \si{m}$.

\textbf{Étape 2 — Isolation :} $T = \dfrac{\num{2,90e-3}\ \si{m.K}}{\lambda_{max}}$.

\textbf{Étape 3 — Calcul :}
\[ T = \frac{\num{2,90e-3}}{\num{5,00e-7}} = \num{5,80e3}\ \si{K} = \SI{5800}{K} \]
soit environ \SI{5500}{\celsius}. L'ordre de grandeur (quelques milliers de kelvins pour une étoile) est cohérent. C'est bien parce que le maximum solaire tombe dans le visible que nos yeux, adaptés au fil de l'évolution, y sont les plus sensibles !
\end{exemplebox}

\begin{exoresolu}[Dans quel domaine le corps humain émet-il le plus ?]
\textbf{Énoncé.} La température du corps humain est $\theta = \SI{37}{\celsius}$.
\begin{enumerate}
\item Convertir cette température en kelvins.
\item Calculer la longueur d'onde $\lambda_{max}$ du maximum d'émission du corps humain.
\item En déduire le domaine des ondes électromagnétiques dans lequel le corps humain émet le plus. Citer une application technologique.
\end{enumerate}
On donne : $\lambda_{max} \times T = \num{2,90e-3}\ \si{m.K}$.
\tcblower
\textbf{Solution détaillée.}
\begin{enumerate}
\item Conversion en kelvins :
\[ T = \theta + 273 = 37 + 273 = \SI{310}{K} \]
\item D'après la loi de Wien :
\[ \lambda_{max} = \frac{\num{2,90e-3}}{T} = \frac{\num{2,90e-3}}{310} \approx \num{9,4e-6}\ \si{m} = \SI{9,4}{\micro m} \]
\item Or $\SI{9,4}{\micro m} = \SI{9400}{nm}$, valeur très supérieure à \SI{780}{nm} : le corps humain émet principalement dans l'\textbf{infrarouge}. C'est pourquoi il est invisible dans le noir pour nos yeux, mais parfaitement « visible » pour une caméra thermique.

\textbf{Applications :} les caméras thermiques utilisées la nuit, les thermomètres frontaux infrarouges employés dans les hôpitaux du Cameroun (par exemple lors des contrôles sanitaires), et les détecteurs de présence exploitent tous ce rayonnement infrarouge du corps humain.
\end{enumerate}
\end{exoresolu}

\courssub{Couleur des étoiles et température : le complément Probatoire}

La loi de Wien fournit un pont direct entre deux grandeurs que l'on peut observer à l'œil nu dans le ciel : la \textbf{couleur} d'une étoile et sa \textbf{température de surface}.

\begin{propriete}[Couleur dominante et température d'une étoile]
La couleur dominante d'une étoile est fixée par la position de $\lambda_{max}$ dans (ou autour du) domaine visible :
\begin{itemize}
\item une étoile \textbf{rouge} a son maximum dans le rouge ou l'infrarouge : c'est une étoile relativement \textbf{froide} ($T \approx \SI{3000}{K}$ à $\SI{4000}{K}$) ;
\item une étoile \textbf{jaune} (comme le Soleil) a son maximum au milieu du visible : $T \approx \SI{5800}{K}$ ;
\item une étoile \textbf{bleue} a son maximum dans l'ultraviolet et émet beaucoup dans le bleu : c'est une étoile très \textbf{chaude} ($T \gtrsim \SI{10000}{K}$).
\end{itemize}
Contrairement à l'intuition quotidienne (le bleu de la « flamme froide »), \textbf{en astronomie, le bleu signifie chaud et le rouge signifie (relativement) froid}.
\end{propriete}

\begin{center}
\begin{tabularx}{\linewidth}{|l|c|c|X|}
\hline
\rowcolor{popblueL} \textbf{Étoile} & \textbf{Couleur} & $\boldsymbol{T}$ \textbf{(K)} & $\boldsymbol{\lambda_{max}}$ \textbf{calculée} \\
\hline
Bételgeuse (Orion) & rouge & $\approx 3500$ & $\dfrac{\num{2,90e-3}}{3500} \approx \SI{830}{nm}$ (proche IR) \\
\hline
Soleil & jaune-blanc & $\approx 5800$ & $\approx \SI{500}{nm}$ (visible) \\
\hline
Rigel (Orion) & bleue & $\approx 11000$ & $\dfrac{\num{2,90e-3}}{11000} \approx \SI{260}{nm}$ (UV) \\
\hline
\end{tabularx}
\end{center}

\begin{savaistu}[La loi de Wien, thermomètre des astronomes]
Comment mesurer la température d'une étoile située à des centaines d'années-lumière, sans jamais pouvoir s'en approcher ? Grâce à la loi de Wien ! Les astronomes dispersent la lumière de l'étoile avec un spectroscope, repèrent la longueur d'onde $\lambda_{max}$ du maximum du spectre continu, et en déduisent la température de surface. C'est ainsi que l'on sait que \textbf{Bételgeuse}, l'épaule rougeoyante de la constellation d'Orion visible à l'œil nu depuis le Cameroun, est une supergéante rouge à environ \SI{3500}{K}, tandis que \textbf{Rigel}, le pied bleuté de la même constellation, culmine à environ \SI{11000}{K}. La même méthode, appliquée au fond diffus cosmologique (le rayonnement fossile de l'Univers), donne un maximum d'émission vers $\lambda_{max} \approx \SI{1,06}{mm}$ dans le domaine des micro-ondes, soit une température de l'Univers d'environ \SI{2,73}{K} : l'Univers se refroidit depuis le Big Bang !
\end{savaistu}

\begin{exercice}[Pour t'entraîner]
\begin{enumerate}
\item Un filament de tungstène est porté à $\theta = \SI{2427}{\celsius}$. Calculer $\lambda_{max}$ et préciser le domaine d'émission maximal. La lumière de la lampe paraît-elle plutôt blanche ou jaunâtre ?
\item Une étoile présente un maximum d'émission à $\lambda_{max} = \SI{290}{nm}$. Calculer sa température de surface et prévoir sa couleur dominante.
\end{enumerate}
\end{exercice}

\begin{corrige}[Exercice]
\begin{enumerate}
\item $T = 2427 + 273 = \SI{2700}{K}$. Alors $\lambda_{max} = \dfrac{\num{2,90e-3}}{2700} \approx \num{1,07e-6}\ \si{m} = \SI{1070}{nm}$ : le maximum est dans l'\textbf{infrarouge}. Seule la partie du spectre située dans le visible, dominée par les grandes longueurs d'onde, atteint nos yeux : la lumière paraît \textbf{jaunâtre} (et la lampe gaspille beaucoup d'énergie en chaleur, d'où leur remplacement par les DEL).
\item $T = \dfrac{\num{2,90e-3}}{\num{2,90e-7}} = \num{1,00e4}\ \si{K} = \SI{10000}{K}$. Le maximum est dans l'ultraviolet, donc l'étoile émet beaucoup dans le bleu : c'est une étoile \textbf{bleue}, très chaude, comparable à Rigel.
\end{enumerate}
\end{corrige}

\begin{aretenir}[L'essentiel de la section]
\begin{itemize}
\item Tout corps chauffé émet un \textbf{spectre continu} dont le profil ne dépend que de sa \textbf{température}.
\item Quand $T$ augmente, le maximum d'émission se déplace vers les \textbf{courtes longueurs d'onde} : rouge $\rightarrow$ orange $\rightarrow$ jaune $\rightarrow$ blanc $\rightarrow$ bleu.
\item \textbf{Loi de Wien} : $\lambda_{max} \times T = \num{2,90e-3}\ \si{m.K}$, avec $\lambda_{max}$ en mètres et $T$ en \textbf{kelvins} ($T = \theta + 273$).
\item Elle permet de mesurer la température des étoiles et de prévoir le domaine d'émission d'un corps (le corps humain à \SI{310}{K} émet dans l'infrarouge).
\item Elle ne s'applique ni aux lampes spectrales (raies), ni aux lasers, ni aux DEL.
\end{itemize}
\end{aretenir}

\coursec{Démarche expérimentale : lumières colorées}

Dans la section précédente, nous avons vu que la couleur d'un corps chauffé dépend de sa température : la loi de Wien $\lambda_{\max} T = \num{2,90e-3}\ \mathrm{m\,K}$ relie la radiation dominante émise à la température de surface. Mais dans la vie quotidienne, la plupart des lumières colorées que nous rencontrons — feux tricolores, projecteurs de spectacle, enseignes lumineuses de Douala ou de Yaoundé — ne proviennent pas de corps chauffés à des températures différentes. Elles sont obtenues à partir d'une \emph{même} source blanche. Comment ? C'est toute la question de cette dernière section, que nous allons traiter en suivant pas à pas la \cle{démarche expérimentale}, compétence exigible au Probatoire.

\courssub{Le problème : obtenir une lumière colorée à partir d'une source blanche}

\begin{experience}[Situation-problème]
Sur la scène d'un spectacle au Palais des Sports de Yaoundé, un unique type de projecteur à lumière blanche permet d'obtenir, tour à tour, un éclairage rouge, vert ou bleu. Devant un feu de circulation, la même ampoule blanche éclaire successivement derrière un verre rouge, orange ou vert.

\textbf{Problème scientifique :} comment obtenir une lumière colorée à partir d'une source de lumière blanche ? Le filtre coloré « teinte-t-il » la lumière, ou bien se contente-t-il de \emph{sélectionner} certaines radiations déjà présentes dans la lumière blanche ?
\end{experience}

Deux hypothèses rivales s'affrontent naturellement :
\begin{itemize}
\item \textbf{Hypothèse 1 :} le filtre coloré « transforme » la lumière blanche en lumière colorée, comme un colorant teinterait de l'eau ;
\item \textbf{Hypothèse 2 :} le filtre coloré \cle{sélectionne} certaines radiations déjà présentes dans la lumière blanche : il transmet certaines longueurs d'onde et absorbe les autres.
\end{itemize}

Pour trancher entre ces deux hypothèses, il faut une expérience capable de \emph{décomposer} la lumière avant et après le filtre. L'outil idéal, nous l'avons appris dans cette leçon, c'est le \cle{prisme}.

\courssub{L'expérience : filtre coloré, prisme et spectre}

\begin{experience}[Analyse au prisme de la lumière transmise par un filtre]
\textbf{Matériel :} lampe à incandescence (source blanche), fente fine, prisme en verre, écran blanc, filtres colorés (rouge puis vert).

\textbf{Protocole :}
\begin{enumerate}
\item Éclairer la fente avec la lampe blanche et disperser la lumière avec le prisme : observer le spectre sur l'écran.
\item Interposer un filtre rouge entre la lampe et la fente. Observer le nouveau spectre.
\item Recommencer avec un filtre vert. Comparer les trois spectres.
\end{enumerate}

\textbf{Observations :}
\begin{itemize}
\item Sans filtre : spectre continu, du violet au rouge (toutes les radiations visibles).
\item Avec le filtre rouge : seule une \emph{bande rouge-orangée} subsiste sur l'écran ; les radiations bleues, vertes et violettes ont disparu.
\item Avec le filtre vert : seule une \emph{bande verte} (parfois débordant vers le jaune et le cyan) subsiste.
\end{itemize}

\textbf{Interprétation :} les radiations vertes existaient déjà dans la lumière blanche avant le filtre rouge — le prisme les révélait. Après le filtre rouge, elles ont disparu du spectre : le filtre les a donc \emph{absorbées}, et non « transformées ». L'hypothèse 2 est validée, l'hypothèse 1 est réfutée.
\end{experience}

\begin{popfigure}
\begin{center}
\begin{tikzpicture}[scale=1.0, font=\small]
% Lampe
\draw[fill=popgoldL] (0,0.35) circle (0.35);
\draw[thick] (-0.15,-0.05) -- (-0.15,-0.35) -- (0.15,-0.35) -- (0.15,-0.05);
\draw[thick] (-0.12,-0.35) -- (0.12,-0.35);
\node[below] at (0,-0.42) {lampe blanche};
% Rayons blancs
\draw[thick,popmuted,->] (0.4,0.2) -- (1.5,0.2);
\draw[thick,popmuted,->] (0.4,0.05) -- (1.5,0.05);
\draw[thick,popmuted,->] (0.4,-0.1) -- (1.5,-0.1);
\node[above,popmuted] at (0.95,0.22) {\scriptsize lumière blanche};
% Filtre rouge
\draw[fill=red!70,draw=red!80!black,thick] (1.6,-0.6) rectangle (1.85,0.7);
\node[below,red!80!black] at (1.72,-0.65) {\scriptsize filtre rouge};
% Rayons rouges
\draw[thick,red,->] (1.9,0.15) -- (2.9,0.15);
\draw[thick,red,->] (1.9,0.0) -- (2.9,0.0);
\node[above,red] at (2.4,0.17) {\scriptsize bande rouge};
% Prisme
\draw[fill=popblueL,draw=popblue,thick] (3.0,-0.55) -- (3.7,0.65) -- (4.4,-0.55) -- cycle;
\node[below,popblue] at (3.7,-0.62) {\scriptsize prisme};
% Spectre partiel sur écran
\draw[thick] (6.1,-1.0) -- (6.1,1.3);
\fill[red!80] (6.1,-0.35) rectangle (6.35,0.35);
\node[right] at (6.4,0) {\scriptsize bande rouge};
\node[right,popmuted] at (6.4,0.75) {\scriptsize (rien : bleu, vert absorbés)};
\node[below] at (6.2,-1.05) {\scriptsize écran};
\end{tikzpicture}
\end{center}
\legende{Dispositif expérimental : la lumière blanche traverse un filtre rouge ; le prisme révèle que seule la bande rouge du spectre est transmise.}
\end{popfigure}

\begin{propriete}[Action d'un filtre coloré]
Un \cle{filtre coloré} transmet certaines radiations de la lumière qui le traverse et absorbe les autres. Il ne « colore » pas la lumière : il \textbf{sélectionne} des radiations déjà présentes dans la lumière incidente. La lumière transmise est \cle{polychromatique} (bande de longueurs d'onde plus ou moins large) ou, pour les filtres très sélectifs, \emph{quasi-monochromatique}.
\end{propriete}

\begin{attention}[Un filtre rouge ne donne pas une lumière monochromatique]
Erreur classique au Probatoire : écrire que « la lumière traversant un filtre rouge est monochromatique de longueur d'onde $\lambda = \SI{650}{nm}$ ». Faux ! Le spectre de la lumière transmise est une \textbf{bande continue} s'étendant typiquement de $\SI{600}{nm}$ à $\SI{700}{nm}$ : c'est une lumière \emph{polychromatique} dont le profil spectral est centré sur le rouge. Seul un laser (ou une raie d'une lampe spectrale isolée par une fente) fournit une lumière rigoureusement monochromatique. Retiens : \textbf{filtre = bande passante plus ou moins large, jamais une unique raie}.
\end{attention}

\courssub{Comparaison des spectres transmis : filtre rouge et filtre vert}

Analysons quantitativement. Un filtre est caractérisé par sa \cle{courbe de transmission} $T(\lambda)$, fraction de l'intensité lumineuse transmise à chaque longueur d'onde. Le profil spectral de la lumière transmise s'obtient en multipliant le profil de la source par la courbe du filtre :
\[ I_{\text{transmise}}(\lambda) = T(\lambda) \times I_{\text{source}}(\lambda). \]

\begin{exemplebox}[Bandes passantes de deux filtres]
Un filtre rouge de gélatine transmet plus de $80\,\%$ de l'intensité pour $\lambda \in [\SI{610}{nm}\,;\,\SI{700}{nm}]$ et moins de $5\,\%$ pour $\lambda < \SI{580}{nm}$. Un filtre vert transmet surtout la bande $[\SI{490}{nm}\,;\,\SI{560}{nm}]$. Placés devant la même lampe blanche (spectre continu), ils donnent deux lumières transmises de couleurs différentes, toutes deux polychromatiques. Si l'on superpose les deux filtres, aucune bande ne leur est commune : presque aucune lumière ne passe — l'ensemble paraît noir.
\end{exemplebox}

Ce dernier résultat est un test décisif supplémentaire : si le filtre « teintait » la lumière, deux filtres superposés donneraient une lumière mélange des deux teintes ; or on obtient l'obscurité, ce qui confirme le mécanisme de \emph{sélection-absorption}.

\courssub{La synthèse additive : reconstruire la lumière blanche}

Nous venons de \emph{décomposer} la lumière blanche. Peut-on faire l'opération inverse, c'est-à-dire \emph{composer} de la lumière blanche à partir de lumières colorées ? Oui : c'est la \cle{synthèse additive}.

\begin{experience}[Superposition de trois faisceaux colorés]
Trois projecteurs munis respectivement d'un filtre rouge, d'un filtre vert et d'un filtre bleu éclairent un même écran blanc, en se recouvrant partiellement. On observe :
\begin{itemize}
\item rouge $+$ vert $=$ \textbf{jaune} ;
\item vert $+$ bleu $=$ \textbf{cyan} ;
\item bleu $+$ rouge $=$ \textbf{magenta} ;
\item rouge $+$ vert $+$ bleu $=$ \textbf{blanc} (zone commune aux trois faisceaux).
\end{itemize}
L'œil, stimulé simultanément par les trois bandes de radiations, perçoit la même sensation que la lumière blanche.
\end{experience}

\begin{popfigure}
\begin{center}
\begin{tikzpicture}[scale=0.95]
\begin{scope}[blend group=screen]
\fill[red] (90:1.1) circle (1.6);
\fill[green] (210:1.1) circle (1.6);
\fill[blue] (330:1.1) circle (1.6);
\end{scope}
\node[white] at (90:2.35) {\textbf{Rouge}};
\node[white] at (210:2.35) {\textbf{Vert}};
\node[white] at (330:2.35) {\textbf{Bleu}};
\node[black] at (0,0) {\textbf{Blanc}};
\node[black] at (0,1.55) {\scriptsize Jaune};
\node[black,rotate=60] at (-1.35,-0.8) {\scriptsize Magenta};
\node[black,rotate=-60] at (1.35,-0.8) {\scriptsize Cyan};
\end{tikzpicture}
\end{center}
\legende{Synthèse additive : la superposition des trois faisceaux rouge, vert et bleu restitue du blanc au centre ; les recouvrements deux à deux donnent le jaune, le cyan et le magenta.}
\end{popfigure}

\begin{savaistu}[Les écrans de ton téléphone : la synthèse additive à l'œuvre]
L'écran de ton smartphone est constitué de millions de \cle{pixels}, chacun divisé en trois sous-pixels : un rouge, un vert, un bleu (technologie \textbf{RVB}, ou RGB en anglais). En réglant l'intensité lumineuse de chaque sous-pixel de 0 à 255, l'écran reproduit $256^3 \approx \num{16,7}$ millions de couleurs ! Pour afficher du jaune, l'écran allume les sous-pixels rouge et vert et éteint le bleu. Observe l'écran avec une goutte d'eau faisant loupe (ou une simple loupe) : tu verras les trois petites bandes lumineuses. Même principe sur les écrans géants des stades lors des matchs des Lions Indomptables.
\end{savaistu}

\courssub{Rédiger le compte rendu d'une démarche expérimentale}

Au Probatoire comme en TP, tu dois savoir rédiger un compte rendu structuré. Voici la méthode complète, appliquée à notre étude.

\begin{methode}[Conduire et rédiger une démarche expérimentale]
\begin{enumerate}
\item \textbf{Problème :} formuler la question scientifique (ici : comment obtenir une lumière colorée à partir d'une source blanche ?).
\item \textbf{Hypothèse :} proposer une réponse vérifiable (le filtre sélectionne des radiations).
\item \textbf{Protocole :} décrire le dispositif et les étapes (lampe blanche, filtre, prisme, écran) en précisant ce qu'on mesure ou observe.
\item \textbf{Résultats :} rapporter les observations \emph{brutes}, sans interprétation (spectre complet sans filtre ; bande rouge avec le filtre rouge ; bande verte avec le filtre vert).
\item \textbf{Conclusion :} confronter les résultats à l'hypothèse, valider ou réfuter, puis énoncer la loi générale (un filtre transmet certaines radiations et absorbe les autres).
\end{enumerate}
\end{methode}

\begin{exoresolu}[Exploiter un profil spectral après filtrage]
\textbf{Énoncé.} Une lampe blanche émet un spectre continu d'intensité à peu près uniforme entre $\SI{400}{nm}$ et $\SI{700}{nm}$. On place devant elle un filtre dont la transmission vaut $T = 0{,}90$ pour $\lambda \in [\SI{490}{nm}\,;\,\SI{560}{nm}]$ et $T \approx 0$ ailleurs.

1. Quelle est la couleur de la lumière transmise ?
2. Cette lumière est-elle monochromatique ? Justifier.
3. On analyse la lumière transmise avec un prisme. Décrire ce qu'on observe sur l'écran.
4. On superpose ce filtre à un filtre rouge transmettant la bande $[\SI{610}{nm}\,;\,\SI{700}{nm}]$. Qu'observe-t-on ?

\tcblower
\textbf{Solution détaillée.}

1. La bande transmise $[\SI{490}{nm}\,;\,\SI{560}{nm}]$ correspond au domaine du \textbf{vert} (avec un peu de cyan). La lumière transmise est donc \textbf{verte}.

2. \textbf{Non}, elle n'est pas monochromatique : elle contient une infinité de radiations de longueurs d'onde comprises entre $\SI{490}{nm}$ et $\SI{560}{nm}$, soit une bande de largeur $\Delta\lambda = 560 - 490 = \SI{70}{nm}$. C'est une lumière \textbf{polychromatique} à bande étroite (on dit parfois « quasi-monochromatique » si la bande est très fine, ce qui n'est pas le cas ici).

3. Le prisme disperse les radiations transmises : on observe sur l'écran une \textbf{bande continue verte} (dégradé du cyan au vert), et rien dans les régions où apparaissaient le violet, le bleu, le jaune, l'orange et le rouge avec la lampe seule. Le spectre complet est devenu un spectre \emph{partiel}.

4. Les bandes passantes des deux filtres sont disjointes : $[\SI{490}{nm}\,;\,\SI{560}{nm}] \cap [\SI{610}{nm}\,;\,\SI{700}{nm}] = \varnothing$. Le filtre rouge absorbe la bande verte transmise par le premier filtre : \textbf{aucune lumière n'est transmise}, l'écran reste noir. Cela confirme que le filtre sélectionne des radiations au lieu de « teinter » la lumière.
\end{exoresolu}

\courssub{Applications : éclairage scénique, signalisation, écrans}

\begin{aretenir}[L'essentiel de la démarche]
\begin{itemize}
\item Un filtre coloré \textbf{transmet} certaines radiations et \textbf{absorbe} les autres : il sélectionne, il ne colore pas.
\item La lumière transmise par un filtre est \textbf{polychromatique} (bande passante), jamais rigoureusement monochromatique.
\item Le prisme est l'instrument qui permet d'\emph{analyser} la lumière transmise et de valider l'hypothèse.
\item La \textbf{synthèse additive} des trois couleurs rouge, vert et bleu restitue la lumière blanche ; c'est le principe des écrans et de l'éclairage scénique.
\end{itemize}
\end{aretenir}

Ces principes expliquent de nombreuses technologies familières :
\begin{itemize}
\item \textbf{Feux de signalisation :} une même source blanche derrière des verres filtrants rouge, orange et vert ; les longueurs d'onde choisies garantissent une visibilité maximale, de jour comme sous la pluie.
\item \textbf{Éclairage scénique :} les projecteurs à filtres (ou à DEL rouges, vertes, bleues) créent toutes les ambiances colorées des concerts et des cérémonies.
\item \textbf{Écrans (téléphones, téléviseurs, panneaux publicitaires) :} mosaïque de sous-pixels RVB exploitant la synthèse additive.
\item \textbf{Photographie et vidéoprojection :} les capteurs et les vidéoprojecteurs codent et recomposent les images en trois canaux rouge, vert, bleu.
\end{itemize}

\begin{exercice}[Pour t'entraîner]
On interpose un filtre magenta (qui transmet le rouge et le bleu, absorbe le vert) entre une lampe blanche et un prisme.

1. Décrire le spectre observé sur l'écran.
2. La lumière transmise est-elle monochromatique ?
3. Quelle couleur obtient-on en superposant sur un écran blanc un faisceau cyan (vert $+$ bleu) et un faisceau rouge ? Justifier par la synthèse additive.
\end{exercice}

\begin{corrige}[Exercice : filtre magenta et synthèse additive]
1. Le spectre complet de la lumière blanche est amputé de sa partie verte : on observe une bande bleue-violette et une bande rouge-orangée, séparées par une zone obscure à l'emplacement du vert.

2. Non : elle contient deux bandes de radiations (bleu et rouge) ; c'est une lumière polychromatique.

3. Cyan $=$ vert $+$ bleu ; en ajoutant le rouge, on obtient rouge $+$ vert $+$ bleu, c'est-à-dire du \textbf{blanc} par synthèse additive des trois couleurs primaires.
\end{corrige}

\newpage

\coursec{Activité d'intégration}

\coursec{Spectres lumineux et sources de lumière}

\courssub{Sources de lumière et obtention d'un spectre}

\begin{definition}[Source de lumière]
Une \cle{source de lumière} est un système qui émet de l'énergie lumineuse. On distingue deux catégories fondamentales :
\begin{itemize}
\item Les \cle{sources primaires} (ou sources propres) : elles émettent leur propre lumière par conversion d'une autre forme d'énergie (thermique, électrique, chimique, nucléaire). Exemples : le Soleil (fusion nucléaire), une lampe à incandescence (effet Joule), une DEL (électroluminescence), un laser (émission stimulée), une flamme (réaction chimique).
\item Les \cle{sources secondaires} (ou objets éclairés) : elles ne font que diffuser ou réfléchir la lumière reçue d'une source primaire. Exemples : la Lune, un écran de cinéma, la route la nuit, un manuel scolaire.
\end{itemize}
\end{definition}

\begin{experience}[Décomposition de la lumière blanche par un prisme — TP classique]
\textbf{Matériel :} Source de lumière blanche (lampe à incandescence ou projecteur à fente), prisme en verre optique, écran blanc, pièce obscurcie.\\
\textbf{Protocole :}
\begin{enumerate}
\item Placer la source face à une fente fine pour obtenir un faisceau étroit.
\item Orienter le prisme de manière à dévier le faisceau (angle d'incidence $\neq 0$).
\item Observer la tache lumineuse sur l'écran placé derrière le prisme.
\item Faire varier l'orientation du prisme pour observer le passage par le \cle{minimum de déviation}.
\end{enumerate}
\textbf{Observations :} Sur l'écran apparaît une bande colorée continue : le \cle{spectre} de la lumière blanche. Les couleurs se succèdent dans l'ordre : rouge, orange, jaune, vert, bleu, indigo, violet (mnemonic : \og ROYGBIV \fg{} ou \og ROUGE VIOLET \fg{}). La déviation est plus forte pour le violet (courtes longueurs d'onde) que pour le rouge (longues longueurs d'onde).
\textbf{Interprétation :} L'indice de réfraction $n$ du verre dépend de la longueur d'onde $\lambda$ (\cle{dispersion}) : $n_{\text{violet}} > n_{\text{rouge}}$. La lumière blanche est un \cle{mélange} de toutes les radiations visibles. Le prisme réalise une \cle{dispersion} angulaire qui sépare les composantes spectrales.
\end{experience}

\begin{propriete}[Spectre d'une source]
Le \cle{spectre} d'une source est la répartition de l'énergie lumineuse (ou de l'éclat, ou de l'intensité spectrale) en fonction de la longueur d'onde $\lambda$ (ou de la fréquence $\nu$). On l'obtient expérimentalement à l'aide d'un \cle{spectroscope} (à prisme ou à réseau de diffraction) muni d'un détecteur (œil, capteur CCD, photodiode).
\end{propriete}

\begin{savaistu}[Contexte camerounais : Électrification et éclairage]
Au Cameroun, le réseau \textsc{ENEO}/\textsc{SONATREL} alimente l'éclairage public et domestique. Le passage progressif des lampes à incandescence (filament de tungstène à $\approx 3000$ K) aux \cle{DEL} (diodes électroluminescentes) et aux lampes fluocompactes répond à un impératif d'efficacité énergétique : une DEL consomme 6 à 10 fois moins d'électricité pour le même flux lumineux, soulageant le réseau (barrages de \textsc{Lom Pangar}, \textsc{Nachtigal}, \textsc{Edéa}) et la facture des ménages. Les lampadaires solaires autonomes (panneau photovoltaïque + batterie + DEL) se déploient dans l'électrification rurale (\textsc{ADER}).
\end{savaistu}

\courssub{Sources monochromatiques et polychromatiques}

\begin{definition}[Source monochromatique]
Une source est dite \cle{monochromatique} si son spectre ne présente qu'une seule raie (ou une bande extrêmement étroite) : elle émet une radiation de longueur d'onde $\lambda_0$ unique (dans le vide). 
\begin{itemize}
\item \textbf{Laser} (Light Amplification by Stimulated Emission of Radiation) : source quasi-monochromatique, cohérente, unidirectionnelle. Exemple : laser He-Ne ($\lambda = 632,8$ nm, rouge), laser à semi-conducteur (pointeur laser vert $\lambda \approx 532$ nm, lecteur Blu-ray $\lambda = 405$ nm).
\item \textbf{Lampe à vapeur de sodium} (basse pression) : doublet jaune très fin $\lambda_1 = 589,0$ nm, $\lambda_2 = 589,6$ nm (quasi-monochromatique).
\end{itemize}
\end{definition}

\begin{definition}[Source polychromatique]
Une source est \cle{polychromatique} si son spectre s'étend sur une plage continue ou comporte de nombreuses raies. 
\begin{itemize}
\item \textbf{Spectre continu} : corps chauffés (Soleil, filament tungstène, forge). L'énergie est répartie sur tout le visible (et au-delà).
\item \textbf{Spectre de raies} : lampes à décharge (néon, mercure, tubes fluorescents). Chaque élément chimique émet un ensemble de raies caractéristiques (\og empreinte digitale \fg{}).
\item \textbf{Spectre de bandes} : molécules (gaz d'ambiance, flammes).
\end{itemize}
\end{definition}

\begin{aretenir}[Critère de distinction]
\begin{center}
\begin{tabularx}{\linewidth}{|l|X|X|}
\hline
\rowcolor{popblueL} \textbf{Caractère} & \textbf{Monochromatique} & \textbf{Polychromatique} \\
\hline
Profil spectral & Une raie unique (largeur $\Delta\lambda \ll \lambda$) & Bande continue ou multiples raies \\
Longueur d'onde & Caractérisée par $\lambda_0$ (dans le vide) & Intervalle $[\lambda_{\min}, \lambda_{\max}]$ \\
Exemples & Laser, Na basse pression, LED (bande étroite $\sim 20$ nm) & Soleil, lampe à incandescence, néon, flamme \\
\hline
\end{tabularx}
\end{center}
\end{aretenir}

\courssub{Analyse des profils spectraux}

\begin{definition}[Profil spectral]
Le \cle{profil spectral} $I(\lambda)$ (ou $E(\lambda)$) représente l'éclat spectral (ou l'irradiance spectrale) en fonction de la longueur d'onde. C'est la \og carte d'identité \fg{} de la source.
\end{definition}

\begin{propriete}[Grandeurs caractéristiques d'un profil spectral]
\begin{itemize}
\item \textbf{Longueur d'onde du maximum} $\lambda_{\max}$ : abscisse du pic principal. Détermine la couleur dominante pour une source thermique.
\item \textbf{Largeur spectrale} $\Delta\lambda$ (ou $\Delta\nu$) : largeur à mi-hauteur (FWHM -- \emph{Full Width at Half Maximum}). Quantifie la \og pureté \fg{} spectrale.
\item \textbf{Domaine spectral} : UV ($\lambda < 400$ nm), Visible ($400 \le \lambda \le 800$ nm), IR ($\lambda > 800$ nm).
\end{itemize}
\end{propriete}

\begin{exemplebox}[Profils types]
\begin{center}
\begin{tikzpicture}[scale=0.8, domain=300:1000, samples=200]
% Axes
\draw[->] (0,0) -- (7.5,0) node[right] {$\lambda$ (nm)};
\draw[->] (0,0) -- (0,3.5) node[above] {$I(\lambda)$ (u.a.)};
\foreach \x/\label in {3/400, 4/500, 5/600, 6/700, 7/800} \draw (\x,0.05) -- (\x,-0.05) node[below] {\label};
% Visible region shading
\fill[popblue!10] (1.33,0) rectangle (5.33,3.5); % 400-800 nm approx mapping
% Source thermique (continu, pic vers 500nm)
\draw[popblue, thick] plot ({(\x-100)/100}, {2.5*exp(-((\x-500)/150)^2)});
\node[popblue, right] at (7, 2.5) {Corps noir (Soleil)};
% Source monochromatique (raie fine)
\draw[poporange, thick] plot ({(\x-100)/100}, {3*exp(-((\x-650)/5)^2)});
\node[poporange, right] at (7, 1.8) {Laser / Na};
% Source à raies (néon)
\foreach \pos/\h in {3.5/1.5, 4.2/2.2, 5.0/1.0, 5.8/2.5, 6.5/1.2} {
    \draw[popgreen, thick] (\pos, 0) -- (\pos, \h);
}
\node[popgreen, right] at (7, 1.0) {Néon (raies)};
\end{tikzpicture}
\end{center}
\legende{Profils spectraux schématiques : continu (thermique), raie fine (monochromatique), raies discrètes (décharge gazeuse).}
\end{exemplebox}

\courssub{Les ondes électromagnétiques}

\begin{propriete}[Nature de la lumière]
La lumière est une \cle{onde électromagnétique} : propagation d'une perturbation couplée du champ électrique $\vec{E}$ et du champ magnétique $\vec{B}$, transversale, se propageant dans le vide à la vitesse $c = 3,00 \times 10^8\ \text{m}\cdot\text{s}^{-1}$. En milieu transparent d'indice $n$, la vitesse est $v = c/n$ et la longueur d'onde $\lambda = \lambda_0/n$ ($\lambda_0$ longueur d'onde dans le vide). La fréquence $\nu$ est invariante.
\[ c = \lambda_0 \nu \qquad ; \qquad E = h\nu = \frac{hc}{\lambda_0} \]
où $h = 6,63 \times 10^{-34}\ \text{J}\cdot\text{s}$ (constante de Planck).
\end{propriete}

\begin{aretenir}[Spectre électromagnétique — Ordres de grandeur]
\begin{center}
\begin{tabularx}{\linewidth}{|l|c|c|X|}
\hline
\rowcolor{popblueL} \textbf{Domaine} & \textbf{$\lambda_0$ (m)} & \textbf{$\nu$ (Hz)} & \textbf{Applications / Contexte camerounais} \\
\hline
Ondes radio & $> 1$ & $< 3 \times 10^8$ & Radio FM (90--108 MHz), TV, \textsc{MTN}/\textsc{Orange}/\textsc{CAMTEL} 2G/3G/4G/5G (700 MHz -- 3,5 GHz) \\
Micro-ondes & $10^{-3} \text{ à } 1$ & $3 \times 10^8 \text{ à } 3 \times 10^{11}$ & Four micro-ondes, liaisons hertziennes (backbone \textsc{CAMTEL}), radar météo (aéroport Nsimalen) \\
Infrarouge (IR) & $7,5 \times 10^{-7} \text{ à } 10^{-3}$ & $3 \times 10^{11} \text{ à } 4 \times 10^{14}$ & Télécommande, vision nocturne, séchage agricole (solaire), télécoms par fibre optique (1310, 1550 nm) \\
\rowcolor{popyellow!30} \textbf{Visible} & \textbf{$4 \times 10^{-7} \text{ à } 8 \times 10^{-7}$} & \textbf{$4 \times 10^{14} \text{ à } 7,5 \times 10^{14}$} & \textbf{Vision humaine, photosynthèse, éclairage (DEL, lampes)} \\
Ultraviolet (UV) & $10^{-8} \text{ à } 4 \times 10^{-7}$ & $7,5 \times 10^{14} \text{ à } 3 \times 10^{16}$ & Stérilisation eau/hôpital, détection billets (UV), bronzage, couche d'ozone \\
Rayons X & $10^{-11} \text{ à } 10^{-8}$ & $3 \times 10^{16} \text{ à } 3 \times 10^{19}$ & Radiographie (hôpitaux Yaoundé/Douala), scanner, contrôle soudures (pipelines \textsc{SNH}) \\
Rayons $\gamma$ & $< 10^{-11}$ & $> 3 \times 10^{19}$ & Radiothérapie (cancer, hôpital général), stérilisation matériel médical, radioactivité naturelle \\
\hline
\end{tabularx}
\end{center}
\end{aretenir}

\begin{attention}[Piège : Longueur d'onde dans le vide vs dans le milieu]
Toutes les longueurs d'onde caractéristiques ($\lambda_{\max}$, raies spectrales, domaine visible) sont \textbf{toujours données dans le vide} (ou l'air, $n_{\text{air}} \approx 1,0003 \approx 1$) sauf mention explicite contraire. Dans un milieu d'indice $n$, $\lambda = \lambda_0/n$ mais la couleur perçue ne change pas car l'œil analyse la fréquence $\nu$.
\end{attention}

\courssub{Couleur des corps chauffés et loi de Wien}

\begin{experience}[Expérience de la forge / filament de lampe — Observation qualitative]
\textbf{Matériel :} Alimentation variable (variac), lampe à incandescence à filament tungstène (culot E27) démunie de son bulbe (ou lampe témoin sur variateur), thermocouple type K (optionnel), pièce sombre.\\
\textbf{Protocole :} Augmenter progressivement la tension aux bornes de la lampe de 0 V à la tension nominale (230 V). Observer la couleur du filament.\\
\textbf{Observations :}
\begin{itemize}
\item $T < 800$ K : filament noir (rayonnement infrarouge invisible).
\item $T \approx 1000$ K : rouge sombre.
\item $T \approx 1500$ K : rouge vif.
\item $T \approx 2000$ K : orange.
\item $T \approx 2500$ K : jaune.
\item $T \approx 3000$ K : blanc chaud (lampe nominale).
\item $T > 5000$ K : blanc bleuté (inatteignable pour le tungstène qui fond à 3695 K).
\end{itemize}
\textbf{Interprétation :} Conformément à la loi de Wien, l'augmentation de $T$ décale $\lambda_{\max}$ vers les courtes longueurs d'onde. Le spectre s'élargit et envahit tout le visible $\rightarrow$ lumière blanche.
\end{experience}

\begin{propriete}[Loi de Wien (loi du déplacement de Wien)]
Pour un \cle{corps noir} (modèle idéal de source thermique), la longueur d'onde $\lambda_{\max}$ correspondant au maximum d'émission spectrale (en énergie par unité de longueur d'onde) est inversement proportionnelle à la température thermodynamique $T$ :
\[ \boxed{\lambda_{\max} \cdot T = b = 2,90 \times 10^{-3}\ \text{m}\cdot\text{K}} \]
où $b$ est la \cle{constante de Wien}.
\begin{itemize}
\item \textbf{Condition d'application :} Source thermique (corps chauffé à l'équilibre thermodynamique). Ne s'applique \textbf{pas} aux lasers, DEL, lampes à décharge (sources non thermiques / luminescence).
\item \textbf{Analyse dimensionnelle :} $[\lambda_{\max}] = \text{m}$, $[T] = \text{K}$ $\Rightarrow$ $[b] = \text{m}\cdot\text{K}$. Cohérent.
\item \textbf{Conséquence :} $T \uparrow \ \Rightarrow\ \lambda_{\max} \downarrow$ (décalage vers le bleu / UV). $T \downarrow \ \Rightarrow\ \lambda_{\max} \uparrow$ (décalage vers le rouge / IR).
\end{itemize}
\end{propriete}

\begin{methode}[Utilisation de la loi de Wien]
\begin{enumerate}
\item Identifier que la source est \textbf{thermique} (corps chauffé, étoile, filament).
\item Relever $\lambda_{\max}$ sur le profil spectral (donnée ou lecture graphique). \textbf{Convertir en mètres} ($1\ \text{nm} = 10^{-9}\ \text{m}$).
\item Appliquer $T = \dfrac{b}{\lambda_{\max}}$ avec $b = 2,90 \times 10^{-3}\ \text{m}\cdot\text{K}$.
\item Exprimer $T$ en kelvins (K) avec 2 ou 3 chiffres significatifs.
\item Situer $\lambda_{\max}$ dans le domaine spectral (UV, Visible, IR) pour interpréter la couleur apparente.
\end{enumerate}
\end{methode}

\begin{savaistu}[Température de couleur et éclairage]
La \cle{température de couleur} (en K) d'une source non thermique (DEL, fluocompacte) est la température du corps noir dont la couleur *apparaît* la plus proche. 
\begin{itemize}
\item $\approx 2700$ K : \og Blanc chaud \fg{} (ambiance salon, restaurant) — imite l'incandescence.
\item $\approx 4000$ K : \og Blanc neutre \fg{} (bureau, école).
\item $\approx 6500$ K : \og Blanc froid / Jour \fg{} (hôpital, atelier, extérieur).
\end{itemize}
Au Cameroun, le choix de la température de couleur influence le confort visuel et la consommation (efficacité lumineuse $\eta$ en lm/W dépend de la technologie, pas seulement de $T$).
\end{savaistu}

\courssub{Démarche expérimentale : lumières colorées}

\textbf{Objectifs de l'activité}

À l'issue de cette activité d'intégration, tu seras capable de :
\begin{itemize}
\item lire et exploiter le \cle{profil spectral} d'une source de lumière pour en extraire la longueur d'onde $\lambda_{\max}$ du maximum d'émission ;
\item appliquer la \cle{loi de Wien} pour déterminer la température de surface d'une source thermique (étoile, lampe à incandescence) ;
\item situer un rayonnement dans les \cle{domaines des ondes électromagnétiques} (ultraviolet, visible, infrarouge) ;
\item classer des sources par température et établir la relation entre la \cle{couleur} d'un corps chauffé et sa \cle{température} ;
\item travailler en groupe, rédiger une conclusion scientifique et la présenter oralement.
\end{itemize}

\textbf{Situation}

\textbf{Météo des étoiles : à la recherche des températures stellaires.}

Lors d'une nuit claire sur le campus de l'Université de Yaoundé I, le club d'astronomie observe le ciel à l'œil nu et avec de petites lunettes. Les élèves remarquent que les étoiles n'ont pas toutes la même couleur : certaines brillent d'une lumière bleutée, d'autres tirent sur le rouge, et le Soleil, lui, nous apparaît jaune-blanc en plein jour. Intrigués, ils se demandent : \emph{la couleur d'une étoile est-elle liée à sa température de surface ? Peut-on faire la \og météo \fg{} des étoiles sans jamais y aller ?}

Pour répondre, l'enseignant leur fournit les \textbf{profils spectraux simplifiés} (courbes de l'éclat lumineux en fonction de la longueur d'onde) de quatre sources, obtenus à l'aide d'un spectroscope à prisme :

\begin{center}
\begin{tabularx}{\linewidth}{|l|X|c|}
\hline
\rowcolor{popblueL} \textbf{Source} & \textbf{Description} & $\lambda_{\max}$ (nm) \\
\hline
Soleil & étoile jaune, notre référence & $500$ \\
\hline
Étoile A (type Bételgeuse) & étoile rougeoyante & $830$ \\
\hline
Étoile B (type Rigel) & étoile bleutée & $290$ \\
\hline
Lampe à incandescence & filament de tungstène (ampoule classique) & $970$ \\
\hline
\end{tabularx}
\end{center}

On rappelle que le domaine visible s'étend approximativement de $400$ nm (violet) à $800$ nm (rouge), que l'ultraviolet (UV) correspond à $\lambda < 400$ nm et l'infrarouge (IR) à $\lambda > 800$ nm. On donne la \textbf{loi de Wien} :
\[ \lambda_{\max} \times T = 2{,}90 \times 10^{-3}\ \text{m}\cdot\text{K} \]
où $T$ est la température absolue de la surface émettrice (en kelvins) et $\lambda_{\max}$ la longueur d'onde du maximum d'émission (en mètres).

Le club se souvient aussi du forgeron du quartier Melen : quand il chauffe une barre de fer dans sa forge, celle-ci devient d'abord rouge sombre, puis rouge vif, orange, et enfin presque blanche lorsqu'elle est très chaude. Comment expliquer ce changement de couleur ?

\textbf{Consignes}

Par groupes de quatre élèves, et en rédigeant soigneusement chaque étape :
\begin{enumerate}
\item Pour chacune des quatre sources, relever sur le profil spectral (ou dans le tableau) la valeur de $\lambda_{\max}$ et la convertir en mètres.
\item Calculer la température de surface $T$ de chaque source à l'aide de la loi de Wien. Vérifier l'homogénéité de la formule par une analyse dimensionnelle rapide.
\item Indiquer, pour chaque source, dans quel domaine des ondes électromagnétiques (UV, visible, IR) se situe le maximum d'émission. Pour le Soleil, préciser à quelle couleur du spectre visible correspond $\lambda_{\max}$.
\item Classer les quatre sources de la plus froide à la plus chaude. En déduire la relation générale entre la couleur dominante d'une source thermique et sa température.
\item Expliquer, à l'aide des résultats précédents, pourquoi le fer du forgeron passe du rouge au blanc lorsque sa température augmente, et pourquoi une lampe à incandescence est peu efficace pour l'éclairage.
\item Rédiger une conclusion collective (10 lignes maximum) puis la présenter oralement en 3 minutes ; les autres groupes évaluent l'exposé à l'aide de la grille fournie (exactitude des calculs, rigueur du raisonnement, qualité de la présentation).
\end{enumerate}

\textbf{Ressources à mobiliser}

\begin{aretenir}[Ce qu'il faut avoir en tête]
\begin{itemize}
\item \textbf{Loi de Wien} : $\lambda_{\max} \cdot T = 2{,}90 \times 10^{-3}\ \text{m}\cdot\text{K}$, donc $T = \dfrac{2{,}90 \times 10^{-3}}{\lambda_{\max}}$ avec $\lambda_{\max}$ en mètres.
\item Conversion : $1\ \text{nm} = 10^{-9}\ \text{m}$.
\item Domaines : UV ($\lambda < 400$ nm) ; visible ($400$ nm à $800$ nm, du violet au rouge) ; IR ($\lambda > 800$ nm).
\item Plus un corps chauffé est chaud, plus son maximum d'émission se déplace vers les \textbf{courtes} longueurs d'onde (du rouge vers le bleu).
\end{itemize}
\end{aretenir}

\textbf{Démarche et corrigé commenté}

\begin{exoresolu}[Météo des étoiles : corrigé complet]
Relever $\lambda_{\max}$ pour chaque source, calculer $T$ par la loi de Wien, situer le domaine spectral, classer les sources et expliquer la couleur du fer chauffé.
\tcblower
\textbf{Étape 1 : lecture et conversion de $\lambda_{\max}$.}

On convertit chaque longueur d'onde en mètres ($1\ \text{nm} = 10^{-9}\ \text{m}$) :
\begin{align*}
\text{Soleil : } & \lambda_{\max} = 500\ \text{nm} = 5{,}00 \times 10^{-7}\ \text{m} \\
\text{Étoile A : } & \lambda_{\max} = 830\ \text{nm} = 8{,}30 \times 10^{-7}\ \text{m} \\
\text{Étoile B : } & \lambda_{\max} = 290\ \text{nm} = 2{,}90 \times 10^{-7}\ \text{m} \\
\text{Lampe : } & \lambda_{\max} = 970\ \text{nm} = 9{,}70 \times 10^{-7}\ \text{m}
\end{align*}

\textbf{Étape 2 : application de la loi de Wien.}

La formule $T = \dfrac{2{,}90 \times 10^{-3}}{\lambda_{\max}}$ est homogène : le numérateur est en m$\cdot$K, le dénominateur en m, le quotient est bien en kelvins.

Pour le Soleil :
\[ T = \frac{2{,}90 \times 10^{-3}}{5{,}00 \times 10^{-7}} = 5{,}80 \times 10^{3}\ \text{K} \approx 5800\ \text{K}. \]

Pour l'étoile A (rouge) :
\[ T = \frac{2{,}90 \times 10^{-3}}{8{,}30 \times 10^{-7}} \approx 3{,}49 \times 10^{3}\ \text{K} \approx 3500\ \text{K}. \]

Pour l'étoile B (bleue) :
\[ T = \frac{2{,}90 \times 10^{-3}}{2{,}90 \times 10^{-7}} = 1{,}00 \times 10^{4}\ \text{K} = 10\,000\ \text{K}. \]

Pour la lampe à incandescence :
\[ T = \frac{2{,}90 \times 10^{-3}}{9{,}70 \times 10^{-7}} \approx 2{,}99 \times 10^{3}\ \text{K} \approx 3000\ \text{K}. \]

\textbf{Étape 3 : domaines spectraux.}
\begin{itemize}
\item Soleil : $\lambda_{\max} = 500$ nm $\in [400\,;\,800]$ nm : domaine \textbf{visible}, lumière verte-bleue (cyan) ; l'ensemble du spectre large donne une lumière blanche-jaune perçue.
\item Étoile A : $830$ nm $> 800$ nm : maximum dans l'\textbf{infrarouge}, mais la partie visible de son spectre est dominée par le rouge (queue de distribution).
\item Étoile B : $290$ nm $< 400$ nm : maximum dans l'\textbf{ultraviolet} ; sa partie visible est dominée par le bleu/violet.
\item Lampe : $970$ nm $> 800$ nm : maximum dans l'\textbf{infrarouge}.
\end{itemize}

\textbf{Étape 4 : classement et relation couleur--température.}

De la plus froide à la plus chaude :
\[ \text{Lampe (3000 K)} < \text{Étoile A rouge (3500 K)} < \text{Soleil (5800 K)} < \text{Étoile B bleue (10\,000 K)}. \]

Conclusion de l'étape : \textbf{plus la température d'une source thermique est élevée, plus son maximum d'émission se déplace vers les courtes longueurs d'onde} : rouge (source \og froide \fg), puis jaune-blanc, puis bleu (source très chaude). Contrairement à l'intuition courante (robinet rouge = chaud, bleu = froid), une étoile bleue est beaucoup plus chaude qu'une étoile rouge !

\textbf{Étape 5 : le fer du forgeron et la lampe.}

Quand le fer se réchauffe, $\lambda_{\max}$ diminue (loi de Wien) : le rayonnement émis passe de l'infrarouge (invisible, fer noir) au rouge ($T \approx 1000$ K), puis le spectre s'enrichit en orange et en jaune ; à très haute température, toutes les couleurs du visible sont émises avec des intensités comparables et le fer paraît \textbf{blanc}. C'est exactement le passage du rouge au blanc observé à la forge de Melen.

Pour la lampe à incandescence, le maximum est dans l'infrarouge ($970$ nm) : l'essentiel de l'énergie électrique est convertie en rayonnement thermique (chaleur) et non en lumière visible, d'où son faible rendement lumineux ($\approx 15$ lm/W) et son remplacement progressif par les DEL ($\approx 100$--$150$ lm/W), sources non thermiques beaucoup plus économes — un enjeu réel pour la facture d'électricité des ménages camerounais et la gestion du réseau \textsc{ENEO}.
\end{exoresolu}

\begin{attention}[Pièges fréquents au Probatoire]
\begin{itemize}
\item Ne jamais appliquer la loi de Wien avec $\lambda_{\max}$ en nanomètres : la convertir d'abord en mètres ($ \times 10^{-9} $) !
\item Ne pas confondre le langage courant et la physique : une étoile rouge est \emph{froide} (relativement, $\sim 3000$ K), une étoile bleue est \emph{chaude} ($> 10\,000$ K).
\item La loi de Wien ne s'applique qu'aux \textbf{sources thermiques} (corps chauffés au sens thermodynamique) : pas au laser, pas aux DEL, pas aux tubes fluorescents (luminescence).
\item Attention à l'unité de la constante de Wien : $b = 2,90 \times 10^{-3}\ \text{m}\cdot\text{K}$ (et non $\text{nm}\cdot\text{K}$).
\end{itemize}
\end{attention}

\textbf{Bilan de l'activité}

Cette activité a permis de mettre en évidence que :
\begin{itemize}
\item le \cle{profil spectral} d'une source thermique est une véritable \og carte d'identité \fg{} de sa température : la lecture de $\lambda_{\max}$ suffit, grâce à la loi de Wien, à déterminer $T$ à distance, sans aucun contact — c'est le principe de la \og météo des étoiles \fg{} (pyrométrie, astrophysique) ;
\item la couleur d'un corps chauffé est directement liée à sa température : du rouge (quelques milliers de kelvins) vers le blanc puis le bleu (dizaines de milliers de kelvins) ;
\item le maximum d'émission peut se situer hors du visible (UV ou IR) : l'œil ne perçoit qu'une partie du rayonnement, ce qui explique le faible rendement des lampes à incandescence (majoritairement IR) ;
\item la démarche expérimentale complète (observation, mesure, calcul, classement, conclusion argumentée, restitution orale) a été mobilisée, conformément à l'approche par les compétences : tu sais désormais analyser le spectre d'une source et en tirer des conclusions physiques rigoureuses.
\end{itemize}

\newpage

\coursec{Méthodes et exercices résolus}

\coursec{Spectres lumineux et sources de lumière}

\courssub{Sources de lumière et obtention d'un spectre}

\begin{definition}[Source de lumière, spectre lumineux]
Une \cle{source de lumière} est un objet qui émet de la lumière (source primaire) ou la renvoie (source secondaire). Le \cle{spectre lumineux} d'une source est la répartition de l'énergie lumineuse en fonction de la longueur d'onde $\lambda$ (ou de la fréquence $\nu$). On l'obtient en décomposant la lumière à l'aide d'un prisme ou d'un réseau de diffraction.
\end{definition}

\begin{methode}[Dispersion de la lumière par un prisme]
Pour obtenir et exploiter le spectre d'une source lumineuse :
\begin{enumerate}
\item \textbf{Réaliser un faisceau fin} : placer devant la source une fente étroite (ou un diaphragme) éclairée, afin d'obtenir un pinceau lumineux bien délimité.
\item \textbf{Placer le prisme} sur le trajet du faisceau, arête réfringente verticale, base en bas. Le faisceau doit frapper une face du prisme sous incidence oblique.
\item \textbf{Observer sur un écran} placé après le prisme : chaque radiation est déviée d'autant plus qu'elle est \emph{plus réfractée} par le verre.
\item \textbf{Se souvenir du sens de déviation} : le prisme dévie davantage le \cle{violet} que le \cle{rouge} (l'indice du verre diminue quand la longueur d'onde augmente : $n_{\text{violet}} > n_{\text{rouge}}$). C'est le phénomène de \cle{dispersion}.
\item \textbf{Conclure sur la nature de la lumière} :
\begin{itemize}
\item spectre \cle{continu} (bande multicolore sans interruption) $\Rightarrow$ source \cle{polychromatique} ;
\item une \cle{seule raie} (ou quelques raies colorées séparées) $\Rightarrow$ source monochromatique (une raie) ou à spectre de raies (plusieurs radiations).
\end{itemize}
\end{enumerate}
\textbf{Réflexe rédaction} : toujours préciser « dans le verre, l'indice dépend de la longueur d'onde : c'est le phénomène de \cle{dispersion} ».
\end{methode}

\begin{popfigure}
\begin{tikzpicture}[scale=1.1]
% Prism
\draw[thick, fill=popblue!10] (0,0) -- (3,1.5) -- (0,3) -- cycle;
% Apex angle label
\draw (0.3,0) arc (0:60:0.3) node[midway, right] {$\alpha$};
% Incident ray
\draw[poporange, thick, ->] (-1.5, 1.5) -- (0, 1.5) node[midway, above] {Lumière blanche};
% Normal (dashed)
\draw[dashed, thin, popmuted] (0,1.5) -- (-0.5, 1.5);
% Emergent rays (simplified geometry: parallel inside, diverging out)
\draw[poppurple, thick, ->] (2.8, 1.2) -- (4.5, 0.5) node[right] {Violet ($\lambda \approx 400\,\text{nm}$)};
\draw[popblue, thick, ->] (2.8, 1.3) -- (4.5, 0.9) node[right] {Bleu};
\draw[popgreen, thick, ->] (2.8, 1.4) -- (4.5, 1.3) node[right] {Vert};
\draw[popgold, thick, ->] (2.8, 1.5) -- (4.5, 1.7) node[right] {Jaune};
\draw[poporange, thick, ->] (2.8, 1.6) -- (4.5, 2.1) node[right] {Orange};
\draw[popred, thick, ->] (2.8, 1.7) -- (4.5, 2.5) node[right] {Rouge ($\lambda \approx 800\,\text{nm}$)};
% Screen
\draw[thick, decorate, decoration={zigzag, segment length=4pt, amplitude=2pt}] (4.5, 0.3) -- (4.5, 2.7) node[midway, right, align=left] {Écran\\Spectre continu};
% Internal rays (optional, show refraction at first face)
\draw[poporange, thick] (-0.1, 1.5) -- (2.8, 1.5);
\draw[popmuted, thin, ->] (1.5, 1.5) -- (1.5, 1.0) node[below] {Dispersion};
\end{tikzpicture}
\legende{Dispersion de la lumière blanche par un prisme : l'indice de réfraction $n$ diminue quand $\lambda$ augmente ($n_{\text{violet}} > n_{\text{rouge}}$), donc le violet est plus dévié que le rouge.}
\end{popfigure}

\begin{exoresolu}[Spectre d'une lampe à incandescence et d'un laser]
On éclaire successivement une fente fine avec : (a) une lampe à incandescence (lampe de poche à filament) ; (b) un pointeur laser rouge. Le faisceau traverse ensuite un prisme en verre et l'on observe sur un écran.
\begin{enumerate}
\item Décrire ce que l'on observe sur l'écran dans chaque cas.
\item Conclure sur la nature (monochromatique ou polychromatique) de chaque source.
\item Dans le cas (a), quelle couleur est la plus déviée ? Justifier.
\end{enumerate}
\tcblower
\textbf{1. Observations.}

(a) Avec la lampe à incandescence, on observe sur l'écran une \textbf{bande colorée continue}, allant du rouge au violet en passant par l'orange, le jaune, le vert et le bleu : c'est le \cle{spectre continu} de la lumière blanche.

(b) Avec le laser, on observe un \textbf{unique spot rouge}, simplement dévié par le prisme : aucune décomposition en plusieurs couleurs n'apparaît.

\textbf{2. Nature des sources.}

(a) Le spectre continu montre que la lumière de la lampe à incandescence est formée d'une \textbf{infinité de radiations} de longueurs d'onde différentes : c'est une source \cle{polychromatique}.

(b) Le laser n'émet qu'\textbf{une seule radiation}, caractérisée par une seule longueur d'onde dans le vide ($\lambda \approx \SI{650}{nm}$ pour un pointeur rouge) : c'est une source \cle{monochromatique}.

\textbf{3. Radiation la plus déviée.}

C'est le \textbf{violet} qui est le plus dévié. En effet, l'indice de réfraction du verre dépend de la longueur d'onde (dispersion) : il est d'autant plus grand que la longueur d'onde est petite. Comme $\lambda_{\text{violet}} < \lambda_{\text{rouge}}$, on a $n_{\text{violet}} > n_{\text{rouge}}$, donc d'après les lois de Snell-Descartes la radiation violette est plus réfractée, donc plus déviée que la rouge.

\textbf{Conclusion} : la lampe à incandescence est polychromatique (spectre continu), le laser est monochromatique (une seule radiation rouge), et le prisme dévie davantage le violet que le rouge.
\end{exoresolu}

\courssub{Sources monochromatiques et polychromatiques}

\begin{definition}[Radiation monochromatique, source monochromatique]
Une \cle{radiation monochromatique} est une onde électromagnétique sinusoïdale de fréquence $\nu$ unique (donc de longueur d'onde dans le vide $\lambda_0 = c/\nu$ unique). Une \cle{source monochromatique} n'émet qu'une seule radiation monochromatique. \textbf{Au Probatoire, on considère qu'un laser est une source monochromatique.}
\end{definition}

\begin{definition}[Source polychromatique]
Une \cle{source polychromatique} émet simultanément plusieurs radiations de longueurs d'onde différentes. Son spectre peut être \textbf{continu} (bande ininterrompue, ex. : Soleil, filament incandescent) ou \textbf{discontinu} (raies d'émission, ex. : lampe à vapeur de sodium, tube spectral).
\end{definition}

\begin{methode}[Longueur d'onde, fréquence, célérité]
\begin{enumerate}
\item \textbf{Identifier la grandeur caractéristique} : une radiation monochromatique est caractérisée par sa \cle{longueur d'onde dans le vide} $\lambda_0$ (en mètres, souvent en nanomètres : $\SI{1}{nm} = 10^{-9}\ \text{m}$).
\item \textbf{Utiliser la relation fondamentale} :
\[ \lambda_0 = \frac{c}{\nu} \qquad \text{soit} \qquad \nu = \frac{c}{\lambda_0} \]
avec $c = \SI{3,00e8}{m.s^{-1}}$ la célérité de la lumière dans le vide et $\nu$ la fréquence en hertz (Hz).
\item \textbf{Vérifier l'homogénéité} : $\dfrac{[\text{m.s}^{-1}]}{[\text{Hz}] = [\text{s}^{-1}]} = [\text{m}]$ : la relation est bien homogène.
\item \textbf{Attention aux conversions} : convertir systématiquement $\lambda$ en mètres avant tout calcul ($\SI{589}{nm} = 589 \times 10^{-9}\ \text{m}$).
\item \textbf{Retenir} : la fréquence $\nu$ d'une radiation est \emph{invariante} (elle ne change pas quand la lumière change de milieu) ; c'est la longueur d'onde qui change ($\lambda = \lambda_0 / n$ dans un milieu d'indice $n$).
\end{enumerate}
\end{methode}

\begin{exoresolu}[La radiation jaune du sodium]
Les lampes à vapeur de sodium, utilisées pour l'éclairage de certaines rues de Yaoundé, émettent une lumière jaune quasi monochromatique de longueur d'onde dans le vide $\lambda_0 = \SI{589}{nm}$.
\begin{enumerate}
\item Calculer la fréquence $\nu$ de cette radiation.
\item Cette radiation traverse du verre d'indice $n = 1,50$. Sa fréquence change-t-elle ? Calculer sa longueur d'onde $\lambda$ dans le verre.
\end{enumerate}
\tcblower
\textbf{1. Fréquence de la radiation.}

Conversion : $\lambda_0 = \SI{589}{nm} = 589 \times 10^{-9}\ \text{m} = \SI{5,89e-7}{m}$.

La relation entre longueur d'onde dans le vide et fréquence s'écrit :
\[ \nu = \frac{c}{\lambda_0} = \frac{3,00 \times 10^{8}}{5,89 \times 10^{-7}} \]
\[ \nu = \frac{3,00}{5,89} \times 10^{8-(-7)} = 0,509 \times 10^{15}\ \text{Hz} \]
\[ \boxed{\nu \approx \SI{5,09e14}{Hz}} \]

Vérification dimensionnelle : $\dfrac{\text{m.s}^{-1}}{\text{m}} = \text{s}^{-1} = \text{Hz}$. Correct.

\textbf{2. Passage dans le verre.}

La \textbf{fréquence} d'une radiation est une caractéristique de l'onde émise par la source : elle \textbf{ne change pas} lors du passage dans un autre milieu. Donc $\nu = \SI{5,09e14}{Hz}$ dans le verre.

En revanche, la célérité dans le verre vaut $v = \dfrac{c}{n}$, donc la longueur d'onde devient :
\[ \lambda = \frac{v}{\nu} = \frac{c}{n\,\nu} = \frac{\lambda_0}{n} = \frac{589}{1,50} \]
\[ \boxed{\lambda \approx \SI{393}{nm}} \]

\textbf{Conclusion} : la fréquence est invariante ($\nu \approx \SI{5,09e14}{Hz}$) ; la longueur d'onde dans le verre vaut environ $\SI{393}{nm}$, plus petite que dans le vide car la lumière s'y propage moins vite.
\end{exoresolu}

\courssub{Analyse des profils spectraux}

\begin{definition}[Profil spectral]
Un \cle{profil spectral} (ou distribution spectrale d'énergie) est la courbe donnant l'intensité lumineuse (ou l'éclairement relatif, ou la luminance) en fonction de la longueur d'onde $\lambda$. Il permet de visualiser la composition spectrale d'une source.
\end{definition}

\begin{methode}[Lecture d'un profil spectral]
Pour analyser un profil spectral :
\begin{enumerate}
\item \textbf{Repérer l'allure générale} :
\begin{itemize}
\item courbe \emph{continue et large} s'étendant sur tout le visible $\Rightarrow$ spectre continu (corps chaud : Soleil, filament) ;
\item \emph{pics fins} superposés ou non à un fond $\Rightarrow$ spectre de raies d'émission (lampe à vapeur, tube spectral, gaz sous faible pression) ;
\item spectre continu \emph{troué} de raies sombres $\Rightarrow$ spectre d'absorption (atmosphère d'une étoile, filtre).
\end{itemize}
\item \textbf{Identifier les longueurs d'onde des pics} : chaque raie correspond à une radiation émise par une espèce chimique ; c'est sa « signature ». Comparer aux valeurs de référence pour identifier l'élément.
\item \textbf{Repérer le maximum d'émission} $\lambda_{\max}$ pour un corps chauffé : il renseigne sur la température (loi de Wien, voir méthode 5).
\item \textbf{Conclure sur la composition} de la source : présence de telle ou telle raie $\Rightarrow$ présence de tel élément chimique.
\end{enumerate}
\textbf{Réflexe} : une DEL émet un pic relativement étroit mais centré sur une longueur d'onde dominante ; un laser donne un pic extrêmement fin (quasi monochromatique).
\end{methode}

\begin{popfigure}
\begin{tikzpicture}
\begin{axis}[
    width=\linewidth, height=7cm,
    xlabel={$\lambda$ (nm)}, ylabel={Intensité relative},
    xmin=380, xmax=780, ymin=0, ymax=1.1,
    xtick={380,450,500,550,600,650,700,780},
    ytick=\empty,
    axis lines=middle,
    enlargelimits=0.02,
    domain=380:780, samples=300,
    clip=false
]
% Continuous spectrum (blackbody approx 5800K) - broad peak ~500nm
\addplot[popblue, thick, domain=380:780] { exp(-(x-500)^2/8000) * (500/x)^2 };
% LED Green - narrow peak ~530nm
\addplot[popgreen, thick, domain=380:780] { 0.7*exp(-(x-530)^2/800) };
% Laser - very narrow peak (approximated as very narrow gaussian)
\addplot[popred, very thick, domain=531.5:532.5] { 1.0 };
% Emission lines (Mercury) - vertical lines
\addplot[poppurple, very thick, samples=2, domain=404.5:405.5] { 0.5 };
\addplot[poppurple, very thick, samples=2, domain=435.5:436.5] { 0.6 };
\addplot[poppurple, very thick, samples=2, domain=545.5:546.5] { 0.8 };
\addplot[poppurple, very thick, samples=2, domain=577.5:578.5] { 0.55 };
% Labels
\node[popblue, anchor=south] at (axis cs:500,1.05) {Continu (Soleil, filament)};
\node[popgreen, anchor=north] at (axis cs:530,0.75) {DEL verte};
\node[popred, anchor=north] at (axis cs:532,1.05) {Laser};
\node[poppurple, anchor=south] at (axis cs:436,0.65) {Raies Hg};
% Visible region shading
\fill[popblue!5] (380,0) rectangle (780,1.1);
\end{axis}
\end{tikzpicture}
\legende{Profils spectraux typiques : continu (corps chaud), bande étroite (DEL), raie fine (laser), raies d'émission (vapeur de mercure).}
\end{popfigure}

\begin{exoresolu}[Identification d'un gaz par son spectre]
On réalise le spectre de la lumière émise par une lampe à décharge contenant un gaz inconnu. On observe quatre raies brillantes sur fond noir, de longueurs d'onde : $\SI{405}{nm}$ (violette), $\SI{436}{nm}$ (bleue), $\SI{546}{nm}$ (verte) et $\SI{578}{nm}$ (jaune).

On dispose des raies de référence suivantes :
\begin{center}
\begin{tabular}{|l|c|}
\hline
\rowcolor{popblueL} Élément & Raies principales (nm) \\
\hline
Hydrogène & 410 ; 434 ; 486 ; 656 \\
\hline
Mercure & 405 ; 436 ; 546 ; 578 \\
\hline
Néon & 585 ; 614 ; 640 ; 703 \\
\hline
\end{tabular}
\end{center}
\begin{enumerate}
\item Ce spectre est-il continu ou de raies ? La source est-elle monochromatique ?
\item Identifier le gaz contenu dans la lampe. Justifier.
\end{enumerate}
\tcblower
\textbf{1. Nature du spectre.}

Le spectre est constitué de \textbf{raies brillantes isolées sur fond noir} : c'est un \cle{spectre de raies d'émission}, caractéristique d'un gaz sous faible pression traversé par une décharge électrique.

La source émet \textbf{plusieurs} radiations distinctes (quatre raies) : elle n'est donc \textbf{pas monochromatique} ; c'est une source \cle{polychromatique} à spectre discontinu (ou « spectre de raies »).

\textbf{2. Identification du gaz.}

Comparons les raies observées aux raies de référence :
\begin{itemize}
\item Hydrogène : raies à 410, 434, 486, $\SI{656}{nm}$ — aucune ne coïncide exactement avec les raies observées (405, 436, 546, 578) ;
\item Mercure : raies à 405, 436, 546, $\SI{578}{nm}$ — \textbf{coïncidence parfaite} avec les quatre raies observées ;
\item Néon : raies dans le rouge-orange (585, 614, 640, 703) — aucune correspondance.
\end{itemize}

Chaque élément chimique possède un spectre de raies qui lui est propre : c'est sa « carte d'identité » lumineuse (spectroscopie d'émission). La coïncidence des quatre raies impose la conclusion :

\textbf{Conclusion} : le gaz contenu dans la lampe est la \textbf{vapeur de mercure}. C'est d'ailleurs le principe des lampes fluorescentes (basse consommation) et de certains lampadaires publics.
\end{exoresolu}

\courssub{Les ondes électromagnétiques}

\begin{propriete}[Spectre électromagnétique]
La lumière visible n'est qu'une infime partie du spectre des ondes électromagnétiques. Toutes ces ondes se propagent dans le vide à la même célérité $c = \SI{3,00e8}{m.s^{-1}}$. Elles se distinguent par leur longueur d'onde $\lambda_0$ (ou leur fréquence $\nu$).
\begin{center}
\begin{tabularx}{\linewidth}{|l|X|X|}
\hline
\rowcolor{popblueL} \textbf{Domaine} & \textbf{Longueur d'onde $\lambda_0$ (approx.)} & \textbf{Applications / Exemples au Cameroun} \\
\hline
Rayons $\gamma$ & $\lambda < \SI{0,01}{nm}$ & Radiothérapie (hôpitaux), stérilisation \\
\hline
Rayons X & $\SI{0,01}{nm} \text{ à } \SI{10}{nm}$ & Radiographie médicale (hôpitaux), contrôle soudures (barrage Nachtigal) \\
\hline
Ultraviolet (UV) & $\SI{10}{nm} \text{ à } \SI{400}{nm}$ & Détection billets de banque, fluorescence, bronzage, stérilisation eau \\
\hline
\textbf{Visible} & \textbf{$\SI{400}{nm} \text{ à } \SI{800}{nm}$} & \textbf{Vision humaine, éclairage (ENEO), photosynthèse} \\
\hline
Infrarouge (IR) & $\SI{800}{nm} \text{ à } \SI{1}{mm}$ & Télécommandes, caméras thermiques, vision nocturne, fibre optique (CAMTEL) \\
\hline
Micro-ondes & $\SI{1}{mm} \text{ à } \SI{30}{cm}$ & Fours micro-ondes, \textbf{téléphonie mobile 4G/5G (MTN, Orange)}, radar, satellite \\
\hline
Ondes radio & $> \SI{30}{cm}$ & Radiodiffusion (CRTV), télévision, radio amateur, GPS \\
\hline
\end{tabularx}
\end{center}
\textbf{Ordre de grandeur à retenir} : $\lambda_{\text{visible}} \sim 10^{-7}\ \text{m}$ ; $\lambda_{\text{radio}} \sim 1\ \text{m}$ à $10^3\ \text{m}$.
\end{propriete}

\begin{methode}[Situer une onde électromagnétique dans son domaine]
\begin{enumerate}
\item \textbf{Connaître les limites du visible} : le spectre visible s'étend environ de $\SI{400}{nm}$ (violet) à $\SI{800}{nm}$ (rouge). En dessous de $\SI{400}{nm}$ : \cle{ultraviolet} (UV) ; au-dessus de $\SI{800}{nm}$ : \cle{infrarouge} (IR).
\item \textbf{Classer par longueur d'onde croissante} : rayons $\gamma$ — rayons X — UV — visible — IR — micro-ondes — ondes radio (radiofréquences).
\item \textbf{Convertir si nécessaire} en fréquence avec $\nu = c/\lambda_0$ : plus $\lambda$ est grande, plus $\nu$ est petite.
\item \textbf{Associer les applications} : rayons X (radiographie médicale), UV (détection, fluorescence), IR (télécommandes, caméras thermiques), micro-ondes (fours, téléphonie mobile MTN/Orange), ondes radio (radiodiffusion, télévision).
\end{enumerate}
\end{methode}

\begin{exoresolu}[Classement de radiations]
On considère les radiations suivantes, caractérisées par leur longueur d'onde dans le vide :
\begin{center}
\begin{tabular}{|l|c|}
\hline
\rowcolor{popblueL} Radiation & $\lambda_0$ \\
\hline
A (signal d'une antenne 4G) & $\SI{0,16}{m}$ \\
\hline
B (lumière d'une DEL verte) & $\SI{530}{nm}$ \\
\hline
C (rayonnement d'un scanner médical) & $\SI{0,030}{nm}$ \\
\hline
D (télécommande de téléviseur) & $\SI{940}{nm}$ \\
\hline
\end{tabular}
\end{center}
\begin{enumerate}
\item Attribuer à chaque radiation son domaine électromagnétique.
\item Calculer la fréquence de la radiation B.
\end{enumerate}
\tcblower
\textbf{1. Domaines.}

Comparons chaque longueur d'onde aux limites du visible ($\SI{400}{nm} < \lambda < \SI{800}{nm}$) :
\begin{itemize}
\item \textbf{A} : $\lambda_0 = \SI{0,16}{m} = 1,6 \times 10^{8}\ \text{nm} \gg \SI{800}{nm}$ : domaine des \textbf{ondes radio / micro-ondes} (radiofréquences utilisées par la téléphonie mobile MTN/Orange) ;
\item \textbf{B} : $\SI{400}{nm} < \SI{530}{nm} < \SI{800}{nm}$ : domaine du \textbf{visible} (lumière verte) ;
\item \textbf{C} : $\lambda_0 = \SI{0,030}{nm} \ll \SI{400}{nm}$, extrêmement petite : domaine des \textbf{rayons X} (imagerie médicale, scanner) ;
\item \textbf{D} : $\SI{940}{nm} > \SI{800}{nm}$ : domaine de l'\textbf{infrarouge} (proche IR, télécommandes).
\end{itemize}

\textbf{2. Fréquence de la radiation B.}

\[ \nu = \frac{c}{\lambda_0} = \frac{3,00 \times 10^{8}}{530 \times 10^{-9}} = \frac{3,00}{5,30} \times 10^{8+7} \]
\[ \boxed{\nu \approx \SI{5,66e14}{Hz}} \]

\textbf{Conclusion} : A est une radiofréquence, B est visible (vert, $\nu \approx \SI{5,66e14}{Hz}$), C est un rayonnement X, D est un infrarouge.
\end{exoresolu}

\courssub{Couleur des corps chauffés et loi de Wien}

\begin{savaistu}[Du forgeron de Douala aux étoiles]
Depuis l'Antiquité, les forgerons (comme ceux du quartier Bessengue à Douala) savent juger la température d'un métal à sa couleur : rouge sombre ($\sim 800\,\text{K}$), rouge cerise ($\sim 1000\,\text{K}$), orange ($\sim 1200\,\text{K}$), jaune ($\sim 1500\,\text{K}$), blanc ($\sim 2000\,\text{K}$). En astrophysique, la loi de Wien permet de déterminer la température de surface des étoiles rien qu'en analysant leur couleur : le Soleil (jaune-blanc, $\sim 5800\,\text{K}$), Bételgeuse (rouge, $\sim 3500\,\text{K}$), Rigel (bleu, $\sim 12000\,\text{K}$).
\end{savaistu}

\begin{propriete}[Loi de Wien — Démonstration non exigible au Probatoire]
Pour un \cle{corps noir} (absorbeur parfait, émetteur thermique idéal), la longueur d'onde $\lambda_{\max}$ correspondant au maximum de l'émissivité spectrale (ou de l'intensité lumineuse émise) est liée à la température absolue $T$ par :
\[ \boxed{\lambda_{\max} \times T = 2,90 \times 10^{-3}\ \text{m.K}} \quad \text{(constante de Wien)} \]
\textbf{Conditions impératives} :
\begin{itemize}
\item $T$ en \cle{kelvins} (K) : convertir avec $T(\text{K}) = \theta(°\text{C}) + 273$ ;
\item $\lambda_{\max}$ en \textbf{mètres}.
\end{itemize}
\textbf{Conséquence} : plus le corps est chaud, plus $\lambda_{\max}$ est petite : le maximum se déplace du rouge vers le bleu quand la température augmente.
\end{propriete}

\begin{popfigure}
\pgfmathdeclarefunction{wien}{1}{\pgfmathparse{2.898e6/#1}}
\pgfmathdeclarefunction{shape}{2}{\pgfmathparse{exp( - ( #1 - wien(#2) )^2 / ( 0.15 * wien(#2) )^2 ) * (wien(#2)/#1)^5 * (#2/5800)^5}}
\begin{tikzpicture}
\begin{axis}[
    width=\linewidth, height=7cm,
    xlabel={$\lambda$ (nm)}, ylabel={Émissivité spectrale (u.a.)},
    xmin=200, xmax=1500, ymin=0, ymax=1.1,
    xtick={400, 500, 700, 1000},
    ytick=\empty,
    axis lines=middle,
    enlargelimits=0.02,
    domain=200:1500, samples=300,
]
% Visible band
\addplot[popblue!10, draw=none, domain=400:800] {1.1} \closedcycle;
\node[popblue, anchor=south, rotate=90] at (axis cs:380,0.5) {Visible};

% T = 3000 K (Betelgeuse)
\addplot[popred, thick] {shape(x,3000)};
\node[popred, anchor=south] at (axis cs:966, {shape(966,3000)+0.05}) {$T=3000\,\text{K}$ (Bételgeuse)};

% T = 4000 K
\addplot[poporange, thick] {shape(x,4000)};
\node[poporange, anchor=south] at (axis cs:725, {shape(725,4000)+0.05}) {$T=4000\,\text{K}$};

% T = 5800 K (Sun)
\addplot[popgold, thick] {shape(x,5800)};
\node[popgold, anchor=south] at (axis cs:500, {shape(500,5800)+0.05}) {$T=5800\,\text{K}$ (Soleil)};

% Wien displacement line (dashed)
\draw[popmuted, dashed, thick] (axis cs:500,0) -- (axis cs:500,1.05);
\draw[popmuted, dashed, thick] (axis cs:725,0) -- (axis cs:725,0.85);
\draw[popmuted, dashed, thick] (axis cs:966,0) -- (axis cs:966,0.6);
\end{axis}
\end{tikzpicture}
\legende{Déplacement du maximum d'émission vers les courtes longueurs d'onde quand la température augmente (loi de Wien). La zone bleue représente le domaine visible.}
\end{popfigure}

\begin{methode}[Exploiter la loi de Wien]
\begin{enumerate}
\item \textbf{Écrire la loi} : $\lambda_{\max} \cdot T = 2,90 \times 10^{-3}\ \text{m.K}$.
\item \textbf{Convertir les données} : $\lambda_{\max}$ en \textbf{mètres}, $T$ en \textbf{kelvins} (K).
\item \textbf{Isoler la grandeur cherchée} : $\lambda_{\max} = \dfrac{2,90 \times 10^{-3}}{T}$ ou $T = \dfrac{2,90 \times 10^{-3}}{\lambda_{\max}}$.
\item \textbf{Calculer} en notant les puissances de 10 séparément.
\item \textbf{Interpréter} : situer $\lambda_{\max}$ par rapport au domaine visible (400–800 nm) pour déduire la couleur apparente.
\end{enumerate}
\end{methode}

\begin{exoresolu}[Température de surface du Soleil et couleur d'une étoile]
\textbf{Données} : constante de Wien : $\lambda_{\max} \cdot T = 2,90 \times 10^{-3}\ \text{m.K}$ ; $c = \SI{3,00e8}{m.s^{-1}}$.
\begin{enumerate}
\item Le spectre du Soleil présente un maximum d'émission pour $\lambda_{\max} = \SI{500}{nm}$. Calculer la température de surface du Soleil.
\item L'étoile Bételgeuse, qui apparaît rougeâtre dans le ciel, a une température de surface $T' = \SI{3500}{K}$. Calculer son $\lambda'_{\max}$ et situer ce maximum dans un domaine du spectre. Justifier sa couleur.
\item Un forgeron de Douala chauffe une barre de fer. Initialement rouge, elle devient jaune-orangé lorsqu'il prolonge le chauffage. Expliquer.
\end{enumerate}
\tcblower
\textbf{1. Température de surface du Soleil.}

D'après la loi de Wien : $T = \dfrac{2,90 \times 10^{-3}}{\lambda_{\max}}$ avec $\lambda_{\max} = \SI{500}{nm} = 5,00 \times 10^{-7}\ \text{m}$.
\[ T = \frac{2,90 \times 10^{-3}}{5,00 \times 10^{-7}} = \frac{2,90}{5,00} \times 10^{-3+7} = 0,580 \times 10^{4}\ \text{K} \]
\[ \boxed{T \approx \SI{5,80e3}{K} \approx \SI{5800}{K}} \]

\textbf{2. Maximum d'émission de Bételgeuse.}

\[ \lambda'_{\max} = \frac{2,90 \times 10^{-3}}{T'} = \frac{2,90 \times 10^{-3}}{3,50 \times 10^{3}} = 8,29 \times 10^{-7}\ \text{m} \]
\[ \boxed{\lambda'_{\max} \approx \SI{829}{nm}} \]

Cette valeur est \textbf{supérieure à $\SI{800}{nm}$} : le maximum se situe dans le \textbf{proche infrarouge}, mais la « queue » du spectre dans le visible est dominée par les grandes longueurs d'onde (rouge). L'étoile nous apparaît donc \textbf{rougeâtre}, ce qui est cohérent avec une température plus faible que celle du Soleil.

\textbf{3. Couleur du fer chauffé.}

Lorsque la température $T$ du fer augmente, la loi de Wien ($\lambda_{\max} \cdot T = \text{constante}$) montre que $\lambda_{\max}$ \textbf{diminue} : le maximum d'émission se déplace des grandes longueurs d'onde (rouge) vers des longueurs d'onde plus courtes (orange, jaune). De plus, l'intensité totale émise augmente fortement (loi de Stefan-Boltzmann, hors programme). La barre passe donc du rouge sombre au jaune-orangé, puis au blanc si le chauffage continue.

\textbf{Conclusion} : la loi de Wien relie quantitativement la couleur d'un corps chauffé à sa température : $T_{\text{Soleil}} \approx \SI{5800}{K}$, et Bételgeuse ($\SI{3500}{K}$) émet surtout dans l'infrarouge, d'où sa teinte rouge.
\end{exoresolu}

\courssub{Démarche expérimentale : lumières colorées}

\begin{experience}[Mettre en évidence la composition d'une lumière colorée]
\textbf{Matériel} : source de lumière colorée (projecteur à filtre, DEL, laser), fente fine, prisme (ou réseau de diffraction), écran blanc.
\textbf{Protocole} :
\begin{enumerate}
\item Placer la fente devant la source pour obtenir un faisceau fin.
\item Faire traverser le faisceau par le prisme (arête verticale).
\item Observer le spectre formé sur l'écran.
\item \emph{Variante} : interposer un filtre coloré (ex. rouge, vert, bleu) sur le trajet d'une lumière blanche (lampe à incandescence), puis disperser la lumière transmise.
\end{enumerate}
\textbf{Observations et interprétation} :
\begin{itemize}
\item \textbf{Une seule raie} de la couleur observée $\Rightarrow$ lumière \cle{monochromatique} (ex. laser, DEL de bonne qualité).
\item \textbf{Plusieurs raies} ou \textbf{bande continue étroite} $\Rightarrow$ lumière \cle{polychromatique} dont la couleur perçue résulte de la superposition de radiations (ex. DEL standard, lumière filtrée).
\item \textbf{Filtre} : un filtre \emph{absorbe} certaines radiations et \emph{transmet} les autres. Un filtre rouge transmet le rouge ($\lambda > \SI{600}{nm}$) et absorbe le vert/bleu. Disperser la lumière transmise permet de vérifier les longueurs d'onde transmises.
\end{itemize}
\textbf{Conclusion} : on distingue toujours la \emph{couleur perçue} (sensation physiologique) de la \emph{composition spectrale} (réalité physique mesurable).
\end{experience}

\begin{exoresolu}[La lumière d'une DEL verte est-elle monochromatique ?]
Lors d'une séance de travaux pratiques au lycée, un groupe d'élèves étudie la lumière émise par une DEL verte. Ils envoient le faisceau sur une fente fine puis sur un prisme. Sur l'écran, ils observent une bande verte étroite mais continue, s'étendant approximativement de $\SI{510}{nm}$ à $\SI{560}{nm}$, centrée vers $\SI{530}{nm}$. Ils recommencent avec un laser vert : ils n'observent alors qu'une seule raie à $\SI{532}{nm}$.
\begin{enumerate}
\item La DEL est-elle une source rigoureusement monochromatique ? Justifier.
\item Comparer qualitativement les deux sources.
\item Les élèves interposent un filtre rouge entre la DEL et le prisme. Prévoir qualitativement l'observation.
\end{enumerate}
\tcblower
\textbf{1. Nature de la lumière de la DEL.}

Le spectre de la DEL est une \textbf{bande continue étroite} (de 510 à $\SI{560}{nm}$) et non une raie unique : la DEL émet donc un \textbf{ensemble de radiations} de longueurs d'onde voisines. Elle n'est \textbf{pas rigoureusement monochromatique} ; on dit qu'elle est \emph{quasi monochromatique}, caractérisée par sa longueur d'onde dominante $\lambda_{\text{dom}} \approx \SI{530}{nm}$.

\textbf{2. Comparaison.}

Le laser vert donne \textbf{une seule raie} à $\SI{532}{nm}$ : il est (quasi parfaitement) \textbf{monochromatique}. La DEL possède une bande spectrale plus large : sa lumière est moins « pure » spectralement, même si l'œil perçoit les deux lumières comme vertes. Cela illustre la différence entre couleur perçue et composition spectrale.

\textbf{3. Interposition d'un filtre rouge.}

Un filtre rouge ne transmet que les radiations rouges (grandes longueurs d'onde, $\lambda \gtrsim \SI{600}{nm}$) et absorbe les autres. Or la DEL n'émet \textbf{aucune radiation rouge} (son spectre s'arrête vers $\SI{560}{nm}$). Le filtre absorbe donc pratiquement toute la lumière : \textbf{l'écran reste éteint} (ou très faiblement éclairé).

\textbf{Conclusion} : la DEL verte est quasi monochromatique (bande 510--$\SI{560}{nm}$) contrairement au laser (raie unique à $\SI{532}{nm}$), et un filtre rouge bloque sa lumière car il n'y a aucune composante rouge à transmettre.
\end{exoresolu}

\begin{attention}[Les 5 erreurs qui coûtent des points au Probatoire]
\begin{enumerate}
\item \textbf{Oublier de convertir les unités.} Dans $\nu = c/\lambda_0$ et dans la loi de Wien, la longueur d'onde doit être en \textbf{mètres} : $\SI{589}{nm} = 5,89 \times 10^{-7}\ \text{m}$. Une erreur de conversion fausse tout le calcul et fait perdre les points d'application numérique.
\item \textbf{Utiliser la température en degrés Celsius dans la loi de Wien.} La loi $\lambda_{\max} \cdot T = 2,90 \times 10^{-3}\ \text{m.K}$ exige $T$ en \textbf{kelvins} : $T(\text{K}) = \theta(°\text{C}) + 273$. Écrire $T = 25$ au lieu de $\SI{298}{K}$ est une faute de condition d'application.
\item \textbf{Inverser le sens de déviation du prisme.} C'est le \textbf{violet} qui est le plus dévié, pas le rouge, car l'indice du verre est plus grand pour les petites longueurs d'onde. Beaucoup de copies affirment le contraire.
\item \textbf{Confondre « monochromatique » et « colorée ».} Une lumière colorée n'est pas forcément monochromatique : seule l'analyse du spectre (une seule raie) permet de conclure. Une DEL verte est \emph{quasi} monochromatique, pas rigoureusement.
\item \textbf{Dire que la fréquence change de milieu.} Lors d'un changement de milieu, c'est la \textbf{longueur d'onde} (et la célérité) qui change ; la \textbf{fréquence} reste constante. Écrire $\lambda = \lambda_0/n$ et non $\nu = \nu_0/n$.
\end{enumerate}
\end{attention}

\newpage

\coursec{Exercices}

\courssub{Je vérifie mes connaissances}

\begin{exercice}[QCM : une seule réponse exacte]
Pour chaque question, choisir la bonne réponse en justifiant brièvement.
\begin{enumerate}
\item Le spectre de la lumière blanche obtenu avec un prisme est :
\begin{enumerate}
\item un spectre de raies d'émission ;
\item un spectre continu d'émission ;
\item un spectre de bandes d'absorption.
\end{enumerate}
\item Une source monochromatique est caractérisée par :
\begin{enumerate}
\item une seule couleur perçue ;
\item une seule longueur d'onde dans le vide ;
\item une seule fréquence audible.
\end{enumerate}
\item La lumière d'un laser rouge de longueur d'onde $\lambda = \SI{633}{nm}$ a pour fréquence (avec $c = \SI{3,00e8}{m.s^{-1}}$) :
\begin{enumerate}
\item $\nu = \SI{4,74e14}{Hz}$ ;
\item $\nu = \SI{4,74e11}{Hz}$ ;
\item $\nu = \SI{2,11e-15}{Hz}$.
\end{enumerate}
\item Quand la température d'un corps chauffé augmente, la longueur d'onde $\lambda_{max}$ de son maximum d'émission :
\begin{enumerate}
\item augmente ;
\item diminue ;
\item ne change pas.
\end{enumerate}
\end{enumerate}
\end{exercice}

\begin{exercice}[Vrai ou faux ? Justifier]
Répondre par vrai ou faux en justifiant chaque réponse par un argument physique.
\begin{enumerate}
\item La lumière émise par une lampe à vapeur de sodium est polychromatique.
\item Le spectre d'une étoile permet de déterminer la composition chimique de son atmosphère.
\item Dans le vide, la radiation rouge se propage plus vite que la radiation bleue.
\item Un corps chauffé à très haute température émet un maximum de lumière dans l'infrarouge.
\item Les ondes radio et la lumière visible sont de même nature physique.
\end{enumerate}
\end{exercice}

\begin{exercice}[Définitions et vocabulaire]
\begin{enumerate}
\item Définir : spectre lumineux ; source monochromatique ; source polychromatique ; profil spectral.
\item Citer deux sources de lumière monochromatiques et deux sources polychromatiques utilisées au quotidien au Cameroun.
\item Énoncer la loi de Wien en précisant la signification et l'unité SI de chaque grandeur.
\end{enumerate}
\end{exercice}

\begin{exercice}[Calculs directs]
On donne : $c = \SI{3,00e8}{m.s^{-1}}$ ; constante de Wien : $\lambda_{max} \times T = \SI{2,90e-3}{m.K}$.
\begin{enumerate}
\item Calculer la fréquence des radiations limites du domaine visible : $\lambda_1 = \SI{400}{nm}$ (violet) et $\lambda_2 = \SI{800}{nm}$ (rouge). En déduire le domaine de fréquences du visible.
\item Une étoile présente un maximum d'émission à $\lambda_{max} = \SI{435}{nm}$. Calculer sa température de surface en kelvins, puis en degrés Celsius. Prévoir sa couleur apparente.
\end{enumerate}
\end{exercice}

\courssub{Je m'entraîne}

\begin{exercice}[Classement des ondes électromagnétiques]
Classer par longueur d'onde croissante les ondes suivantes, puis calculer la grandeur manquante (longueur d'onde $\lambda$ ou fréquence $\nu$) pour chacune. On donne $c = \SI{3,00e8}{m.s^{-1}}$.
\begin{enumerate}
\item Les rayons X d'une radiographie à l'hôpital de district : $\nu = \SI{3,0e18}{Hz}$ ;
\item la lumière d'un laser rouge : $\lambda = \SI{633}{nm}$ ;
\item les ondes radio de la CRTV : $\nu = \SI{100}{MHz}$ ;
\item le rayonnement infrarouge d'une télécommande : $\lambda = \SI{940}{nm}$.
\end{enumerate}
\end{exercice}

\begin{exercice}[Lampe halogène]
Le filament de tungstène d'une lampe halogène utilisée pour l'éclairage d'un chantier à Douala est porté à la température $T = \SI{3000}{K}$. On donne $\lambda_{max} \times T = \SI{2,90e-3}{m.K}$.
\begin{enumerate}
\item Calculer la longueur d'onde $\lambda_{max}$ du maximum d'émission du filament.
\item Situer ce maximum dans le spectre des ondes électromagnétiques.
\item Expliquer pourquoi, malgré ce résultat, la lampe paraît émettre une lumière blanche-jaune.
\item Le filament d'une lampe à incandescence classique est à $\SI{2700}{K}$. Comparer qualitativement la couleur des deux lampes.
\end{enumerate}
\end{exercice}

\begin{exercice}[Identification de lampes par leurs profils spectraux]
On donne les profils spectraux (intensité lumineuse en fonction de la longueur d'onde) de trois lampes :
\begin{itemize}
\item \textbf{Lampe A} : une courbe continue, en cloche, s'étendant de $\SI{400}{nm}$ à $\SI{800}{nm}$ avec un maximum vers $\SI{600}{nm}$ ;
\item \textbf{Lampe B} : un pic unique très étroit à $\SI{589}{nm}$ ;
\item \textbf{Lampe C} : un pic unique extrêmement fin et intense à $\SI{650}{nm}$, sans aucune autre émission.
\end{itemize}
\begin{enumerate}
\item Identifier chaque lampe parmi : lampe à incandescence, lampe à vapeur de sodium, laser.
\item Justifier chaque identification en exploitant la forme du profil spectral.
\item Préciser quelle(s) lampe(s) est (sont) monochromatique(s).
\item Calculer la fréquence de la radiation émise par la lampe B. On donne $c = \SI{3,00e8}{m.s^{-1}}$.
\end{enumerate}
\end{exercice}

\begin{exercice}[Spectre d'une lampe à vapeur de mercure]
Un professeur du lycée de Nkongsamba veut faire observer à ses élèves le spectre de la lumière émise par une lampe à vapeur de mercure.
\begin{enumerate}
\item Décrire le protocole expérimental permettant d'obtenir ce spectre à l'aide d'un prisme : matériel nécessaire, schéma légendé du dispositif, étapes de la manipulation.
\item Préciser le rôle de la fente fine placée devant la lampe et celui du prisme.
\item Le spectre observé présente des raies brillantes (notamment une raie violette vers $\SI{436}{nm}$, une raie verte vers $\SI{546}{nm}$ et une raie jaune vers $\SI{578}{nm}$) sur fond noir. Quel type de spectre obtient-on ? Que révèle-t-il sur la nature de la lumière émise ?
\item Calculer la fréquence de la raie verte. On donne $c = \SI{3,00e8}{m.s^{-1}}$.
\end{enumerate}
\end{exercice}

\courssub{Je me prépare au Probatoire}

\begin{exercice}[Les étoiles, des corps chauffés]
\textit{(D'après Probatoire C — durée conseillée : 30 min)}

On donne : $c = \SI{3,00e8}{m.s^{-1}}$ ; $h = \SI{6,63e-34}{J.s}$ ; loi de Wien : $\lambda_{max} \times T = \SI{2,90e-3}{m.K}$ ; $T(\text{K}) = \theta(°C) + 273$.

Les étoiles se comportent, en première approximation, comme des corps chauffés. Leur spectre, obtenu à l'aide d'un spectroscope à prisme, présente un fond continu surmonté de raies sombres.

\begin{enumerate}
\item \textbf{Obtention du spectre.}
\begin{enumerate}
\item Décrire le principe d'obtention du spectre de la lumière d'une étoile à l'aide d'un prisme. Faire un schéma légendé du dispositif.
\item Quel phénomène physique est mis en jeu lors de la traversée du prisme ? Pourquoi les différentes radiations sont-elles séparées ?
\end{enumerate}
\item \textbf{Température de surface.} L'étoile Bételgeuse (constellation d'Orion) présente un maximum d'émission à $\lambda_{max} = \SI{830}{nm}$.
\begin{enumerate}
\item Calculer sa température de surface en kelvins, puis en degrés Celsius.
\item Dans quel domaine des ondes électromagnétiques se situe ce maximum ?
\item Prévoir la couleur apparente de Bételgeuse. Justifier.
\end{enumerate}
\item \textbf{Comparaison avec le Soleil.} Le Soleil a une température de surface $T_S = \SI{5800}{K}$.
\begin{enumerate}
\item Calculer $\lambda_{max}$ pour le Soleil et situer cette radiation dans le spectre.
\item Comparer les couleurs apparentes du Soleil et de Bételgeuse. Quelle relation générale peut-on énoncer entre la température d'une étoile et sa couleur ?
\end{enumerate}
\item \textbf{Composition chimique.} Le spectre du Soleil présente des raies sombres (raies de Fraunhofer). Expliquer l'origine de ces raies et ce qu'elles permettent de déterminer.
\end{enumerate}
\end{exercice}

\begin{exercice}[Caméras thermiques et détection de la fièvre]
\textit{(Durée conseillée : 25 min)}

On donne : $c = \SI{3,00e8}{m.s^{-1}}$ ; loi de Wien : $\lambda_{max} \times T = \SI{2,90e-3}{m.K}$ ; $T(\text{K}) = \theta(°C) + 273$.

Lors d'une épidémie, des caméras thermiques sont installées à l'aéroport international de Douala pour détecter les passagers fiévreux. Le corps humain émet un rayonnement assimilable à celui d'un corps chauffé.

\begin{enumerate}
\item La température normale du corps humain est $\theta_1 = \SI{37}{\celsius}$.
\begin{enumerate}
\item Convertir cette température en kelvins.
\item Calculer la longueur d'onde $\lambda_{max,1}$ du maximum d'émission du corps humain.
\item Dans quel domaine des ondes électromagnétiques se situe ce rayonnement ? Pourquoi ne le voit-on pas à l'œil nu ?
\end{enumerate}
\item Un passager fiévreux a une température corporelle $\theta_2 = \SI{39}{\celsius}$. Calculer $\lambda_{max,2}$ et comparer à $\lambda_{max,1}$. Conclure sur l'évolution de $\lambda_{max}$ quand la température augmente.
\item Expliquer le principe de fonctionnement d'une caméra thermique : quel rayonnement détecte-t-elle et comment distingue-t-elle un passager fiévreux d'un passager sain ?
\item Calculer les fréquences $\nu_1$ et $\nu_2$ associées à $\lambda_{max,1}$ et $\lambda_{max,2}$.
\item Un technicien affirme : « La caméra thermique mesure la température en émettant un rayonnement infrarouge vers le passager. » Cette affirmation est-elle correcte ? Justifier.
\end{enumerate}
\end{exercice}

\begin{exercice}[Éclairage public et économie d'énergie à Yaoundé]
\textit{(Durée conseillée : 30 min)}

On donne : $c = \SI{3,00e8}{m.s^{-1}}$ ; loi de Wien : $\lambda_{max} \times T = \SI{2,90e-3}{m.K}$ ; domaine du visible : de $\SI{400}{nm}$ à $\SI{800}{nm}$.

La mairie de Yaoundé remplace progressivement les anciennes lampes à vapeur de sodium de l'éclairage public par des lampes à DEL. Un élève de Première C analyse les profils spectraux des deux types de lampes.

\begin{enumerate}
\item \textbf{La lampe à vapeur de sodium.} Son profil spectral présente essentiellement une raie jaune intense à $\lambda = \SI{589}{nm}$.
\begin{enumerate}
\item Cette lampe est-elle monochromatique ou polychromatique ? Justifier.
\item Calculer la fréquence de la radiation émise.
\item Expliquer pourquoi les objets éclairés par cette lampe paraissent tous jaunâtres, quelle que soit leur couleur réelle.
\end{enumerate}
\item \textbf{La lampe à DEL blanche.} Son profil spectral présente un pic étroit dans le bleu vers $\SI{450}{nm}$ et une large bande continue de $\SI{500}{nm}$ à $\SI{700}{nm}$.
\begin{enumerate}
\item Cette lampe est-elle monochromatique ? Justifier.
\item Expliquer pourquoi les objets éclairés par cette lampe retrouvent des couleurs proches de leurs couleurs en plein jour.
\end{enumerate}
\item \textbf{Comparaison avec une lampe à incandescence.} Une vieille lampe à incandescence a un filament porté à $T = \SI{2800}{K}$.
\begin{enumerate}
\item Calculer $\lambda_{max}$ et situer ce maximum dans le spectre électromagnétique.
\item Quelle fraction de l'énergie émise se trouve en dehors du domaine visible ? Justifier qualitativement à l'aide du profil spectral en cloche d'un corps chauffé.
\item En déduire pourquoi les lampes à incandescence sont énergivores et pourquoi ENEO encourage leur remplacement par des DEL.
\end{enumerate}
\item \textbf{Synthèse.} Rédiger un court texte (5 à 8 lignes) comparant les trois sources (sodium, DEL, incandescence) du point de vue de la nature de leur lumière, de leur spectre et de leur efficacité pour l'éclairage public.
\end{enumerate}
\end{exercice}

\newpage

\coursec{Corrigés des exercices}

\begin{corrige}[Exercice 1 — QCM]
\begin{enumerate}
\item \textbf{Réponse b : \cle{spectre continu d'émission}.} La lumière blanche contient toutes les radiations visibles ; le prisme les sépare en un dégradé continu du violet au rouge, sans raie ni bande sombre.
\item \textbf{Réponse b : une seule \cle{longueur d'onde dans le vide}.} C'est la définition rigoureuse d'une \cle{source monochromatique}. La couleur perçue (réponse a) est une conséquence physiologique, non la caractéristique physique ; la lumière n'est pas une onde sonore (réponse c).
\item \textbf{Réponse a.} $\nu = \dfrac{c}{\lambda} = \dfrac{\num{3,00e8}}{\num{633e-9}} = \SI{4,74e14}{Hz}$. Les réponses b et c correspondent à des erreurs de conversion du nanomètre en mètre.
\item \textbf{Réponse b : diminue.} D'après la \cle{loi de Wien} $\lambda_{\max} \times T = \text{constante}$, $\lambda_{\max}$ et $T$ sont inversement proportionnelles : plus le corps est chaud, plus son maximum d'émission se déplace vers les courtes longueurs d'onde.
\end{enumerate}
\end{corrige}

\begin{corrige}[Exercice 2 — Vrai ou faux]
\begin{enumerate}
\item \textbf{Faux.} La lampe à vapeur de sodium émet essentiellement une radiation jaune ($\lambda \approx \SI{589}{nm}$, en réalité un \cle{doublet} très proche à \SI{589,0}{nm} et \SI{589,6}{nm}) : sa lumière est \cle{quasi monochromatique}.
\item \textbf{Vrai.} Le spectre d'une étoile présente des \cle{raies d'absorption} (raies sombres) caractéristiques des éléments chimiques présents dans son atmosphère : chaque élément absorbe des longueurs d'onde précises qui constituent sa « signature » spectrale.
\item \textbf{Faux.} Dans le vide, toutes les \cle{ondes électromagnétiques} se propagent à la même \cle{célerité} $c = \SI{3,00e8}{m.s^{-1}}$, quelle que soit leur longueur d'onde. (C'est dans un milieu \cle{dispersif} comme le verre que les vitesses diffèrent.)
\item \textbf{Faux.} D'après la loi de Wien, quand $T$ augmente, $\lambda_{\max}$ diminue : le maximum d'émission se déplace vers le visible, puis vers l'ultraviolet, et s'\textbf{éloigne} de l'infrarouge.
\item \textbf{Vrai.} Ce sont toutes deux des ondes électromagnétiques : elles ne diffèrent que par leur \cle{fréquence} (et donc leur longueur d'onde).
\end{enumerate}
\end{corrige}

\begin{corrige}[Exercice 3 — Définitions et vocabulaire]
\begin{enumerate}
\item \textbf{\cle{Spectre lumineux}} : figure obtenue en décomposant une lumière à l'aide d'un système \cle{dispersif} (prisme ou réseau) ; elle montre les radiations qui composent cette lumière. \textbf{\cle{Source monochromatique}} : source émettant une seule radiation, caractérisée par une seule longueur d'onde dans le vide $\lambda$. \textbf{\cle{Source polychromatique}} : source émettant plusieurs radiations de longueurs d'onde différentes. \textbf{\cle{Profil spectral}} : courbe donnant l'intensité lumineuse émise par une source en fonction de la longueur d'onde.
\item \textbf{Monochromatiques} : le laser d'un pointeur, la DEL rouge d'un appareil électroménager (téléviseur, radio). \textbf{Polychromatiques} : la lampe torche à incandescence, la lampe à tube fluorescent (ou DEL blanche) des salles de classe, la lumière du Soleil.
\item \textbf{Loi de Wien} : $\lambda_{\max} \times T = \SI{2,90e-3}{m.K}$, où $\lambda_{\max}$ est la longueur d'onde du maximum d'émission du corps chauffé (en mètres, m) et $T$ sa \cle{température absolue} (en kelvins, K). La constante $\SI{2,90e-3}{m.K}$ s'exprime en mètres-kelvins.
\end{enumerate}
\end{corrige}

\begin{corrige}[Exercice 4 — Calculs directs]
\begin{enumerate}
\item Pour le violet ($\lambda_1 = \SI{400}{nm}$) : $\nu_1 = \dfrac{c}{\lambda_1} = \dfrac{\num{3,00e8}}{\num{400e-9}} = \SI{7,50e14}{Hz}$. Pour le rouge ($\lambda_2 = \SI{800}{nm}$) : $\nu_2 = \dfrac{c}{\lambda_2} = \dfrac{\num{3,00e8}}{\num{800e-9}} = \SI{3,75e14}{Hz}$. Comme $\nu = c/\lambda$ est une fonction décroissante de $\lambda$, le domaine de fréquences du visible s'étend de $\SI{3,75e14}{Hz}$ (rouge) à $\SI{7,50e14}{Hz}$ (violet).
\item Loi de Wien : $T = \dfrac{\num{2,90e-3}}{\num{435e-9}} = \SI{6,67e3}{K}$, soit $\theta = 6667 - 273 \approx \SI{6,39e3}{\celsius}$. Le maximum se situe dans le bleu-violet ($\SI{435}{nm}$) : l'étoile paraît \textbf{bleutée}, car les courtes longueurs d'onde dominent son émission visible (exemple : Rigel, étoiles de type spectral B).
\end{enumerate}
\end{corrige}

\begin{corrige}[Exercice 5 — Classement des ondes électromagnétiques]
Calculs préalables avec $\lambda = c/\nu$ ou $\nu = c/\lambda$ ($c = \SI{3,00e8}{m.s^{-1}}$) :
\begin{itemize}
\item Rayons X : $\nu = \num{3,0e18}\,Hz \Rightarrow \lambda = \dfrac{\num{3,00e8}}{\num{3,0e18}} = \SI{1,0e-10}{m} = \SI{0,10}{nm}$ ;
\item Laser rouge : $\lambda = \SI{633}{nm} \Rightarrow \nu = \dfrac{\num{3,00e8}}{\num{633e-9}} = \SI{4,74e14}{Hz}$ ;
\item Radio CRTV : $\nu = \SI{100}{MHz} = \num{1,00e8}\,Hz \Rightarrow \lambda = \dfrac{\num{3,00e8}}{\num{1,00e8}} = \SI{3,00}{m}$ ;
\item Infrarouge (télécommande) : $\lambda = \SI{940}{nm} \Rightarrow \nu = \dfrac{\num{3,00e8}}{\num{940e-9}} = \SI{3,19e14}{Hz}$.
\end{itemize}
Classement par \cle{longueur d'onde croissante} :
\[ \text{rayons X } (\SI{0,10}{nm}) \;<\; \text{laser rouge } (\SI{633}{nm}) \;<\; \text{infrarouge } (\SI{940}{nm}) \;<\; \text{radio } (\SI{3,00}{m}). \]
\emph{Remarque :} l'infrarouge de la télécommande ($\SI{940}{nm}$) est bien au-delà du rouge visible ($\SI{800}{nm}$), ce qui confirme qu'il est invisible à l'œil nu.
\end{corrige}

\begin{corrige}[Exercice 6 — Lampe halogène]
\begin{enumerate}
\item $\lambda_{\max} = \dfrac{\num{2,90e-3}}{3000} = \SI{9,67e-7}{m} = \SI{967}{nm}$.
\item $\SI{967}{nm} > \SI{800}{nm}$ : le maximum d'émission se situe dans l'\textbf{infrarouge}, invisible à l'œil humain.
\item Le profil spectral du filament est une \cle{courbe en cloche} (loi de Planck) très large : bien que le maximum soit dans l'infrarouge, une part importante de l'émission couvre tout le domaine visible (de $\SI{400}{nm}$ à $\SI{800}{nm}$), avec une prédominance du rouge-orangé sur le bleu. La superposition de toutes ces radiations visibles donne une lumière blanche tirant sur le jaune.
\item À $\SI{2700}{K}$, $\lambda_{\max} = \dfrac{\num{2,90e-3}}{2700} = \SI{1074}{nm}$, encore plus profondément dans l'infrarouge. La lampe classique émet relativement moins de bleu que l'halogène : sa lumière paraît plus \textbf{jaune-orangée} (« chaude »), tandis que l'halogène, plus chaude, donne une lumière plus blanche (meilleur \cle{rendu des couleurs}).
\end{enumerate}
\end{corrige}

\begin{corrige}[Exercice 7 — Identification de lampes par leurs profils spectraux]
\begin{enumerate}
\item Lampe A = \textbf{lampe à incandescence} (filament) ; Lampe B = \textbf{lampe à vapeur de sodium} ; Lampe C = \textbf{laser}.
\item \textbf{Justifications :} A : courbe continue en cloche couvrant tout le visible, signature d'un \cle{corps chauffé} (filament) dont l'émission est continue (loi de Planck). B : pic unique à $\SI{589}{nm}$, \cle{raie jaune} caractéristique du sodium (doublet non résolu). C : pic extrêmement fin et intense à $\SI{650}{nm}$, sans aucune autre émission : seul le laser produit une radiation aussi \cle{pure} (bande passante très étroite).
\item Les lampes B et C sont \cle{monochromatiques} (une seule longueur d'onde émise, de façon quasi stricte pour le laser). La lampe A est \cle{polychromatique}.
\item $\nu = \dfrac{c}{\lambda} = \dfrac{\num{3,00e8}}{\num{589e-9}} = \SI{5,09e14}{Hz}$.
\end{enumerate}
\end{corrige}

\begin{corrige}[Exercice 8 — Spectre d'une lampe à vapeur de mercure]
\begin{enumerate}
\item \textbf{Matériel} : lampe à vapeur de mercure, support, \cle{fente fine} réglable éclairée par la lampe, prisme en verre, écran blanc, pièce obscurcie. \textbf{Protocole} : placer la fente fine devant la lampe ; disposer le prisme après la fente ; placer l'écran en sortie du prisme ; régler la position de l'écran et l'orientation du prisme jusqu'à observer des raies colorées nettes sur fond noir.
\begin{popfigure}
\begin{center}
\begin{tikzpicture}[scale=0.9]
\draw[fill=popgold!40,draw=popdark] (0,0.4) rectangle (0.5,1.6);
\node[below] at (0.25,0.4) {\footnotesize Lampe \text{Hg}};
\draw[fill=popdark] (1.8,0.2) rectangle (2.1,1.8);
\draw[white,line width=1.5pt] (1.95,0.85) -- (1.95,1.15);
\node[below] at (1.95,0.2) {\footnotesize Fente};
\draw[fill=popblueL,draw=popblue] (4.2,0.1) -- (5.2,0.1) -- (4.7,1.7) -- cycle;
\node[below] at (4.7,0.1) {\footnotesize Prisme};
\draw[fill=popcream,draw=popdark] (8.2,0) rectangle (8.5,2.2);
\node[below] at (8.35,0) {\footnotesize Écran};
\draw[popmuted,->] (0.5,1.0) -- (1.95,1.0);
\draw[popmuted,->] (2.1,1.0) -- (4.55,0.95);
\draw[popgold,thick] (4.85,0.9) -- (8.2,1.5);
\draw[popgreen,thick] (4.85,0.9) -- (8.2,1.1);
\draw[poppurple,thick] (4.85,0.9) -- (8.2,0.7);
\node[popgold,right] at (8.5,1.5) {\footnotesize jaune};
\node[popgreen,right] at (8.5,1.1) {\footnotesize verte};
\node[poppurple,right] at (8.5,0.7) {\footnotesize violette};
\end{tikzpicture}
\end{center}
\legende{Dispositif d'observation du spectre d'émission de la lampe à vapeur de mercure.}
\end{popfigure}
\item La \textbf{fente fine} sélectionne un pinceau lumineux étroit afin que les images des différentes radiations ne se recouvrent pas : chaque raie observée est l'image de la fente pour une longueur d'onde donnée. Le \textbf{prisme} dévie la lumière par \cle{réfraction} ; comme l'indice du verre dépend de la longueur d'onde (\cle{dispersion}), chaque radiation est déviée différemment (le violet plus que le rouge) : les radiations sont \cle{séparées}.
\item On obtient un \textbf{spectre de raies d'émission} : quelques raies colorées brillantes sur fond noir. La lumière émise est donc \textbf{polychromatique} mais constituée d'un nombre limité de radiations bien précises, caractéristiques de l'atome de mercure (\cle{spectre de raies}).
\item $\nu = \dfrac{c}{\lambda} = \dfrac{\num{3,00e8}}{\num{546e-9}} = \SI{5,49e14}{Hz}$ (raie verte du mercure).
\end{enumerate}
\end{corrige}

\begin{corrige}[Exercice 9 — Les étoiles, des corps chauffés]
\textbf{1. Obtention du spectre.}
\begin{enumerate}
\item La lumière de l'étoile, collectée par une lunette (ou télescope), traverse une fente fine puis un prisme (ou un réseau de diffraction) ; le spectre est observé sur un écran ou enregistré par un capteur (CCD). Le schéma est identique dans son principe à celui de l'exercice 8 (source $\to$ fente $\to$ élément dispersif $\to$ détecteur).
\item Le phénomène mis en jeu est la \textbf{réfraction}, avec \textbf{dispersion} : l'indice du verre dépend de la longueur d'onde ($n_{\text{violet}} > n_{\text{rouge}}$), donc chaque radiation est déviée d'un angle différent, ce qui sépare les radiations spatialement.
\end{enumerate}
\textbf{2. Température de Bételgeuse ($\alpha$ Orionis).}
\begin{enumerate}
\item $T = \dfrac{\num{2,90e-3}}{\num{830e-9}} = \SI{3,49e3}{K}$, soit $\theta = 3494 - 273 \approx \SI{3,22e3}{\celsius}$.
\item $\SI{830}{nm} > \SI{800}{nm}$ : le maximum se situe dans l'\textbf{infrarouge proche}.
\item Le maximum étant dans l'infrarouge, la partie visible de l'émission est dominée par les grandes longueurs d'onde (rouge-orangé) : Bételgeuse paraît \textbf{rouge-orangée}, ce que confirme son observation à l'œil nu dans la constellation d'Orion.
\end{enumerate}
\textbf{3. Comparaison avec le Soleil.}
\begin{enumerate}
\item $\lambda_{\max} = \dfrac{\num{2,90e-3}}{5800} = \SI{5,00e-7}{m} = \SI{500}{nm}$ : dans le \textbf{visible} (vert-bleu), au milieu du domaine de sensibilité de l'œil.
\item Le Soleil ($\SI{5800}{K}$, maximum à $\SI{500}{nm}$) émet de façon équilibrée toutes les radiations visibles : il paraît blanc-jaune. Bételgeuse, plus froide ($\SI{3500}{K}$), paraît rouge. Relation générale : \textbf{plus une étoile est chaude, plus son maximum d'émission se déplace vers les courtes longueurs d'onde et plus sa couleur tire vers le bleu ; plus elle est froide, plus elle paraît rouge.} (Séquence spectrale OBAFGKM).
\end{enumerate}
\textbf{4. Composition chimique.} Le fond continu est émis par la surface chaude de l'étoile (photosphère) ; en traversant son atmosphère plus froide, certains éléments (hydrogène, hélium, sodium, calcium, fer...) absorbent leurs radiations caractéristiques, créant des \textbf{raies sombres d'absorption} (raies de Fraunhofer). La position de ces raies permet d'\textbf{identifier les éléments chimiques présents} dans l'atmosphère de l'étoile (spectroscopie d'absorption).

\medskip
\textbf{Barème indicatif (sur 10)} : 1.a : 1,5 pt ; 1.b : 1 pt ; 2.a : 1,5 pt ; 2.b : 0,5 pt ; 2.c : 1 pt ; 3.a : 1 pt ; 3.b : 1,5 pt ; 4 : 2 pts.

\textbf{Points de vigilance} : citer la loi de Wien avant de l'appliquer ; convertir les nanomètres en mètres ; la température de la loi de Wien est obligatoirement en kelvins ; pour la couleur, raisonner sur la position de $\lambda_{\max}$ par rapport au domaine visible, pas seulement sur la valeur numérique.
\end{corrige}

\begin{corrige}[Exercice 10 — Caméras thermiques et détection de la fièvre]
\begin{enumerate}
\item 
\begin{enumerate}
\item $T_1 = 37 + 273 = \SI{310}{K}$.
\item $\lambda_{\max,1} = \dfrac{\num{2,90e-3}}{310} = \SI{9,35e-6}{m} = \SI{9350}{nm} = \SI{9,35}{\micro m}$.
\item Ce rayonnement se situe dans l'\textbf{infrarouge} (très au-delà de $\SI{800}{nm}$) : il est invisible à l'œil humain, dont la sensibilité s'arrête vers $\SI{800}{nm}$.
\end{enumerate}
\item $T_2 = 39 + 273 = \SI{312}{K}$ ; $\lambda_{\max,2} = \dfrac{\num{2,90e-3}}{312} = \SI{9,29e-6}{m} = \SI{9,29}{\micro m}$. On a $\lambda_{\max,2} < \lambda_{\max,1}$ : quand la température augmente, $\lambda_{\max}$ diminue (loi de Wien), et l'intensité totale émise augmente (loi de Stefan-Boltzmann).
\item La caméra thermique est un \textbf{récepteur passif} : elle détecte le rayonnement infrarouge \textbf{émis spontanément par le corps} du passager (tout corps à $T > 0\,\text{K}$ rayonne). Un passager fiévreux émet un rayonnement plus intense (et de maximum légèrement décalé vers les courtes longueurs d'onde) : la caméra mesure cette intensité (radiométrie), la convertit en température via un étalonnage et signale les températures supérieures au seuil.
\item $\nu_1 = \dfrac{c}{\lambda_{\max,1}} = \dfrac{\num{3,00e8}}{\num{9,35e-6}} = \SI{3,21e13}{Hz}$ ; $\nu_2 = \dfrac{c}{\lambda_{\max,2}} = \dfrac{\num{3,00e8}}{\num{9,29e-6}} = \SI{3,23e13}{Hz}$.
\item \textbf{Affirmation fausse.} La caméra thermique n'émet rien (ou seulement un très faible rayonnement propre) : elle \textbf{reçoit} le rayonnement infrarouge naturellement émis par le corps humain (corps noir à $\approx \SI{310}{K}$). C'est ce rayonnement \emph{reçu} qui permet de déduire la température de surface. C'est de la \cle{thermographie infrarouge passive}.
\end{enumerate}
\textbf{Barème indicatif (sur 10)} : 1.a : 0,5 pt ; 1.b : 1,5 pt ; 1.c : 1 pt ; 2 : 2 pts ; 3 : 2 pts ; 4 : 1,5 pt ; 5 : 1,5 pt.

\textbf{Points de vigilance} : conversion $^\circ\text{C} \to \text{K}$ obligatoire avant la loi de Wien ; distinguer clairement \cle{émetteur} (le corps) et \cle{récepteur} (la caméra) dans la question 5 ; exprimer $\lambda_{\max}$ en micromètres ($\mu\text{m}$) pour la lisibilité dans le domaine infrarouge.
\end{corrige}

\begin{corrige}[Exercice 11 — Éclairage public et économie d'énergie à Yaoundé]
\textbf{1. Lampe à vapeur de sodium (basse pression).}
\begin{enumerate}
\item Elle est (quasi) \textbf{monochromatique} : son profil ne présente qu'une seule raie intense à $\SI{589}{nm}$ (doublet non résolu), donc une seule longueur d'onde significative.
\item $\nu = \dfrac{c}{\lambda} = \dfrac{\num{3,00e8}}{\num{589e-9}} = \SI{5,09e14}{Hz}$.
\item Un objet ne peut \cle{diffuser} (réémettre) que les radiations qu'il reçoit (\cle{synthèse soustractive}). La lampe n'émettant que du jaune ($\SI{589}{nm}$) : un objet bleu, par exemple, ne reçoit aucune radiation bleue à diffuser et paraît sombre ou jaunâtre. Toutes les couleurs sont donc « écrasées » vers le jaune : le \cle{rendu des couleurs} est très mauvais (IRC $\approx 0$).
\end{enumerate}
\textbf{2. Lampe à DEL blanche (à phosphore).}
\begin{enumerate}
\item Elle est \textbf{polychromatique} : son spectre comporte un pic bleu étroit à $\SI{450}{nm}$ (puce InGaN) et une large bande continue de $\SI{500}{nm}$ à $\SI{700}{nm}$ (phosphore YAG:Ce convertissant le bleu en jaune), donc de nombreuses radiations.
\item Son spectre couvre presque tout le domaine visible : chaque objet reçoit les radiations qu'il est capable de diffuser (rouge, vert, bleu...) et restitue donc des couleurs proches de celles qu'il montre en plein jour. On dit que la DEL blanche a un bon \cle{rendu des couleurs} (IRC $> 80$).
\end{enumerate}
\textbf{3. Comparaison avec l'incandescence (remplacée par ENEO).}
\begin{enumerate}
\item $\lambda_{\max} = \dfrac{\num{2,90e-3}}{2800} = \SI{1,04e-6}{m} = \SI{1036}{nm}$ : dans l'\textbf{infrarouge}.
\item Le profil spectral en cloche du filament (loi de Planck) est centré dans l'infrarouge à $\SI{1036}{nm}$ et s'étend largement de part et d'autre : seule la « queue » de la courbe, entre $\SI{400}{nm}$ et $\SI{800}{nm}$, tombe dans le visible. La \textbf{majeure partie de l'énergie (environ 90-95\%) est donc émise en dehors du visible}, essentiellement en infrarouge (chaleur).
\item L'énergie infrarouge ne sert pas à éclairer : elle est perdue en chaleur. Le \cle{rendement lumineux} de l'incandescence est donc très faible ($\approx \SI{15}{lm.W^{-1}}$), ce qui la rend très énergivore. Les DEL convertissent l'énergie électrique presque exclusivement en lumière visible (rendement $> \SI{100}{lm.W^{-1}}$) : les remplacer réduit fortement la consommation pour un même flux lumineux, d'où les campagnes de remplacement menées par ENEO.
\end{enumerate}
\textbf{4. Synthèse (exemple de rédaction).} La lampe à vapeur de sodium émet une lumière quasi monochromatique jaune ($\SI{589}{nm}$) : son spectre est une raie unique, son rendu des couleurs est très mauvais, mais son efficacité énergétique est correcte ($\approx \SI{150}{lm.W^{-1}}$). La lampe à incandescence émet un spectre continu de corps chauffé, centré dans l'infrarouge à $\SI{2800}{K}$ : elle gaspille l'essentiel de l'énergie en chaleur et son rendement est très faible ($\approx \SI{15}{lm.W^{-1}}$). La DEL blanche émet un spectre polychromatique couvrant tout le visible (pic bleu + bande continue phosphorée) : elle offre un excellent rendu des couleurs et la meilleure efficacité énergétique ($> \SI{100}{lm.W^{-1}}$). C'est donc la solution la plus adaptée à un éclairage public moderne, économe et offrant une bonne visibilité des couleurs pour la sécurité routière.

\medskip
\textbf{Barème indicatif (sur 10)} : 1.a : 1 pt ; 1.b : 1 pt ; 1.c : 1 pt ; 2.a : 1 pt ; 2.b : 1 pt ; 3.a : 1 pt ; 3.b : 1 pt ; 3.c : 1 pt ; 4 : 2 pts.

\textbf{Points de vigilance} : pour 1.c et 2.b, expliciter le mécanisme « un objet diffuse les radiations qu'il reçoit » (synthèse soustractive) ; pour 3.b, l'argument doit s'appuyer sur la forme en cloche du profil spectral (loi de Planck) et la position de $\lambda_{\max}$ hors du visible ; la synthèse doit mentionner les trois critères : nature de la lumière (mono/poly), profil spectral, efficacité énergétique (rendement lumineux).
\end{corrige}

\newpage

\coursec{Fiche bilan}

\begin{popfigure}
\centering
\begin{center}\tcbox[colback=poporange,colframe=poporange,coltext=white,arc=9pt,boxsep=4pt,left=6pt,right=6pt]{\bfseries\small Spectres \& Sources}\end{center}
\vspace{-2pt}
\begin{tcbraster}[raster columns=3,raster equal height=rows,raster column skip=6pt,raster row skip=6pt,enhanced,blankest]
\begin{tcolorbox}[enhanced,colback=white,colframe=poporange,boxrule=0.9pt,arc=3pt,title={Sources \& Spectres},fonttitle=\bfseries\scriptsize,coltitle=white,colbacktitle=poporange,left=4pt,right=4pt,top=2pt,bottom=2pt,boxsep=1pt,toptitle=1.5pt,bottomtitle=1.5pt,halign=left]
\begin{itemize}[nosep,leftmargin=*,itemsep=1pt,font=\scriptsize]
\item Source primaire
\item Source secondaire
\item Prisme / Réseau
\end{itemize}
\end{tcolorbox}
\begin{tcolorbox}[enhanced,colback=white,colframe=popgreen,boxrule=0.9pt,arc=3pt,title={Mono / Poly},fonttitle=\bfseries\scriptsize,coltitle=white,colbacktitle=popgreen,left=4pt,right=4pt,top=2pt,bottom=2pt,boxsep=1pt,toptitle=1.5pt,bottomtitle=1.5pt,halign=left]
\begin{itemize}[nosep,leftmargin=*,itemsep=1pt,font=\scriptsize]
\item Monochromatique $\lambda_0$
\item Polychromatique
\item Ligne / Bande / Continu
\end{itemize}
\end{tcolorbox}
\begin{tcolorbox}[enhanced,colback=white,colframe=poppurple,boxrule=0.9pt,arc=3pt,title={Profils spectraux},fonttitle=\bfseries\scriptsize,coltitle=white,colbacktitle=poppurple,left=4pt,right=4pt,top=2pt,bottom=2pt,boxsep=1pt,toptitle=1.5pt,bottomtitle=1.5pt,halign=left]
\begin{itemize}[nosep,leftmargin=*,itemsep=1pt,font=\scriptsize]
\item Énergie spectrale
\item Couleur dominante
\item Pureté colorimétrique
\end{itemize}
\end{tcolorbox}
\begin{tcolorbox}[enhanced,colback=white,colframe=poppink,boxrule=0.9pt,arc=3pt,title={Ondes EM},fonttitle=\bfseries\scriptsize,coltitle=white,colbacktitle=poppink,left=4pt,right=4pt,top=2pt,bottom=2pt,boxsep=1pt,toptitle=1.5pt,bottomtitle=1.5pt,halign=left]
\begin{itemize}[nosep,leftmargin=*,itemsep=1pt,font=\scriptsize]
\item $c = \lambda \nu$
\item Spectre $\gamma$ à ondes radio
\item Domaine visible $[380;780]$ nm
\end{itemize}
\end{tcolorbox}
\begin{tcolorbox}[enhanced,colback=white,colframe=popgold,boxrule=0.9pt,arc=3pt,title={Loi de Wien},fonttitle=\bfseries\scriptsize,coltitle=white,colbacktitle=popgold,left=4pt,right=4pt,top=2pt,bottom=2pt,boxsep=1pt,toptitle=1.5pt,bottomtitle=1.5pt,halign=left]
\begin{itemize}[nosep,leftmargin=*,itemsep=1pt,font=\scriptsize]
\item $\lambda_{\max} T = b$
\item Corps noir
\item Couleur $\leftrightarrow$ Température
\end{itemize}
\end{tcolorbox}
\begin{tcolorbox}[enhanced,colback=white,colframe=popblue!70!black,boxrule=0.9pt,arc=3pt,title={Expérimentation},fonttitle=\bfseries\scriptsize,coltitle=white,colbacktitle=popblue!70!black,left=4pt,right=4pt,top=2pt,bottom=2pt,boxsep=1pt,toptitle=1.5pt,bottomtitle=1.5pt,halign=left]
\begin{itemize}[nosep,leftmargin=*,itemsep=1pt,font=\scriptsize]
\item Décomposition lumière
\item Synthèse additive
\item Mesure $\lambda_{\max}$
\end{itemize}
\end{tcolorbox}
\end{tcbraster}

\legende{Carte mentale de la leçon : Spectres lumineux et sources de lumière}
\end{popfigure}

\begin{bilan}
\paragraph{Définitions clés}
\begin{itemize}
    \item \cle{Source de lumière} : Objet qui émet de la lumière (primaire : étoile, lampe, LED ; secondaire : objet éclairé qui diffuse).
    \item \cle{Spectre lumineux} : Répartition de l'énergie lumineuse en fonction de la longueur d'onde $\lambda$ (ou fréquence $\nu$). Obtenu par décomposition (prisme, réseau de diffraction).
    \item \cle{Lumière monochromatique} : Une seule longueur d'onde dans le vide $\lambda_0$ (ex. : laser, raie spectrale). Spectre = raie fine.
    \item \cle{Lumière polychromatique} : Plusieurs longueurs d'onde (ex. : Soleil, lampe à incandescence, LED blanche). Spectre = bande, continu ou raies multiples.
    \item \cle{Profil spectral} : Courbe $E(\lambda)$ ou $I(\lambda)$ montrant l'intensité/énergie par unité de longueur d'onde.
    \item \cle{Onde électromagnétique} : Onde transversale se propageant dans le vide à $c \approx \SI{3,00e8}{m/s}$, caractérisée par $\lambda$ et $\nu$ liées par $c = \lambda \nu$.
\end{itemize}

\paragraph{Formules à connaître par cœur}
\begin{center}
\begin{tabularx}{\linewidth}{|>{\centering\arraybackslash}X|>{\centering\arraybackslash}X|>{\centering\arraybackslash}X|}
\hline
\rowcolor{popblueL}
\textbf{Relation} & \textbf{Expression} & \textbf{Unités SI} \\
\hline
Vitesse de la lumière & $c = \lambda \nu$ & $c$ en \si{m/s}, $\lambda$ en \si{m}, $\nu$ en \si{Hz} \\
\hline
Loi de Wien (corps noir) & $\lambda_{\max} T = b$ & $\lambda_{\max}$ en \si{m}, $T$ en \si{K}, $b \approx \SI{2,90e-3}{m.K}$ \\
\hline
Énergie d'un photon & $E = h \nu = \dfrac{hc}{\lambda}$ & $E$ en \si{J}, $h \approx \SI{6,63e-34}{J.s}$ \\
\hline
\end{tabularx}
\end{center}

\paragraph{Méthodes express}
\begin{enumerate}
    \item \textbf{Obtenir un spectre} : Faire traverser la lumière par un prisme (dispersion $n(\lambda)$) ou un réseau de diffraction (interférence). Observer sur écran ou capteur (spectromètre).
    \item \textbf{Classer une source} : Regarder le spectre $\to$ raie unique = monochromatique ; continu / bandes / raies multiples = polychromatique.
    \item \textbf{Analyser un profil spectral} : Identifier $\lambda_{\max}$ (pic), largeur à mi-hauteur (pureté), domaine spectral (UV, visible, IR). Couleur perçue $\approx$ couleur dominante.
    \item \textbf{Appliquer la loi de Wien} : Mesurer $\lambda_{\max}$ du pic d'émission $\to$ déduire $T = b/\lambda_{\max}$ (température de surface d'étoile, four, filament).
\end{enumerate}
\end{bilan}

\begin{aretenir}[Checklist avant le Probatoire]
\begin{itemize}
    \item $\square$ Savoir définir source primaire / secondaire et donner des exemples camerounais (Soleil, lampe torche, écran téléphone, route la nuit).
    \item $\square$ Expliquer le principe de décomposition par prisme (indice $n$ décroissant avec $\lambda$ : $n_{violet} > n_{rouge}$).
    \item $\square$ Différencier sans hésiter spectre de raies, spectre de bandes, spectre continu ; associer à la source (gaz raréfié, molécules, solide/liquide chaud).
    \item $\square$ Écrire et utiliser $c = \lambda \nu$ avec conversions correctes (nm $\to$ m, THz $\to$ Hz) et analyse dimensionnelle $[c] = LT^{-1}$.
    \item $\square$ Situer le domaine visible $[\SI{380}{nm}; \SI{780}{nm}]$ dans le spectre EM (UV à gauche, IR à droite) et ordonner les couleurs (ROJVBLE).
    \item $\square$ Énoncer la loi de Wien $\lambda_{\max} T = b$ et préciser qu'elle s'applique aux \emph{corps noirs} (approximation : étoiles, filaments).
    \item $\square$ Calculer $\lambda_{\max}$ connaissant $T$ (ex. : Soleil $T \approx \SI{5800}{K} \to \lambda_{\max} \approx \SI{500}{nm}$ vert) ou $T$ connaissant $\lambda_{\max}$.
    \item $\square$ Relier la couleur d'un corps chauffé à sa température (rouge $\approx \SI{1000}{K}$, jaune $\approx \SI{3000}{K}$, blanc $\approx \SI{6000}{K}$, bleu $> \SI{10000}{K}$).
    \item $\square$ Décrire la synthèse additive (Rouge + Vert + Bleu = Blanc) et la synthèse soustractive (filtres, pigments).
    \item $\square$ Savoir exploiter un profil spectral $I(\lambda)$ : lire $\lambda_{\max}$, estimer la largeur spectrale, commenter la pureté colorimétrique.
    \item $\square$ Respecter les unités SI et les chiffres significatifs (ex. : $\lambda = \SI{632,8}{nm}$, pas $633$ nm pour un He-Ne ; $b = \SI{2,90e-3}{m.K}$).
\end{itemize}
\end{aretenir}

\findecours

\end{document}
